\documentclass[11pt]{article}
\usepackage[margin=1in]{geometry}
\usepackage{amsmath,amssymb,booktabs,longtable,array}
\usepackage[hidelinks]{hyperref}
\newcommand{\dd}{\mathrm d}
\newcommand{\DE}{\mathrm{DE}}
\newcommand{\Pl}{\mathrm{Pl}}
\newcommand{\Mb}{M_b}
\newcommand{\lp}{\lambda_{HX}}
\newcommand{\lt}{\lambda_{\rm tot}}
\newcommand{\eq}{\mathrm{eq}}
\title{Cold energy, dark energy, and a particle bridge\\
Equations, physical roles, executed routes, and remaining closure}
\author{Zimmerman programme: consolidated CPB-QFG-2026-10-16 dossier}
\date{16 October 2026}
\begin{document}
\maketitle
\section*{Current result and reading guide}
The physical theory remains incomplete. This is one standalone equation
dossier containing the current inventory and the successive calculations
needed to understand its limits. It is not a single unified action assembled
by identifying unrelated fields, nor a claim that every historical candidate
in the repository is physically compatible.

Three physical roles remain distinct: vacuum stress, massive cold excitations,
and collective gravitational response. The proposed link
$a_0=\kappa c\sqrt{G\rho_{\rm DE}}$, with fitted $\kappa$, remains a target.
The particle construction supplies actual interactions, while its
identification with the galaxy-response medium remains unresolved.

Part I gives the equations and initial scalar production construction.
Part II executes its thermal and charge audit, including a nonzero
number-changing amplitude. Part III executes a reduced kinetic checkpoint:
the pole-free reaction $\Phi\Phi^*\leftrightarrow XX^*$ and fast source
conversion can change the final daughter count substantially. An isolated,
energy-conserving reduced gas calculation is executed from specified initial
states. It does not yet derive those states from the Higgs era or include
the complete real-intermediate-state-consistent collision hierarchy.

The original 1-MeV thermal comparison and the separately specified 2-MeV
evolution are different candidates. The latter's moderate initial state
nearly doubles its weighted count in the reduced model; the hottest state
fails the original dilution gate. No observational abundance or coldness
verdict is inferred. That calculation also extends the dust exchange
identity to a warm gas, preventing internal rest/kinetic conversion from
being misidentified as vacuum exchange.

Part IV records the collision audit. The same action has a missing
pole-free source-assisted channel, $X+\Phi+\Phi^*\leftrightarrow\Phi+\Phi$.
Its two conjugate channels add approximately 4.95 times the pair-conversion
coefficient at $T_d=m$ in the chosen common-MB approximation. The old
abundance curves therefore remain reduced-network results. Exact source-pole
factorization identifies further required reactions. A positive interior
pole-matching construction has an explicit smooth-profile error bound;
physical endpoint and overlapping-pole matching remain unfinished.

Part V executes that pole-free source-assisted evolution. Its final weighted
yield ratios are approximately 1.02719, 1.95810 and 4.49358 for the three
specified initial temperatures. The hottest state still fails the original
dilution gate. These remain reduced-network results, with further resonant
channels and the actual production history unfinished. The same part derives
the direct Higgs-induced baryonic contact response: at the chosen galactic
density its fractional shift is about $7\times10^{-48}$. An ideal-gas halo
occupation check also prevents simply assuming a MeV condensate. These are
candidate-specific constraints, not exclusions of collective media.

Part VI derives a finite-time bound for stable canonical quadratic response.
Large static susceptibility from a soft mode does not imply large response
at a fixed time. It applies this test to the existing scalar/polarization
construction with explicit limits on the approximation. A separate one-loop
calculation shows that a tree-zero daughter Higgs coupling acquires a finite
soft contribution. Its chosen galactic effect remains tiny. This soft vertex
is not a hard Higgs decay calculation or a full renormalization result.

Part VII constructs a separately specified NR density bridge to the nodal
medium. Its canonical occupation energy and a chosen repulsive contact
give a sufficient joint convexity condition. Both cold and fermion densities
are varied at fixed inventories in a prescribed external baryon potential.
A neutral cubic-response interval survives, with compensating signed fermion
charge elsewhere. The strong branch has a proved boundary-sensitive
large-radius limit and can lose the desired baryonic mass scaling. This is
an explicit effective candidate, not a matching from the MeV scalar action,
an evolved halo, or a relativistic closure.

Part VIII constructs a charge-even scalar bridge directly to $|\Phi|^2$.
A completed-square potential retains the carrier mass lower bound in the
specified positive-energy density-functional extension. The actual neutral
stock is distinguished from signed charge. Chosen vacuum tree scattering
and a new mediator decay are calculated. A direct band calculation supplies
the density-vertex kernel, its uniform Ward zero, finite-momentum absorption
and a conservative static RPA test. The original mass-vertex response cannot
be substituted for this density vertex. Neither full quantum stability,
neutral-mixture EOS, a populated-mediator thermal history nor galaxy closure
is established.

Part IX computes selected scalar loop effects. A small carrier mass shift
can close the narrow source decay in a diagnostic with fixed MS mass inputs;
finite on-shell matching remains possible and a weaker comparison keeps the
decay open. Tadpole, mass and vacuum matching are separate conditions.
Partial scalar vacuum terms do not predict the cosmological constant.
A known sequestering gravity route is checked conditionally; directly
coupling an acceleration coefficient to its auxiliary vacuum parameter
changes the constraint, so the previous cancellation argument cannot be
inherited without the new variation. Neither residual vacuum magnitude nor
the acceleration normalization is derived.

Part X varies an explicit auxiliary-action extension with a separate
conjugate flux $Q$. A Planck-scaled flux energy with selected power two
survives the projected metric equation, unlike an unscaled constant flux.
A chosen inverse-flux cubic response produces $a_0\propto Q$ while that
flux energy is proportional to $Q^2$. Residual flux/history terms, a
covariant microscopic matter/frame action, the functional choice and the
normalization remain required inputs or obligations. The scalar inverse-square
force is counted in the Newtonian dictionary. The annular response has
logarithmically divergent infinite-tail energy, so it is not a complete halo.
This is a conditional classical action candidate, not physical closure.

Part XI supplies a covariant barotropic material frame and its stress.
A material-scaled negative band sea has radiation-like pressure and fails
by itself; a specified positive host repairs the elementary longitudinal
signs while introducing new EOS inputs. A fixed-velocity material spectrum
makes the cubic acceleration vary with density. Compensating this coefficient
requires a density-dependent critical subtraction and its host rearrangement.
On a prescribed finite-count profile, neither persistent cubic response nor
the compensated regulated tail has finite infinite-radius band energy.
The material alternative returns to an inverse-square toy tail after leaving
the controlled small-gap regime. These tests identify what an actual phase
interface must solve; no full fermion/frame or finite halo is obtained.

Part XII tests a different particle-level host: a sextic charged scalar
phase with a computed planar wall and a chosen reciprocal local response.
Fixed-charge spherical solutions have a cubic-dominated interior annulus
and a finite-energy inverse-square exterior. Two numerical resolutions and
a fixed-charge scaling identity check stationarity; finite radial amplitude
Hessians have positive sampled curvature. This is not the repulsive MeV
carrier or a microscopic nodal-band matching. Its exterior changes the
Newtonian force coefficient substantially, while its charge inventory,
formation, dynamic/metric health and cosmology remain unclosed.

Part XIII supplies a renormalizable two-scalar particle model for the
attractive host phase. Its selected family has an algebraically proved
unique global classical vacuum and a unique charged coexistence density.
Exact matching shows a controlled heavy-mediator limit, including the
amplitude kinetic correction. Improving that tree limit increases a
quartic-cancellation requirement. Quantum thresholds, a derivation of the
response coupling and a new two-field finite-body/formation calculation
remain required. This is a classical particle-potential bridge, not a
completed theory of cold energy, dark energy or gravity.

Part XIV explicitly tests quantum corrections, formation, gravity response
and vacuum normalization. One-loop running exposes the unprotected original
quartic/mass matching. A weak-coupling classical-limit scaling controls the
selected loop sensitivities while increasing the required charge and energy.
Actual charged homogeneous scalar backgrounds exhibit a finite clumping
instability band; nonlinear formation and the initial charged state are
not yet derived. The minimal Lorentz scalar cubic-response completion has
radial characteristics approaching $\sqrt2$ times the metric light speed;
additional physical dynamics are needed for a common-cone completion.
Independent response parameters and flux/history data prevent a unique
prediction of the dark-energy magnitude or acceleration normalization.

The reports retain their exact model boundaries, sources and evidence.
Their local ``next'' sections describe their own checkpoints; the current
next steps are to match the remaining
source-assisted and six-point contributions at actual thermal boundaries,
then compute the stabilized production/thermal history. Galaxy
formation, finite supply, source selection, relativistic lensing, native
cosmology, parameter origins and joint empirical validation remain open.
No global novelty or independent-referee claim is made.

\tableofcontents
\clearpage

\part{Equation inventory and the first production construction}

\begin{abstract}
The physical theory remains open. This report separates vacuum stress,
mass-carrying excitations, and collective response, and records the equations
needed to connect them. Three routes are executed: a charge-preserving Higgs
production extension, an exact cold--vacuum exchange consistency relation,
and an exterior-mass lensing test. The simplest portal extension encounters
mass-potential and kinetic-evolution obstructions at the adopted benchmark;
a source-quartic extension repairs the global classical potential exactly;
the exchange relation supplies a falsifiable joint constraint; positive
spherical exterior mass cannot increase interior excess tangential shear.
These are conditional reductions and applications of known mechanisms,
not a particle discovery, a Standard Model derivation, or global novelty claims.
\end{abstract}

\section{Target, evidence, and conventions}
The requested goal is one admitted physical model producing the galactic law,
its finite mass supply, formation and source selection, the lensing metric,
and cosmological abundance and perturbations. A predictive effective theory
may contain measured constants. The stronger goal of deriving every constant
also requires a mechanism selecting those constants; it is a separate standard.
Neither standard is met by combining successful tests of different actions.

The baseline is Git revision
\texttt{370dfb58006e330f210f5e306297e822188e37ba} in a shared dirty checkout.
Inputs to the calculation are hashed in \texttt{run\_003/manifest.json}.
The controlling recent work order is
\path{sol61_push/THEORY_NEXT_STEPS_2026-10-15.md}, including its October 15
fourth-round successor. The frozen October 14 candidate is preserved.
Older README and retraction prose describe different stages and do not
override the dated candidate-specific results.

\textbf{MODEL} denotes a chosen action/state; \textbf{DERIVED} means conditional
on the displayed hypotheses; \textbf{TARGET} denotes a relation to explain;
\textbf{OPEN} denotes a missing implication. Equations from distinct models
below are not silently parts of one action. Microscopic formulas use
$\hbar=c=k_B=1$, signature $(+---)$; cosmological SI equations explicitly
use mass-equivalent densities $\rho_i$ and $p_i/c^2$. The weak lensing metric
uses $(-+++)$ separately. These conventions must be translated in a future
joined action. $\Phi_c$ below is the complex carrier; $\phi_g$ is gravitational
potential; $\mathbf P$ is polarization; $p$ is pressure.

The two acceleration footings remain separate:
\begin{equation}
 a_{0,\mathrm{can}}=9.3603\,10^{-11}\ \mathrm{m\,s^{-2}},\qquad
 a_{0,\mathrm{alt}}=1.1312\,10^{-10}\ \mathrm{m\,s^{-2}}.
\end{equation}
These are the frozen candidate's values; nearby older documents use slightly
different rounded values. The alternative total-density footing must not
be identified with vacuum density in the equations below. The new microscopic
controls do not choose or combine these galaxy footings.

\section{What cold and dark energy mean}
Cold energy is energy mostly stored as rest mass, with
$|p_c|\ll\rho_c c^2$ in the relevant frame. It is not an identified species.
Dark energy is a component with sufficiently negative pressure to accelerate
expansion. A constant positive vacuum has
\begin{equation}
 \epsilon_v=\rho_v c^2>0,\qquad p_v=-\epsilon_v,\qquad w_v=-1.
\end{equation}
A nonrelativistic isotropic distribution instead has, at leading order,
\begin{equation}
 \rho_c c^2\simeq mn c^2+\tfrac12 mn\langle v^2\rangle,
 \qquad p_c\simeq\tfrac13 mn\langle v^2\rangle.
\end{equation}
Low mean pressure does not establish a small perturbation sound speed,
short free-streaming length, correct abundance, or a galactic attractor.

A simple example of one field carrying both roles is the stipulated potential
\begin{equation}
 V(\Phi_c)=V_0+m^2|\Phi_c|^2+\lambda|\Phi_c|^4,\quad
 \Phi_c\simeq e^{-imt}\psi/\sqrt{2m},\quad n=|\psi|^2.
\end{equation}
In its dilute rotating, weak-interaction limit,
\begin{equation}
 \epsilon\simeq V_0+mn+\frac{\lambda n^2}{4m^2},\qquad
 p\simeq-V_0+\frac{\lambda n^2}{4m^2}.
\end{equation}
Here $mn$ is rest-energy density in natural units. The interaction is
repulsive for $\lambda>0$. The approximate adiabatic sound speed of the
nonvacuum part is
\begin{equation}
 c_{s,c}^2\simeq
 \frac{\lambda n/(2m^2)}{m+\lambda n/(2m^2)}.
\end{equation}
This illustrates coexistence, not the desired galaxy force. Changing $V_0$
leaves the flat-space scalar equation unchanged while changing vacuum gravity.
Thus writing one field does not fix the cold abundance, vacuum magnitude,
or $a_0$. Previously tested linked-square potentials also have excessive
pressure/scale-separation problems; their failures are not erased here.

\section{Dark energy, expansion, and the acceleration scale}
For FLRW geometry,
\begin{align}
 H^2&=\frac{8\pi G}{3}\sum_i\rho_i-\frac{k c^2}{a^2},\\
 \frac{\ddot a}{a}&=-\frac{4\pi G}{3}\sum_i(\rho_i+3p_i/c^2).
\end{align}
Vacuum may be written as density on the right or as $\Lambda$ on the left
of Einstein's equation, once:
\begin{equation}
 \rho_v=\frac{\Lambda c^2}{8\pi G},\quad
 H_\Lambda^2=\frac{\Lambda c^2}{3},\quad
 a_0=\kappa c\sqrt{G\rho_v}
 =\kappa c^2\sqrt{\frac{\Lambda}{8\pi}}.
\label{eq:scale}
\end{equation}
The last link is the programme's \textbf{TARGET}, with $\kappa=1/2$ fitted,
not derived. Consequently
\begin{equation}
 a_0^2=\frac{\Lambda c^4}{32\pi},\quad
 \epsilon_v=\frac{4a_0^2}{G},\quad
 \frac{a_0}{cH_\Lambda}=\frac{1}{2}\sqrt{\frac{3}{8\pi}}.
\end{equation}
These three are algebraic restatements of the adopted coefficient.
They are not three independent derivations. A constant offset of a response
potential remains a free vacuum term unless an actual mechanism fixes it.

For separately conserved dark energy,
\begin{equation}
 \frac{\rho_\DE(z)}{\rho_{\DE0}}
 =\exp\!\left[3\int_0^z\frac{1+w(z')}{1+z'}\dd z'\right],\qquad
 \frac{a_0(z)}{a_0(0)}=\sqrt{\frac{\rho_\DE(z)}{\rho_{\DE0}}}.
\end{equation}
For CPL, $w(a)=w_0+w_a(1-a)$,
\begin{equation}
 F(z)=(1+z)^{3(1+w_0+w_a)}e^{-3w_a z/(1+z)}.
\end{equation}
Constant vacuum gives constant $a_0$. Evolving CPL density, an $H(z)$ law,
and an event-horizon law are distinct hypotheses. An event-horizon model
instead requires a specified future expansion history:
\begin{equation}
 R_E(a)=ca\int_a^\infty\frac{\dd a'}{a'^2 H(a')},\qquad
 a_0(a)/a_0(1)=R_E(1)/R_E(a),
\end{equation}
when the integral exists. No existing seal is changed here.

\section{The galaxy law and the static polarization construction}
For spherical baryons,
\begin{equation}
 g_b=\frac{G\Mb(<r)}{r^2},\quad y=g_b/a_0,\quad
 g=\nu(y)g_b,\quad g\sim\sqrt{a_0g_b},\quad V_f^4=G\Mb a_0.
\end{equation}
Three distinct interpolations appearing in the corpus are
\begin{align}
 \nu_s(y)&=\tfrac12+\sqrt{\tfrac14+1/y},
 &g_b&=\frac{g^2}{g+a_0},\\
 \nu_q(y)&=\sqrt{1+1/y},
 &g^2-g_b^2&=a_0g_b,\\
 \nu_e(y)&=(1-e^{-\sqrt y})^{-1}.
\end{align}
They share a deep limit and differ elsewhere. They must not be called
equivalent. The October 14 static candidate uses $\nu_s$.

Put $K=4\pi G$, $u=K|\mathbf P|<a_0$, and
$\gamma=K L_*^2/2$. Its chosen static action density is
\begin{equation}
 \mathcal L_{S2}=-\frac{|\nabla\phi_g-K\mathbf P|^2}{2K}
 -\rho_b\phi_g-W_{\rm rem}(u)-\frac\gamma2|\nabla\mathbf P|^2,
\end{equation}
where the critical quadratic part has already been included in the square:
\begin{equation}
 W_{\rm rem}(u)=\frac1K\left[-\frac{u^2}{2}-a_0u
 -a_0^2\ln(1-u/a_0)\right],\qquad
 K W'_{\rm rem}=\frac{u^2}{a_0-u}.
\end{equation}
Variation, before a phase gate is added, gives
\begin{align}
 \nabla\cdot(\nabla\phi_g-K\mathbf P)&=K\rho_b,\\
 \nabla\phi_g-K\mathbf P-KW'_{\rm rem}(u)\widehat{\mathbf P}
 +\gamma\nabla^2\mathbf P&=0.
\end{align}
Outside stiffness-dominated regions, in the spherical aligned branch,
\begin{equation}
 g-u=g_b=\frac{u^2}{a_0-u},\qquad
 u=\frac{a_0g}{a_0+g},\qquad g_b=\frac{g^2}{g+a_0}.
\end{equation}
This is \textbf{DERIVED} from a potential chosen to give the law. It does not
derive that potential from microscopic physics. $L_*$ is a bounded input.
The static action alone supplies no inertia, collision operator, physical
clock, or condensed-baryon ownership prescription. A local AQUAL-like
reduction also cannot by itself establish absence of an external-field effect.

\section{Mass supply, currents, and free matter}
If $\mathbf P$ is actual cold-mass displacement, the proposed identification is
\begin{equation}
 \nabla\cdot\mathbf P=\rho_c-\bar\rho_c,\qquad
 \dot{\mathbf P}=-\mathbf j_c.
\end{equation}
Continuity alone implies only
$\nabla\cdot(\mathbf j_c+\dot{\mathbf P})=0$.
For homogeneous positive $\rho_0$, ordinary material kinetic energy requires
\begin{equation}
 T\ge\frac{|\mathbf j_c|^2}{2\rho_0},\qquad
 T_k\ge\frac{|\dot{\delta\rho}_k|^2}{2\rho_0k^2},\qquad
 I_{\rm mass}=\rho_0^{-1},\quad
 c_{\rm bare}^2=\gamma\rho_0=2\pi G\rho_0L_*^2.
\end{equation}
This reproduces the October 15 mass-response audit. Choosing an independent
fast response field changes the model; it does not remove mass transport.
For a localized contrast,
\begin{equation}
 \int\delta\rho\,\dd^3x=\oint\mathbf P\cdot\dd\mathbf S.
\end{equation}
Zero boundary flux implies compensation, not creation of positive mass
from an empty isolated domain.

In the zero-stiffness point-baryon $\nu_s$ branch put
$r_M=\sqrt{G\Mb/a_0}$ and $x=r/r_M$. Then
\begin{equation}
 \frac{M_t(<r)}{\Mb}=\frac{\sqrt{1+4x^2}-1}{2},\quad
 M_t=\frac{r^2(g-g_b)}{G}.
\end{equation}
An allocated stock $S_c=\mathcal R\Mb$ therefore gives
\begin{equation}
 x\le\sqrt{\mathcal R(\mathcal R+1)}=5.8426446\ldots
 \quad(\mathcal R=5.364).
\end{equation}
The numerical stock is an input inferred in a cosmological model, not a
derivation or a measured species identity. More generally
\begin{equation}
 M_{c,\rm target}=(\nu-1)\Mb,\quad
 \rho_{c,\rm target}=\frac{1}{4\pi r^2}\frac{\dd M_{c,\rm target}}{\dd r}.
\end{equation}
Both enclosed positivity and density positivity must hold.
The earlier bookkeeping $M_{\rm tot}=\nu\Mb+k M_{\rm hot}$ has fitted
$k\simeq3$; it is not a physical partition law. The newer work order retains
the failure of a gas-tracing free pool and requires a counted free population.

For a conditional extension where free mass supplies ordinary source
$A=g_b+g_f$ to the same responding phase,
\begin{equation}
 A=\frac{g^2}{g+a_0},\quad a_{\rm app}=\frac{g(g-g_b)}{g_b},\quad
 M_f(<r)=\Mb(<r)\frac{a_{\rm app}-a_0}{g+a_0}.
\end{equation}
The October 15 inverse diagnostic requires $M_f\ge0$ and $M'_f\ge0$.
Unchanged exact baryon-only adherence then leaves no such free mass inside
the radius. A different source/phase assignment requires different equations.

\section{The microscopic cold carrier}
The independent charge-protected scalar model is
\begin{align}
 \mathcal L_c&=|\partial\Phi_c|^2+|\partial X|^2-V_c,\\
 V_c&=m^2|\Phi_c|^2+4m^2|X|^2+
 \left|\sqrt{\Delta_c}X+\frac{g_c}{\sqrt{\Delta_c}}\Phi_c^2\right|^2
 +\lambda_{\rm rep}|\Phi_c|^4,\\
 M&=2m\sqrt{1+u_*^2},\quad \Delta_c=4m^2u_*^2,
 \quad\lambda_{\rm tot}=g_c^2/\Delta_c+\lambda_{\rm rep}.
\end{align}
Here $m,\Delta_c,\lambda_{\rm rep}>0$ and $g_c$ is real. Its equations are
\begin{equation}
 (\Box+m^2)\Phi_c+2g_c X\Phi_c^*
 +2\lambda_{\rm tot}|\Phi_c|^2\Phi_c=0,\qquad
 (\Box+M^2)X+g_c\Phi_c^2=0.
\end{equation}
The symmetry $(\Phi_c,X)\mapsto(e^{i\theta}\Phi_c,e^{2i\theta}X)$ yields
\begin{equation}
 j_Q^\mu=i(\Phi_c^*\partial^\mu\Phi_c-\Phi_c\partial^\mu\Phi_c^*)
 +2i(X^*\partial^\mu X-X\partial^\mu X^*),\quad \partial_\mu j_Q^\mu=0.
\end{equation}
In the isolated classical model $E\ge m|Q|$. This carrier-only bound does
not automatically include gravity, a portal, or a finite-temperature medium.

At a fixed regulator in the nonrelativistic approximation put
$\delta=M-2m$, $c_3=g_c/(2m\sqrt{2M})$, $k_c=\lambda_{\rm tot}/(4m^2)$.
With source annihilator $b$,
\begin{align}
 H&=mQ+D_c,\quad Q=\int(\psi^\dagger\psi+2b^\dagger b)\dd^3x,\\
 D_c&=T+\delta\int\left|b+\frac{c_3}{\delta}\psi^2\right|^2\dd^3x
 +\left(k_c-\frac{c_3^2}{\delta}\right)\int\psi^{\dagger2}\psi^2\dd^3x,\\
 k_c-\frac{c_3^2}{\delta}
 &=\frac{\lambda_{\rm rep}}{4m^2}
 +\frac{g_c^2}{8m^2M(M+2m)}>0.
\end{align}
Thus $0\le\langle T\rangle\le\langle D_c\rangle$. For initially resting
sources $D_c/(mQ)=\delta/(2m)\simeq u_*^2/2$.
The regulated homogeneous pressure bound uses
\begin{equation}
 r_c^2=\frac{c_3^2}{\delta k_c}<1,\quad
 b_\pm=\tfrac12(1\pm(1-r_c^2)^{-1/2}),\quad
 \frac{|p_c|}{\epsilon_c}\le\frac{b_+d}{1+d},\quad
 d=\frac{\langle D_c\rangle}{m\langle Q\rangle}.
\end{equation}
It is not a continuum quantum-field or perturbation theorem.

A resting source decays to two same-charge daughters, with
\begin{equation}
 p_*=mu_*,\quad v_*^2=\frac{u_*^2}{1+u_*^2},\quad
 \frac{p_{\rm free}}{\epsilon_{\rm free}}=\frac{u_*^2}{3(1+u_*^2)},\quad
 \Gamma_c=\frac{g_c^2}{8\pi M}\sqrt{1-\frac{4m^2}{M^2}}.
\end{equation}
At threshold and away from an unresolved resonance,
\begin{align}
 \sigma_0&=\frac{\lambda_{\rm rep}^2}{8\pi m^2},\quad
 x_s=\frac{s-4m^2}{\Delta_c},\quad r_\lambda=\lambda_{\rm rep}\Delta_c/g_c^2,\\
 \lambda_{\rm eff}&=\frac{g_c^2}{\Delta_c}
 \left(r_\lambda-\frac{x_s}{1-x_s}\right),\quad
 \frac{\sigma(s)}{\sigma_0}=\left(\frac{\lambda_{\rm eff}}{\lambda_{\rm rep}}\right)^2
 \frac{4m^2}{s}.
\end{align}
A physical resonance needs a width and inclusive rates without double-counting
real intermediate $X$. The existing screening example is $m=1$ MeV,
$u_*=.02$, $r_\lambda=2$, $\sigma_0/m=.3\ \mathrm{cm^2/g}$,
$g_c\simeq2168$ eV, $\lambda_{\rm tot}\simeq.00881$.
It is a selected benchmark, not an identified particle.

\section{A distinct microscopic response candidate}
The separate nodal-fermion model has
\begin{align}
 \mathcal L_{\varphi P}&=\frac{Z_\varphi}{2}\dot\varphi^2
 +\frac{I_P}{2}|\dot{\mathbf P}|^2
 -\frac C2|\nabla\varphi-\mathbf P/C|^2-q_b\rho_b\varphi,\\
 H_F&=\frac{k_x^2+k_y^2-k_0^2}{2m_*}\Gamma_1+vk_z\Gamma_2
 +h\sum_{j=1}^3P_j\Gamma_{j+2},\quad
 \{\Gamma_i,\Gamma_j\}=2\delta_{ij}.
\end{align}
For $N$ occupied negative bands, zero temperature and zero chemical potential,
let $\Delta_f=h|\mathbf P|$ and $A=Nm_*/(2\pi v)$. With an energy-tube
cutoff $\Lambda_f$,
\begin{equation}
 V_F(\Delta_f)-V_F(0)=-\frac A3
 [(\Lambda_f^2+\Delta_f^2)^{3/2}-\Lambda_f^3-|\Delta_f|^3].
\end{equation}
Subtracting the analytic quadratic term gives leading
$V_{F,\rm ren}=\beta_3|\mathbf P|^3/3$, $\beta_3=Ah^3>0$.
The local cubic approximation gives
\begin{equation}
 V_\infty^4=\mathcal K\Mb,\qquad
 \mathcal K=\frac{q_b^3}{4\pi C^3\beta_3}.
\end{equation}
The required $\mathcal K=Ga_0$ remains an unfixed calibration.
In the separately derived gapless spatial regime,
\begin{equation}
 \mathbf D/C=\eta|\nabla|\mathbf P,\quad
 \eta=Nm_*h^2\pi/32,\quad
 V_\infty^2=\frac{q_b^2\Mb}{2\pi^2 C^2\eta},
\end{equation}
so $V^4\propto\Mb^2$, not the target exponent. The crossover is unsolved.
Finite temperature or occupancy further changes the local response:
\begin{equation}
 \partial_{\Delta_f}V_F=2AT\Delta_f\ln[2\cosh(\Delta_f/(2T))],\qquad
 \left.\partial_{\Delta_f}V_F\right|_{T=0}
 =A\Delta_f\max(\Delta_f,|\mu|).
\end{equation}
For $\Delta_f\ll T$, the leading potential is quadratic. This already-known
obstruction is not a new route discovered here. Positive homogeneous
response/absorption does not establish galaxy-age settling or covariant lensing.
No interaction connects this fermion medium to $\Phi_c,X$ in the current model.

\section{Route A: a Standard Model production bridge, actually screened}
\textbf{Proposal, not an adopted replacement.} Keep $\Phi_c,X$ as Standard
Model singlets with their global charges, and introduce
\begin{equation}
 \mathcal L_{\rm portal}=-\lambda_{HX}|H_{\rm SM}|^2|X|^2
 -\lambda_{H\Phi}|H_{\rm SM}|^2|\Phi_c|^2.
\end{equation}
This preserves the Standard Model gauge group and the dark $U(1)$. Initially
take $\lambda_{H\Phi}=0$ to screen the simplest extension. The broken-phase
channel is $h\to XX^*$ followed by $X\to\Phi_c\Phi_c$ and its conjugate.
It creates four daughter quanta, two of each sign, with net charge zero.
No net charge asymmetry is generated. In that neutral state the bound
$E\ge m|Q|$ supplies no nonzero abundance bound.

Freeze-in is established prior art \cite{hall}; this application supplies
the missing source interaction rather than deriving the Standard Model.
For a Maxwell--Boltzmann thermal Higgs bath, negligible inverse reactions,
radiation domination, constant $g_s,g_\rho$, high enough reheating, unit
cascade branching, and no number-changing dark reprocessing, integration
of the decay collision term gives
\begin{align}
 \dot n_\Phi+3Hn_\Phi&=4\gamma_{h\to XX^*},\\
 \gamma_{h\to XX^*}&=\frac{m_h^2T\Gamma_h}{2\pi^2}K_1(m_h/T),\\
 Y_\Phi&=4\frac{135}{8\pi^3(1.66)g_s\sqrt{g_\rho}}
 \frac{M_\Pl\Gamma_h}{m_h^2},\\
 \Gamma_h&=\frac{\lambda_{HX}^2 v_H^2}{16\pi m_h}
 \sqrt{1-4M^2/m_h^2}.
\end{align}
Here $M_\Pl=G^{-1/2}$, $Y_\Phi$ counts particles plus antiparticles, and
$Y_\Phi=2Y_{X+X^*}$. These assumptions define a production control, not
a solved dark-sector cosmology. With stipulated $\Omega_c h^2=.12$,
the illustrative conversion $\Omega_c h^2=2.742\,10^8(m/\mathrm{GeV})Y_\Phi$
and $m_h=125$ GeV, $v_H=246$ GeV, $g_s=g_\rho=106.75$, the control gives
\begin{equation}
 Y_\Phi=4.3764\,10^{-7},\quad
 \Gamma_h=4.7104\,10^{-19}\ \mathrm{GeV},\quad
 \lambda_{HX}=2.2115\,10^{-10}.
\end{equation}
The coupling is calibrated to abundance; abundance is not derived.
Assuming total Higgs width $4.1$ MeV gives illustrative branching
$1.15\,10^{-16}$, not a promised collider signal.

\subsection{First obstacle: physical masses and the scalar potential}
After electroweak symmetry breaking,
\begin{equation}
 M^2=M_0^2+\lambda_{HX}v_H^2/2,\quad
 m^2=m_0^2+\lambda_{H\Phi}v_H^2/2,\quad
 \Delta_c=\Delta_0+(\lambda_{HX}-4\lambda_{H\Phi})v_H^2/2.
\end{equation}
At the screened point the source shift is $6.6915\,10^{-6}$ GeV$^2$,
or $4182.2\Delta_c$. Moreover $M_0^2=-2.6899\,10^{-6}$ GeV$^2<0$.
The original carrier has no $|X|^4$ stabilizer. Holding $H=\Phi_c=0$ and
sending $|X|\to\infty$ then makes this minimally extended classical
potential unbounded. Thus this particular positive-portal, fixed-physical-mass
benchmark is not an admitted complete scalar model as written.
Forcing $M_0^2\ge0$ bounds
\begin{equation}
 \lambda_{HX}\le2M^2/v_H^2\simeq1.3225\,10^{-10},\qquad
 (\Omega_c h^2)_{\max}\simeq.0429
\end{equation}
within the same simplified production calculation.
The following subsection executes the positive-quartic repair. Thermal
masses, additional production channels, and changed mass/coupling points
remain distinct continuations. Setting $\lambda_{H\Phi}=\lambda_{HX}/4$
removes this tree-level contribution to the gap but is an unprotected
relation here, not an explanation. The dimensional dark-loop diagnostic
$g_c^2/(16\pi^2\Delta_c)=1.86\,10^{-5}$ is not a renormalized loop result.

\subsection{An explicit globally stable scalar extension}
Keep the desired physical masses in $V_c$ and set
$t=|H_{\rm SM}|^2$, $w_H=t-v_H^2/2$. Specify the new candidate
\begin{equation}
 V_{\rm joined}=V_c+\lambda_H w_H^2
 +\lambda_{HX}w_H|X|^2+\lambda_X|X|^4,
 \qquad \lambda_H=\frac{m_h^2}{2v_H^2}.
\end{equation}
This includes the compensating bare source mass and is gauge invariant.
Completion of a square gives the exact identity
\begin{equation}
 V_{\rm joined}=V_c+
 \lambda_H\left(w_H+\frac{\lambda_{HX}|X|^2}{2\lambda_H}\right)^2
 +\left(\lambda_X-\frac{\lambda_{HX}^2}{4\lambda_H}\right)|X|^4.
\end{equation}
Therefore $\lambda_X\ge\lambda_{HX}^2/(4\lambda_H)$ suffices for
$V_{\rm joined}\ge V_c\ge0$ for all scalar configurations, not just
near the chosen vacuum. The vacuum $X=\Phi_c=0$, $t=v_H^2/2$ is a global
minimum. For the benchmark the sufficient floor is
$\lambda_X\ge9.47\,10^{-20}$; choose $\lambda_X=10^{-18}$ as an explicit
new input. The negative bare source mass at $H=0$ no longer implies an
unbounded total potential. This is a classical boundedness result, not
a claim about the early-universe phase or radiative naturalness.

With canonical scalar kinetic terms, the coupled source and Higgs equations
include, before the additional Standard Model gauge/fermion contributions,
\begin{align}
 (\Box+M^2)X+g_c\Phi_c^2+
 \lambda_{HX}w_H X+2\lambda_X|X|^2X&=0,\\
 D_\mu D^\mu H_{\rm SM}+
 (2\lambda_Hw_H+\lambda_{HX}|X|^2)H_{\rm SM}&=0.
\end{align}
The daughter equation is unchanged at tree level in this specified
$\lambda_{H\Phi}=0$ model. Both portal forces come from the same potential;
the scalar subsystem exchanges energy reciprocally and retains the dark
$U(1)$. With vacuum energy subtracted, in isolated flat-space classical
scalar theory, the completed-square lower bound also retains $E\ge m|Q|$:
the charge-one kinetic/mass terms bound $m|Q_\Phi|$, the charge-two
$4m^2|X|^2$ terms bound $m|Q_X|$, and all added energy is nonnegative.
This extension is an admitted classical scalar production candidate, not
yet the polarization medium or a cosmologically viable theory. The portal
mass-gap sensitivity, quantum corrections, actual thermal Higgs phase,
and dark kinetic equations remain to be evaluated.

\subsection{Second obstacle: momentum and dark reprocessing}
In the rest frame of a Higgs, the source boost is $\gamma_X=m_h/(2M)$.
For isotropic source decay,
\begin{equation}
 E_\Phi=\gamma_X(M/2+\beta_X m u_*\cos\theta),\qquad
 \langle E_\Phi\rangle=m_h/4.
\end{equation}
The benchmark gives $\gamma_X=31243.8$ and daughter energies
$30.625$--$31.875$ GeV. Rest-frame coldness is not cosmological birth coldness.
Early redshifting could nevertheless cool these daughters; the boost alone
does not exclude the model.

For a deliberately collisionless monochromatic control produced at
$T_p=m_h/3$, visible entropy conservation gives
\begin{equation}
 a_p=\frac{T_{\gamma0}}{T_p}(g_{s0}/g_{sp})^{1/3},\quad
 a_{\rm nr}=a_p(p_p/m),\quad
 v(a)/c=\frac{a_{\rm nr}}{\sqrt{a^2+a_{\rm nr}^2}}.
\end{equation}
Approximating the entire interval as radiation dominated gives the exact
control integral
\begin{equation}
 \lambda_{\rm fs,RD}=\frac{c a_{\rm nr}}{H_0\sqrt{\Omega_r}}
 \left[\operatorname{asinh}(a_{\rm eq}/a_{\rm nr})
 -\operatorname{asinh}(a_p/a_{\rm nr})\right].
\end{equation}
With the declared standard-cosmology inputs in the script, this is
$4.3728\,10^{-4}$ Mpc and $p/m=1.9890\,10^{-7}$ at equality.
These are kinematic controls, not a transfer function or valid small-scale
structure constraint. The source thermal spectrum and changing degrees of
freedom are not evolved in this control.

It is also not self-consistent to retain collisionless evolution with the
adopted cross section. At visible $T=1$ keV, well below the source resonance,
using the total benchmark number density, isotropic equal-speed control
$\langle v_{\rm rel}\rangle=4v/3$, and applying the threshold cross section
to that mixture gives
\begin{equation}
 \frac{\Gamma_{22}}H=\frac{Y_\Phi s(T)\sigma_0\langle v_{\rm rel}\rangle}{H(T)}
 \simeq1374.
\end{equation}
The particle/antiparticle channel matrix must be derived; even retaining only
the same-sign half-density estimate leaves a rate much greater than $H$.
This rejects the collisionless continuation of this control, not cold
interacting matter. Momentum-resolved dark collisions, inverse conversions
and allowed number-changing reactions must replace it. The required equation is
\begin{equation}
 (\partial_t-Hp\partial_p)f_i(p,t)
 =C_i^{\rm SM}+C_i^{1\leftrightarrow2}+C_i^{2\leftrightarrow2}+cdots,
\end{equation}
with reciprocal collision integrals and energy/charge conservation.
The Higgs bridge now has an explicit tree-stable scalar extension. Its
collisional kinetics and cosmological phase history remain unresolved.

\section{Route B: an exact exchange--acceleration consistency test}
Define $Q$ in mass-density per time, positive from dark energy into cold
matter. With conserved baryons, constant $G,c,\kappa$, and the scale link
in equation~\eqref{eq:scale},
\begin{align}
 \dot\rho_c+3H\rho_c&=Q,\\
 \dot\rho_\DE+3H(1+w)\rho_\DE&=-Q,\\
 2\frac{\dd\ln a_0}{\dd\ln a}
 &=-3(1+w)-\frac{Q}{H\rho_\DE},\\
 \frac{\dd\ln(\rho_c/\rho_b)}{\dd\ln a}&=\frac{Q}{H\rho_c}.
\end{align}
Eliminating $Q$ gives the \textbf{DERIVED} conditional null relation
\begin{equation}
 \boxed{2\frac{\dd\ln a_0}{\dd\ln a}+3(1+w)
 +\frac{\rho_c}{\rho_\DE}
 \frac{\dd\ln(\rho_c/\rho_b)}{\dd\ln a}=0.}
\label{eq:null}
\end{equation}
Proof: differentiate $a_0^2\propto\rho_\DE$, insert the second continuity
equation, then substitute the fourth equation. No cosmological solution
or microscopic choice of $Q$ is required for that algebraic implication.
The code checks 27 rational cases as controls of the implementation.

For a vacuum-like component $w=-1$, exactly constant $a_0$ forces background
$Q=0$. A homogeneous net transfer from cold matter into vacuum raises
$\rho_\DE$ and $a_0$; transfer in the other direction lowers them.
The interacting-vacuum equations are known \cite{wands}; that source defines
$Q$ with the opposite sign. Our equation is their continuity system combined
with this programme's postulated scale link, not a new general cosmology.

Important scope: $Q$ here is a cosmic mean net exchange, not automatically
the local loss of halo kinetic energy. Halo cooling may deposit energy in
dark radiation or other excitations instead of vacuum. Spatially local
exchange, local gravitational work, and cosmological averaging need their own
covariant model. A constant $T_v^{\mu\nu}=-\epsilon_v g^{\mu\nu}$ has zero
covariant divergence. A varying vacuum can exchange energy but requires
$\nabla_\mu T_v^{\mu\nu}=-\nabla^\nu\epsilon_v$ and a compatible receiver.
For variable $G$ or $\kappa$, equation~\eqref{eq:null} acquires the negative
terms $-\dd\ln G/\dd\ln a-2\dd\ln\kappa/\dd\ln a$ on its left.

\textbf{Next discriminator:} derive the actual $Q^\mu$ of the proposed
formation receiver, determine its background average, and evolve abundance
and $a_0$ together. A measurable inconsistency in equation~\eqref{eq:null}
would reject the conjunction of the scale link and exchange assumptions,
not select a particle identity.

\section{Route C: can exterior free mass repair interior lensing?}
This tests one tempting escape from the inverse free-mass constraint:
place all additional positive matter outside the region where the baryonic
law must be preserved. In ordinary spherical weak-field mass lensing,
\begin{equation}
 \Sigma(R)=2\int_R^\infty\frac{\rho(r)r\dd r}{\sqrt{r^2-R^2}},\qquad
 \Delta\Sigma(R)=\frac2{R^2}\int_0^R\Sigma(S)S\dd S-\Sigma(R).
\end{equation}
For a thin shell of mass $M_s>0$ and radius $r_s>R$, put
$x=R/r_s$ and $s_x=\sqrt{1-x^2}$. Direct projection gives
\begin{align}
 \Sigma_s&=\frac{M_s}{2\pi r_s^2s_x},\quad
 M_{\rm proj}(<R)=M_s(1-s_x),\\
 \Delta\Sigma_s
 &=\frac{M_s}{\pi r_s^2}
 \left[\frac{1-s_x}{x^2}-\frac1{2s_x}\right]
 =-\frac{M_s x^2}{2\pi r_s^2s_x(1+s_x)^2}<0.
\end{align}
Newtonian interior acceleration from the shell is zero, but its excess
tangential shear is negative. Its convergence is positive; these are
different observables. By positive superposition, every finite nonnegative
spherical distribution supported at $r\ge r_0$ has
$\Delta\Sigma_{\rm ext}(R)\le0$ for $0<R<r_0$, whenever the integrals exist.
It is strictly negative for nonzero exterior mass at finite radii.
The small-$R$ limit is
\begin{equation}
 \Delta\Sigma_s=-\frac{M_sR^2}{8\pi r_s^4}+O(R^4).
\end{equation}
That leading expansion was already calculated in
\path{campaign_fresh_gravity/CFG3_principle.py}, lines 297--302.
The all-interior sign and positive-superposition application are formal
corollaries, not an independently novel force law.

The script checks five exact rational shell cases and three positive shell
mixtures. Therefore placing all repair mass beyond a protected radius cannot
increase a positive $\Delta\Sigma$ residual inside that radius while leaving
interior dynamics unchanged. This does not reject repair at larger projected
radii, nonspherical environments, interpenetrating matter, correlated
line-of-sight structure, or a different relativistic metric theory.
The next empirical step remains the measured baryon distribution and
isolation-selected lens--matter correlation required by the October 15 work
order. We have not rerun a lensing likelihood or changed an empirical verdict.

\section{Relativity, damping, and what remains to close}
The necessary joined action has the schematic form
\begin{equation}
 S=S_{\rm EH}+S_{\rm SM}+S_c+S_{\rm response}
 +S_{\rm bridge}+S_{\rm receiver}+S_{\rm counterterms}.
\end{equation}
This is an inventory, not a completed action. In a common $c=1$ convention,
\begin{equation}
 \nabla_\mu T_i^{\mu\nu}=Q_i^\nu,\qquad \sum_i Q_i^\nu=0,
 \qquad \nabla_\mu T_{\rm tot}^{\mu\nu}=0.
\end{equation}
A portal producing particles supplies neither the MOND constitutive law nor
its photon metric. In the usual weak-field convention,
\begin{equation}
 \dd s^2=-(1+2\Phi_g/c^2)c^2\dd t^2
 +a^2(1-2\Psi_g/c^2)\dd\mathbf x^2.
\end{equation}
Nonrelativistic gravity probes $\Phi_g$ and lensing probes $\Phi_g+\Psi_g$.
A separately added force on baryons cannot be assumed to bend light.
No same-action relativistic kinetic/interface proof is supplied here.

The existing Newtonian carrier starting point is
\begin{equation}
 \partial_t f+\mathbf v\cdot\nabla_x f
 -\nabla\Phi_g\cdot\nabla_v f=C[f],\quad
 \rho_c=m\int f\dd^3v,\quad
 \nabla^2\Phi_g=4\pi G(\rho_b+\rho_c).
\end{equation}
Elastic scattering redistributes energy; it cannot erase the total formation
energy. Radiation, escaping matter, a reciprocal bath, or phase mixing must
be distinguished. Conservative phase mixing can yield weak coarse settling
without strong convergence to a zero-energy resting minimum. No chosen
relaxation time is promoted to a derived clock.

\begin{longtable}{p{.20\textwidth}p{.36\textwidth}p{.35\textwidth}}
\toprule Obligation & What exists & First unresolved implication\\\midrule
\endhead
O0 source dictionary & Dated candidate equations and input hashes & One common positive-mass, charge and force mapping\\
O1 reciprocal action & Independent carrier, static and current constructions & Microscopic fields generating the adopted polarization law\\
O2 metric health & Scoped static hybrid/Einstein branches in prior work & Constrained dynamics, interfaces and lensing for the joined model\\
O3 formation & Prior finite-stock energy calculations & Actual receiver, rate and basin from independent starts\\
O4 partition & Tether/free prescription; ownership algorithm & Physical phase variable and reciprocal transition\\
O5 cosmology & Production control here; inherited cosmology in older tests & Stable potential, collision integrals, native transfer functions\\
O6 parameters & Fitted $\kappa$, stock, bounded $L_*$ & Mechanism or explicit independently measured input inventory\\
O7 joint tests & Candidate-specific successes and failures & Same candidate across all domains, both footings\\
O8 sealed data & Existing unchanged commitments & Verify public availability and execute frozen scoring when available\\
\bottomrule
\end{longtable}

The current fourth-round lensing calculation passes its relative criterion
but fails its absolute criterion: point-source $\chi^2=165.36/52$ nominal
degrees of freedom; even four fitted baryon radii give $75.73/48$ with
nominal $p=.006516<.01$. These are recorded diagnostics, not recomputed here.
They do not become successes by adding a mathematical stability theorem.
Likewise the prior static Einstein branch does not inherit the hybrid
model's static positivity theorem or prove dynamical stability.

\section{Research decisions and novelty boundary}
\textbf{Priority 1: the stabilized production candidate with full dark kinetics.}
Use the completed-square scalar extension above, determine its thermal
history, then solve the charged momentum distributions including
inverse conversion and the correct scattering channels. Success would be
an admitted potential and reproducible abundance/coldness; the source mass,
portal and initial history remain inputs unless independently fixed.
If the original 1-MeV benchmark is retained, use the explicit potential
repair and replace the failed collisionless continuation. A different mass
point requires new screening.

\textbf{Priority 2: a physical formation receiver.}
Derive reciprocal $Q^\mu$ and check energy, charge and momentum; if vacuum
is the receiver, impose equation~\eqref{eq:null}. If another component is
the receiver, propagate its perturbations and observational consequences.
This route remains open because no receiver has been specified by this report.

\textbf{Priority 3: one finite mass model for rotation and lensing.}
Use measured extended baryons, real-mass positivity, a counted stock,
the inverse constraint, and the exterior-shell sign. Predict the
isolation-selected cross-correlation rather than a fitted matter-auto-power
surrogate. The shell-only interior-repair construction is excluded in its
stated domain; broader finite population constructions remain open.

Particle production through a Higgs portal is known. Dipolar matter joining
dark energy, cold cosmological behavior and MOND-like polarization also
predates this programme \cite{blanchet}. The exchange null relation is a
formal combination of known continuity with this programme's scale postulate.
The shell result follows from standard projection, with its leading term
already in this repository. The contribution of this checkpoint is the
specific joint screening and its actionable constraints. No claim here is
certified globally unique or novel.

Literature discovery used targeted arXiv/primary-publisher web searches on
Higgs freeze-in, decay momenta, dipolar polarization, interacting vacuum,
and exterior-shell shear, checked on 16 October 2026. Primary full text was
read for \cite{hall,wands,blanchet}; abstracts of adjacent clockwork and
two-field portal work were discovery aids only. A cache lookup for
\texttt{arxiv:0911.1120v2} had no match. No PDF was retained here. This was
not a complete ADS/INSPIRE citation-graph or global novelty search.

\section{Verification and continuation checkpoint}
\texttt{checks.py} uses Python 3.9.6 standard library only. The fresh
\texttt{run\_003} records 56 passing bounded controls, including 27 exact
exchange identities, five exact shell identities, three shell mixtures,
two quadratures of the radiation-only distance and nine exact scalar
completion identities. Exact derivations are
given above; passing finite cases is not their universal proof.
The first run preceded the explicit negative-bare-mass check; run 002 added
that check, and run 003 executes the stable-quartic continuation. Earlier
runs are retained with their historical input hashes. Use run 003 for the
current script. Review is self-review, not an independent referee report.

The remaining goal is exactly O1--O7 above, plus the unchanged observational
commitments as data become available. Next execution starts with the
stabilized portal's thermal history and charged collision operators, or with a
specified reciprocal formation receiver. None is currently running in this
checkpoint. The routes are deferred/open, not exhausted and not marked solved.

\begin{thebibliography}{9}
\bibitem{hall} L. J. Hall, K. Jedamzik, J. March-Russell and S. M. West,
\emph{Freeze-In Production of FIMP Dark Matter}, arXiv:0911.1120v2
(8 December 2009), especially section 6.1, equations (6.6)--(6.10).
\url{https://arxiv.org/html/0911.1120v2}.
The decay-yield formula is specialized here with four final daughters.

\bibitem{wands} C. Kaeonikhom, H. Assadullahi, J. Schewtschenko and D. Wands,
\emph{Observational constraints on interacting vacuum energy with linear
interactions}, JCAP 01 (2023) 042, arXiv:2210.05363,
equation (2.7). \url{https://arxiv.org/pdf/2210.05363}.
The web PDF available on the check date was read; its revision was not pinned
locally. Only the displayed continuity equations, with reversed $Q$ sign,
are used, not its observational bounds.

\bibitem{blanchet} L. Blanchet and A. Le Tiec,
\emph{Model of Dark Matter and Dark Energy Based on Gravitational Polarization},
Phys. Rev. D 78, 024031 (2008), arXiv:0804.3518v2,
introduction and action framework.
\url{https://arxiv.org/html/0804.3518v2}.
This establishes overlap of the broad mechanism, not the validity of our action.
\end{thebibliography}

\clearpage

\part{Thermal phase and charge-resolved collision audit}

\begin{abstract}
This continues the explicit scalar Higgs/carrier construction at checkpoint
CPB-2026-10-16. A leading high-temperature Higgs-bath comparison finds an
intermediate instability of the source field at the abundance-calibrated
1-MeV benchmark. An exact sufficient inequality avoids dark condensation in
the stated class of positive bath-mass comparisons; a separately dated
mass-family continuation satisfies it. Vacuum tree amplitudes show that the
portal-produced neutral mixture scatters more strongly than the previous
one-channel estimate. Time dilation also invalidates a prompt-source-decay
shortcut in a separate vacuum control. Neither a physical thermal history,
a relic abundance, the galaxy law, nor full theory closure is established.
\end{abstract}

\section{Continuation contract and model}
The previous goal turn was progress: it produced an explicit reciprocal
scalar potential and identified unresolved thermal/kinetic conditions.
The actual checkout remains
\texttt{370dfb58006e330f210f5e306297e822188e37ba}, with shared uncommitted
work. This continuation reads the actual prior report, its current run 003,
and the October 15 work order. No frozen theory, seal, prior empirical score,
or other lane is modified. The new authoritative computation is
\texttt{thermal\_charge\_run\_002/manifest.json}. This is self-review.

Use $\hbar=c=k_B=1$, $(+---)$, and the same complex scalars $\Phi,X$ with
charges $1,2$. Write $x=|X|$, $n=|\Phi|^2$, $t=H^\dagger H$, and
$w_H=t-v_H^2/2$. The previously specified potential is
\begin{align}
 V_J={}&m^2|\Phi|^2+M^2|X|^2
 +g_c(X^*\Phi^2+X\Phi^{*2})+\lt|\Phi|^4\\
 &+\lambda_Hw_H^2+\lp w_H|X|^2+\lambda_X|X|^4,\nonumber\\
 M^2={}&4m^2(1+u_*^2),\quad \Delta_c=4m^2u_*^2,\quad
 \lt=\frac{g_c^2}{\Delta_c}+\lambda_{\rm rep}.
\end{align}
The positive-square carrier identity and
$\lambda_X\ge\lp^2/(4\lambda_H)$ establish the previous global
zero-temperature bound. They do not prove that the same origin is stable
at every temperature. In particular,
\begin{equation}
 M_0^2=M^2-\lp v_H^2/2
\end{equation}
is the source curvature before the Higgs obtains its vacuum expectation value.

The benchmark fixes $m=10^{-3}$ GeV, $u_*=.02$, $r_\lambda=2$,
$\sigma_0/m=.3\ \mathrm{cm^2/g}$, $\lambda_X=10^{-18}$ and
$\lp=2.21149\,10^{-10}$. The portal was calibrated in the earlier
simplified production control to $\Omega_c h^2=.12$; it is not a derived
abundance. $g_c\simeq2.168\,10^{-6}$ GeV and $\lt\simeq.00881$ follow
from these stipulated microscopic inputs.

\section{A thermal comparison that does not assume a dark heat bath}
Assume the Standard Model bath is thermal, initially no populated dark bath,
and retain only its leading high-temperature real scalar self-energies.
Specify the comparison free energy
\begin{equation}
 V_T=V_J+c_H T^2t+\frac{\lp T^2}{6}x^2,
 \qquad
 c_H=\frac{3g_2^2+g_Y^2}{16}+\frac{\lambda_H}{2}+\frac{y_t^2}{4}.
\label{eq:VT}
\end{equation}
We do not insert $\lt T^2$ or $\lambda_XT^2$ as dark thermal masses:
the dark distributions have not been shown to equilibrate at the visible
temperature. Those corrections would depend on the actual populated states.

The portal tadpole normalization is independently fixed by
\begin{equation}
 \langle H^\dagger H\rangle_T=
 \frac12\sum_{a=1}^4\langle h_a^2\rangle_T=\frac{T^2}{6}
\end{equation}
for four massless real Higgs components with
$\langle h_a^2\rangle_T=T^2/12$. This is the same normalization as
Seller, Sz\'ep and Tr\'ocs\'anyi, Appendix C, equations (C.11)--(C.12)
\cite{seller}: their complex singlet is our $X$, their doublet is our $H$,
and their quartic mixing coefficient is our $\lp$. Their reverse singlet-bath
contribution to the Higgs mass is omitted here because there is no assumed
dark thermal population. The visible-bath diagram remains present even
when the dark field is out of equilibrium.

With illustrative fixed $g_2=.65$, $g_Y=.36$, $y_t=.94$,
$\lambda_H=m_h^2/(2v_H^2)$, $m_h=125$ GeV and $v_H=246$ GeV,
$c_H=.372768$. The comparison Higgs origin changes stability at
\begin{equation}
 T_H=\sqrt{\lambda_Hv_H^2/c_H}=144.77\ \mathrm{GeV}.
\end{equation}
This is a mean-field control scale, not a new prediction for the physical
Standard Model crossover. Running, thermal resummation, field-dependent
thermal integrals, imaginary parts in unstable regions, and non-equilibrium
dark backreaction have not been solved. High-temperature formulas must not
be extrapolated to a few GeV as physical decay-threshold predictions.
Phase and thermal-rate effects in freeze-in are known issues \cite{deromeri}.

\section{Analytic elimination and the baseline instability}
Minimize over the relative scalar phase so that the cubic term is
$-2g_cxn$. At each fixed $x\ge0$ the two other invariants minimize exactly at
\begin{equation}
 n_*(x)=\max\!\left(\frac{2g_cx-m^2}{2\lt},0\right),\quad
 t_*(x,T)=\max\!\left(\frac{\lambda_Hv_H^2-c_HT^2-\lp x^2}
 {2\lambda_H},0\right).
\label{eq:elim}
\end{equation}
Thus the entire minimization of the stipulated homogeneous free energy is
one-dimensional and piecewise quartic. The boundaries are
$x=m^2/(2g_c)$ and, when real,
$x=\sqrt{(\lambda_Hv_H^2-c_HT^2)/\lp}$.
Interior stationary points obey
\begin{equation}
 4a x^3+2b x+c=0.
\end{equation}
For $t_*>0$,
$a=\lambda_X-\lp^2/(4\lambda_H)$ and
$b=M^2+\lp[1/6-c_H/(2\lambda_H)]T^2$.
For $t_*=0$, $a=\lambda_X$ and $b=M_0^2+\lp T^2/6$.
If $n_*>0$, replace $b$ by $b-g_c^2/\lt$ and set
$c=g_cm^2/\lt$; otherwise $c=0$.
Checking all admissible roots, branch boundaries and $x=0$ suffices since
the quartic coefficients are positive. The implementation enumerates those
roots and separately checks a logarithmic amplitude grid.

In the unbroken Higgs branch the source's linear curvature is
\begin{equation}
 M_T^2=M_0^2+\lp T^2/6.
\end{equation}
At the benchmark $M_0^2=-2.68992\,10^{-6}$ GeV$^2$. At $T=200$ GeV,
\begin{equation}
 M_T^2=-1.21559\,10^{-6}\ \mathrm{GeV}^2<0.
\end{equation}
Both visible Higgs and daughter have positive quadratic curvatures at that
origin, so this source instability is an actual negative scalar direction
of the specified comparison, not an artifact of minimizing them out.
A canonical source fluctuation satisfies, locally,
\begin{equation}
 \ddot{\delta X}+3H\dot{\delta X}
 +(k^2/a^2+M_T^2)\delta X=0.
\end{equation}
For long wavelengths $\sqrt{-M_T^2}/H=1.9621\,10^{10}$ with the stated
radiation-era $H$. The exact classical field at precisely zero remains at
zero without a seed; the result is linear instability to admitted nonzero
fluctuations, not spontaneous production proved by a deterministic zero solution.

The stipulated high-temperature source-origin instability occupies
\begin{equation}
 T_H<T<\sqrt{-6M_0^2/\lp}=270.15\ \mathrm{GeV}.
\end{equation}
The exact lower edge of the full comparison's dark phase can extend below
$T_H$; it is not asserted here as a physical transition temperature.
Varying $c_H$ over $.3,.4,.5$ and the source tadpole coefficient over
$\lp/12,\lp/6,\lp/3$ preserves a nonempty unbroken-Higgs instability
interval in all nine stipulated cases. This is coefficient sensitivity,
not a precision resummed calculation.

At 200 GeV the minimized comparison gives
\begin{equation}
 |X|=7.79785\,10^5\ \mathrm{GeV},\quad
 |\Phi|=13.8502\ \mathrm{GeV},\quad H=0,\quad
 \Delta V_T=-3.69742\,10^5\ \mathrm{GeV}^4.
\end{equation}
These are homogeneous mean-field extrema and a free-energy difference,
not a simulated condensate, a released-energy budget, or a relic yield.
Both dark amplitudes are nonzero; the vacuum $U(1)$ is broken on that
comparison branch and the original uncondensed particle spectrum cannot
be reused there. Whether the universe follows this branch, fragments,
restores the symmetry or produces acceptable relics is unresolved.
The earlier zero-temperature positivity proof remains correct. What fails
is promoting it to an unbroken dilute thermal production history.

\section{A sufficient condition for avoiding the dark phase}
There is a useful analytic continuation beyond that failed shortcut.
Suppose $\lp,\lambda_X,\lt>0$, the only dark bath correction is a
nonnegative source-mass correction, and
\begin{equation}
 \boxed{M_0^2\ge g_c^2/\lt.}
\label{eq:sufficient}
\end{equation}
For all $x,n,t\ge0$ the dark part before that bath correction is
\begin{align}
 &m^2n+M_0^2x^2-2g_cxn+\lt n^2+\lp tx^2+\lambda_Xx^4\\
 &\quad=\lt\left(n-\frac{g_cx}{\lt}\right)^2
 +m^2n+\left(M_0^2-\frac{g_c^2}{\lt}\right)x^2
 +\lp tx^2+\lambda_Xx^4\ \ge0.\nonumber
\end{align}
The relative phase can only raise this minimum. For each Higgs field
configuration, $X=\Phi=0$ attains the dark minimum. Therefore no lower
dark-condensed equilibrium exists in this whole specified class, even if
the visible thermal potential is replaced by an arbitrary bounded function
of $t$ with an attained minimum. This is a sufficient inequality, not a
necessary condition, and not a theorem including uncomputed dark thermal
or quantum corrections.

To execute a changed continuation, hold $u_*,r_\lambda,\sigma_0/m$ and the
earlier illustrative final abundance fixed while varying $m$. The prescribed
relations are
\begin{equation}
 \lambda_{\rm rep}=\sqrt{8\pi m^3(\sigma_0/m)},\quad
 g_c^2=\lambda_{\rm rep}\Delta_c/r_\lambda,\quad
 \lp(m)\propto m^{-1/2}(1-4M^2/m_h^2)^{-1/4}.
\end{equation}
The last is abundance calibration within the old decay-yield approximation;
it is not a cosmologically verified prescription. Substitution into
equation~\eqref{eq:sufficient} gives the illustrative sufficient floor
$m\simeq1.22840$ MeV. This number is neither a derived mass nor an observed
particle bound.

\begin{center}
\begin{tabular}{r r l}
\toprule $m$ (MeV) & $M_0^2$ (GeV$^2$) & Condition~\eqref{eq:sufficient}\\\midrule
0.75 & $-5.47580\,10^{-6}$ & not met\\
1.00 & $-2.68992\,10^{-6}$ & not met\\
1.25 & $2.67422\,10^{-7}$ & met\\
2.00 & $1.12748\,10^{-5}$ & met\\
10.0 & $3.98044\,10^{-4}$ & met\\\bottomrule
\end{tabular}
\end{center}
The 1.25-, 2- and 10-MeV rows have no lower dark-condensed state at the
nine tested temperatures; the all-temperature sufficient statement follows
from the square identity, not that finite grid. These are separately
specified mass-family candidates, not a rescue declared for the old point.
They still require thermal rates, full kinetics, empirical tests and the
connection to gravity. More massive candidates can also change independent
laboratory/cosmological constraints, which have not been scored here.

\section{The neutral population has two distinct scattering channels}
The Higgs portal produces $XX^*$ and hence equal positive/negative daughter
populations in the stated charge-symmetric scenario. The conserved charge
is a difference, not their total rest mass. The vacuum tree vertices from
$-\lt|\Phi|^4-g_c(X^*\Phi^2+\mathrm{h.c.})$ are $-4i\lt$ and $-2ig_c$.
With on-shell external daughters and invariant $s$,
\begin{align}
 \mathcal M_{++}(s)&=-4\left[\lt+\frac{g_c^2}{s-M^2}\right],\\
 \mathcal M_{+-}(u)&=-4\left[\lt+\frac{g_c^2}{u-M^2}\right].
\end{align}
Same-charge scattering has an $s$-channel source; the crossed
particle--antiparticle channel has $u$-channel exchange. There is no
same-vertex $X$ $t$-channel because it would violate charge assignment.
These are tree amplitudes; the physical near-pole same-charge channel needs
a width and inclusive accounting. The crossed channel has no such pole for
the physical elastic $u\le0$ used here.

For COM momentum $p$, $s=4(m^2+p^2)$ and
$u=-2p^2(1+\cos\theta)$. For distinguishable opposite-charge daughters,
\begin{equation}
 \sigma_{+-}(s)=\frac{1}{2\pi s}\int_{-1}^1
 \left[\lt-\frac{g_c^2}{M^2+2p^2(1+z)}\right]^2\dd z.
\end{equation}
For identical same-charge final particles the $1/2!$ factor instead gives
\begin{equation}
 \sigma_{++}(s)=\frac{1}{2\pi s}
 \left[\lt+\frac{g_c^2}{s-M^2}\right]^2.
\end{equation}
At threshold,
\begin{align}
 \sigma_{++}(0)&=\frac{\lambda_{\rm rep}^2}{8\pi m^2}\equiv\sigma_0,\\
 \sigma_{+-}(0)&=\frac{(\lt-g_c^2/M^2)^2}{4\pi m^2},\\
 \frac{\sigma_{+-}(0)}{\sigma_0}
 &=2\left[1+\frac{1}{r_\lambda(1+u_*^2)}\right]^2=4.49880056.
\end{align}
The threshold rate per test daughter in an equal mixture of total density
$n_{\rm tot}$ and equal isotropic speed distributions is
\begin{equation}
 \Gamma_{\rm el}=\frac{n_{\rm tot}}{2}
 (\sigma_{++}+\sigma_{+-})\langle v_{\rm rel}\rangle.
\end{equation}
It is $2.74940028$ times the old screen $n_{\rm tot}\sigma_0
\langle v_{\rm rel}\rangle$. Therefore the old 1-keV monochromatic
collisionless control has $\Gamma_{\rm el}/H\simeq3776.90$, rather than
1373.72, in the threshold limit. This refines a hypothetical control and
does not supply the evolved spectrum. These are total collision cross
sections, not automatically a viscosity or momentum-transfer cross section
to insert into an empirical halo bound. Angular weighting needs its own
definition, especially for identical particles.

Nine exact rational threshold identities and numerical angular integrations
at $p/m=0,.00025,.003,.01,.05$ check normalization and continuity. They do
not compute the dense-medium inelastic collision operator.

\section{A short resting lifetime does not prove prompt cosmological decay}
The earlier isolated source lifetime is about $3.52\,10^{-13}$ seconds.
For a monochromatic Higgs initially at rest, $\gamma_{X,p}=m_h/(2M)$,
and the production-era decay rate is reduced by that boost. Keeping vacuum
masses and constant radiation degrees of freedom solely as a separate control,
\begin{equation}
 \frac{\Gamma_X}{\gamma_{X,p}H_p}=.0245331\quad(T_p=m_h/3).
\end{equation}
This is not prompt on the production Hubble scale. If $r=a/a_p$,
$u_{X,p}=\sqrt{\gamma_{X,p}^2-1}$ and $H=H_p/r^2$, the survival law is
\begin{equation}
 P_X(r)=\exp[-\tau(r)],\qquad
 \tau(r)=\frac{\Gamma_X}{H_p}\int_1^r
 \frac{s^2}{\sqrt{s^2+u_{X,p}^2}}\dd s.
\end{equation}
Numerical quadrature gives $\tau=1$ at $r=4.97701$ or $T=8.37182$ GeV,
with the source still highly relativistic. These numbers omit thermal
mass shifts and dark collisions and are not a physical decay epoch for
the unstable thermal benchmark.
While the source is ultrarelativistic, delayed decay approximately preserves
the mean comoving daughter momentum, since $aE_X$ is approximately constant.
Thus this delay alone does not overturn the earlier kinematic redshift
estimate. It does invalidate using a prompt cascade in the rate evolution.

\section{An explicit number-changing amplitude, not just an allowed charge}
The stronger check is the actual reaction
\begin{equation}
 2\Phi+2\Phi^*\longrightarrow\Phi+\Phi^*.
\end{equation}
Take all four incoming particles at rest. Each final particle has energy
$2m$ and momentum magnitude $\sqrt3m$. In the specified scalar theory,
with $\lambda_{H\Phi}=0$ at tree level, the complete six-daughter tree
graphs have three topologies: nine graphs with two daughter quartics,
18 with a daughter quartic and two cubic source vertices, and nine with
four cubic source vertices. The source quartic and Higgs portal cannot
occur in a six-daughter tree at this order: with trivalent/quartic vertices
$n_3+2n_4=4$, those choices leave too few external daughters or unmatched
Higgs/source legs. This does not omit them from loops or other processes.

The sum can be reconstructed without a current recursion. In units $m=1$,
write $d=\Delta_c/m^2$, $L=M^2/m^2+2=6+d$ and $\widehat g=g_c/m$.
An external same-charge pair has source denominator $-d$ if both are
incoming or $-L$ if one is a crossed outgoing particle. Daughter triple
denominators are $8$ or $-4$. For the three topologies the sums of inverse
propagator products are respectively
\begin{equation}
 S_0=-\frac34,\qquad S_2=\frac{3}{2L},\qquad
 S_4=-\frac{1}{4d^2}+\frac{1}{2dL}-\frac{1}{L^2}.
\end{equation}
For $S_0$, four triples give denominator $8$ and five give $-4$.
For $S_2$, endpoint pairs of invariant $4$ cancel after summing their
three allowed attachments; the other pairs give the displayed $3/(2L)$.
For $S_4$, the two endpoint pairs can be incoming/incoming once,
incoming/crossed in four cases, or crossed/crossed in four cases,
giving the three displayed terms. The products of vertex factors are
$16\lt^2$, $16\lt\widehat g^2$, and $16\widehat g^4$, with the
common amplitude sign $-1$. Restoring units,
\begin{equation}
 \mathcal M_{4\to2}=\frac1{m^2}\left[
 12\lt^2-\frac{24\lt\widehat g^2}{L}
 +\widehat g^4\left(\frac4{d^2}-\frac8{dL}+\frac{16}{L^2}\right)
 \right].
\end{equation}
Using $\widehat g^2=\lambda_{\rm rep}d/r_\lambda$ and
$z=d/L\in(0,1)$ gives a useful universal simplification:
\begin{equation}
 \boxed{\mathcal M_{4\to2}
 =\frac{\lambda_{\rm rep}^2}{m^2}
 \left[12+\frac{24(1-z)}{r_\lambda}
 +\frac{16(1-z)^2}{r_\lambda^2}\right].}
\label{eq:six}
\end{equation}
Consequently this vacuum tree amplitude is strictly nonzero for
$m,\Delta_c,g_c^2,\lambda_{\rm rep}>0$. The bracket is greater than 12.
This is a statement about the specified tree theory throughout that
parameter family, not a statement that chemical equilibration occurs.
It also does not establish the in-medium amplitude near a changing phase.

At $u_*^2=1/2500$, $r_\lambda=2$, its bracket is exactly
$393885012/14070001\simeq27.9947$. The computation uses exact rational
momenta and polynomials in $g_c/m$. A labeled off-shell-current recursion
and a separate direct enumeration of the 36 graph topologies agree at
12 rational parameter points. Re-rooting the external labels, the analytic
four-point amplitudes, and the pure-quartic six-point limit supply controls.
The positive closed form proves nonvanishing; the finite checks only
verify those implementations. This remains self-review.

For a dilute nonrelativistic uncorrelated gas, the zero-incoming-momentum
event rate of this one channel is
\begin{equation}
 R_{2+2-\to+-}=\frac{\sqrt3\,|\mathcal M_{4\to2}|^2}
 {1024\pi m^4}n_+^2n_-^2.
\end{equation}
The factors are the initial $1/(2!2!)$, four initial measures $1/(2m)$,
and the distinguishable final two-body phase space $\sqrt3/(16\pi)$.
Since two particles are lost per event, a symmetric gas of total density
$n$ has one-channel loss coefficient
\begin{equation}
 \Gamma_N\equiv\frac{2R}{n}
 =\frac{\sqrt3\,|\mathcal M_{4\to2}|^2 n^3}{8192\pi m^4}.
\end{equation}
Using the earlier stipulated density in a threshold control gives
$\Gamma_N/H=1.07\,10^{-28}$ at visible $T=1$ keV and
$1.24\,10^{-6}$ at 1 MeV. The latter has control $p/m\simeq.35$ and
cannot be read as a precision threshold rate, especially with the narrow
source scale $u_*=.02$. Neither number checks relativistic early production,
inverse gains, other charge channels or medium effects. They show why
``number change exists'' and ``number change is fast'' must be kept separate.

\section{The kinetic system now required}
The actual species are $\Phi,\Phi^*,X,X^*$, the visible bath, and any
coherent phase/response fields. In an uncondensed quasiparticle phase,
\begin{equation}
 (\partial_t-Hp\partial_p)f_i=C_i^{\rm portal}
 +C_i^{1\leftrightarrow2}+C_i^{2\leftrightarrow2}
 +C_i^{\rm number\ change}.
\end{equation}
An isotropic reduction still retains momentum dependence. With
$\dd\Pi_j=\dd^3p_j/[(2\pi)^32E_j]$, a bosonic elastic term has the form
\begin{align}
 C_a(p_1)={}&\frac{1}{2E_1 S_f}\int\dd\Pi_2\dd\Pi_3\dd\Pi_4
 (2\pi)^4\delta^4(p_1+p_2-p_3-p_4)|\mathcal M|^2\\
 &\ \times[f_3f_4(1+f_1)(1+f_2)-f_1f_2(1+f_3)(1+f_4)],\nonumber
\end{align}
where for the displayed tagged-daughter elastic channels $S_f=2!$ for
same-charge identical final daughters and $S_f=1$ for opposite-charge
final daughters. Other processes need their own tagged-species multiplicities.
The source collision term is
\begin{align}
 C_X(p_X)={}&\frac{1}{2E_X\,2!}\int\dd\Pi_1\dd\Pi_2
 (2\pi)^4\delta^4(p_X-p_1-p_2)\,4g_c^2\\
 &\ \times[f_1f_2(1+f_X)-f_X(1+f_1)(1+f_2)].\nonumber
\end{align}
The tagged daughter gain/loss term has the corresponding factor of two
relative to this event count. $X^*$ and the negative-charge daughters have
the conjugate terms. This is how inverse conversion, stimulation and charge
balance enter; a fixed decay rate without inverse terms is insufficient.

The classical global symmetry conserves
\begin{equation}
 q_\Phi=1,\ q_{\Phi^*}=-1,\ q_X=2,\ q_{X^*}=-2,\qquad
 \sum_i q_i\int\frac{\dd^3p}{(2\pi)^3} C_i=0.
\end{equation}
It does not conserve the total number of positive plus negative daughters.
The preceding section proves a nonzero mixed-charge six-point tree amplitude.
Its full relativistic/in-medium collision kernel and the other number-changing
channels have not been calculated, so neither total negligibility nor
chemical equilibrium may be asserted. In chemical equilibrium this reaction
imposes $\mu_\Phi+\mu_{\Phi^*}=0$; a charge-symmetric state then has zero
chemical potential for both. Elastic equilibration alone does not impose that.
The earlier all-positive charge-energy protection cannot silently be used
as a conservation theorem for a neutral population's rest abundance.
In particular, the earlier regulated pressure bound was formulated for the
positive-charge nonrelativistic sector with $d=\langle D_c\rangle/(m\langle Q\rangle)$.
For a charge-symmetric cosmological population $Q=0$ this quantity is
undefined and the bound is not a proof of continued coldness. A new
particle/antiparticle energy and pressure calculation is required. The
reaction above conserves total energy by converting some rest mass into
final kinetic energy; it supplies no energy from nothing.

For fixed masses, the sum of collision energy moments vanishes when every
participating bath/species is included. If quasiparticle masses depend on
the thermal/condensate background, each density also has a mass-work term:
\begin{equation}
 \dot\epsilon_i+3H(\epsilon_i+p_i)
 =\int\frac{\dd^3p}{(2\pi)^3}E_i C_i
 +\int\frac{\dd^3p}{(2\pi)^3}\dot E_i f_i.
\end{equation}
Mean-field and bath stresses must supply the reciprocal work without
double-counting the self-energy. Minima of $V_T$ are not by themselves
conserved mechanical energy or a production calculation.

This also sharpens the dark-energy question. A thermal scalar free-energy
offset is not automatically the observed cosmic vacuum. Any proposal that
net energy enters or leaves that vacuum must derive its $Q^\mu$ and obey
the previous checkpoint's exchange--acceleration consistency relation.
No equation here fixes the vacuum magnitude, $\kappa$, or the polarization
law. A successful production history would close one input to the dark
mass problem, not the gravitational and lensing completion.

\section{What changed and the next discriminating work}
The zero-temperature scalar bridge now has a sharper decision tree.
The original 1-MeV comparison cannot simply start from a thermally stable
empty dark origin and retain the earlier dilute cascade. Its intermediate
phase is a possible production route, but it requires its own non-equilibrium
field/particle calculation. A distinct sufficient-condition mass-family
route avoids that particular phase problem and can proceed to kinetic
production, while keeping its parameters explicitly selected.

Next execute the candidate's collision integrals with full charge channels,
finite-width intermediate states, thermal kinematics and number-changing
rates. For the symmetric mass-family branch, first verify thermal mass/rate
corrections and dark reprocessing before recalibrating abundance. For the
1-MeV phase branch, first derive the condensate and fluctuation evolution,
including the bath work and the changed spectrum. A stable equilibrium
table is not a substitute for either calculation.

Computational evidence: 108 bounded controls in the thermal/elastic run, with
five mass choices and nine temperatures, exact rational threshold identities,
stationary-root/log-grid comparisons, angular quadrature refinement, and
a time-dilated decay control. Its run 001 preceded the added decay control and
retains historical input hashes. The separate number-changing run 002
adds 120 exact/finite controls, including the independent topology sum;
its run 001 preceded that independent enumeration. The universal
sufficient-condition and amplitude derivations are written above; finite checks do not
prove physical validity. No simulations or external jobs remain running.

Literature scope: primary full text of \cite{seller} and the thermal section
of \cite{deromeri} inspected on 16 October 2026. A project-cache query for
the latter had no match. No PDFs retained. Broad singlet phase transitions,
thermal freeze-in and charge crossing are known; the present result is a
specific conditional screen for the current scalar construction. No complete
citation-graph search or global novelty claim is made.

\begin{thebibliography}{9}
\bibitem{seller} K. Seller, Z. Sz\'ep and Z. Tr\'ocs\'anyi,
\emph{Real effective potentials for phase transitions in models with extended
scalar sectors}, JHEP 04 (2023) 096, arXiv:2301.07961v2,
Appendix C.1, equations (C.11)--(C.12).
\url{https://arxiv.org/html/2301.07961v2}.
Only leading visible-bath thermal normalizations and the need to distinguish
them from a full resummed phase calculation are imported.
\bibitem{deromeri} V. De Romeri, D. Karamitros, O. Lebedev and T. Toma,
\emph{Neutrino dark matter and the Higgs portal: improved freeze-in analysis},
arXiv:2003.12606v2, sections 2.1.1 and 6.2.4 and Appendix A.
\url{https://arxiv.org/html/2003.12606v2}.
Its sterile-neutrino model is not identified with this scalar model;
the source supplies methodological overlap, not our abundance or phase verdict.
\end{thebibliography}

\clearpage

\part{Pair-conversion evolution and pressure-aware exchange}

\begin{abstract}
The scalar carrier contains a pole-free pair-conversion reaction that must
be included alongside source decay and inverse decay. Their combination
changes the eventual daughter count while conserving charge and energy.
We derive its tree cross section and thermal rate, and execute an expanding
Maxwell--Boltzmann gas comparison using fast source partial equilibrium.
One specified initial state nearly doubles its weighted daughter count;
a hotter comparison fails the preset dilution gate and is retained as such.
The evolution supplies concrete kinetic evidence, not a primordial production
history, a relic-density fit or a closed gravitational theory. A pressure-aware
generalization of the cold--vacuum exchange identity separates internal
rest-mass conversion from exchange with dark energy.
\end{abstract}

\section{Exact target and provenance}
The full user goal remains a physical completion of the theory, a defensible
particle bridge, and a complete account of cold and dark energy. This
checkpoint advances the missing collision/abundance calculation. It does
not replace that goal with a chosen initial-value problem.

The previous goal turn was progress: it established a thermal obstruction,
a sufficient alternative scalar branch, and a mixed-charge six-point
threshold amplitude. This turn reads the actual previous records before
extending the collision system. The checkout remains
\texttt{370dfb58006e330f210f5e306297e822188e37ba}, with shared dirty state.
Inputs and software are pinned by
\texttt{pair\_conversion\_run\_004/manifest.json}. The frozen theory and
other lanes remain unchanged. This is self-review.

We choose the separately specified $m=2$ MeV family member, not the earlier
thermally unstable 1-MeV point. It satisfies the previous sufficient
condition $M_0^2\ge g_c^2/\lambda_{\rm tot}$ within that thermal comparison.
The choices $u_*=.02$, $r_\lambda=2$ and
$\sigma_0/m=.3\ \mathrm{cm^2/g}$ remain inputs. The source mass is
$M=2m\sqrt{1+u_*^2}$, and $g_c,\lambda_{\rm tot}$ are the same derived
parameter values in the pinned mass-family record. Natural units apply.

\section{A missing reaction and an exact cross section}
The charge-preserving scalar interaction
\begin{equation}
 \mathcal L\supset-g_c(X^*\Phi^2+X\Phi^{*2})
 -\lambda_{\rm tot}|\Phi|^4
\end{equation}
allows
\begin{equation}
 \Phi+\Phi^*\ \rightleftarrows\ X+X^*.
\label{eq:pair}
\end{equation}
The cubic vertex is $-2ig_c$. For the forward process, one internal daughter
joins the two vertices; fixed charges give a single $t$-channel tree:
\begin{equation}
 \mathcal M=-\frac{4g_c^2}{t-m^2}.
\end{equation}
There is no crossed identical-final graph because $X$ and $X^*$ are distinct.
This amplitude is not the same as opposite-charge daughter elastic scattering.

At invariant $s\ge4M^2$ put
$p=\sqrt{s/4-m^2}$, $k=\sqrt{s/4-M^2}$ and
\begin{equation}
 t=m^2+M^2-s/2+2pk\cos\theta.
\end{equation}
Both initial and final energies are $\sqrt s/2$, so
$t=-|\mathbf p-\mathbf k|^2\le0$. The internal daughter is never on shell.
With $A=s/2-M^2$, $B=2pk$,
\begin{equation}
 A^2-B^2=M^4+m^2(s-4M^2)>0,\qquad
 \int_{-1}^1\frac{\dd z}{(A-Bz)^2}=\frac2{A^2-B^2}.
\end{equation}
Direct angular integration therefore gives
\begin{align}
 \sigma_f(s)&=\frac{g_c^4}{\pi s}
 \sqrt{\frac{s-4M^2}{s-4m^2}}
 \frac1{M^4+m^2(s-4M^2)},\\
 \sigma_r(s)&=\frac{g_c^4}{\pi s}
 \sqrt{\frac{s-4m^2}{s-4M^2}}
 \frac1{M^4+m^2(s-4M^2)}.
\end{align}
They obey
$(s-4m^2)\sigma_f=(s-4M^2)\sigma_r$.
There are no initial/final identical-pair factorials for either direction.
The reverse cross section diverges kinematically as $1/k$ at its threshold,
but its physical velocity product has the finite limit
\begin{equation}
 (\sigma_r v_{\rm M\o l})_0
 =\frac{g_c^4}{2\pi M^6}\sqrt{1-m^2/M^2}
 =9.50486\,10^{-8}\ \mathrm{GeV}^{-2}
\end{equation}
for the selected point. Angular quadrature and detailed-balance controls
check the normalization independently of the closed angular formula.

\section{Why a two-to-two reaction changes the final rest count}
Use a charge-symmetric state, with per-sign densities $n=n_\Phi=n_{\Phi^*}$
and $b=n_X=n_{X^*}$. Let $D$ be net source decay rate per sign and $P$
net forward pair rate. After visible production has ceased,
\begin{equation}
 \dot n+3Hn=2D-P,\qquad
 \dot b+3Hb=-D+P.
\end{equation}
The useful weighted count per sign is $w=n+2b$, the daughter count after
all remaining sources decay. It satisfies
\begin{equation}
 \boxed{\dot w+3Hw=P.}
\end{equation}
Forward pair conversion consumes one daughter of each sign and creates
one source of each sign; subsequent decays produce four daughters instead
of the initial two. The inverse cycle reduces the eventual daughter count.
The net conserved charge stays zero. Source decay alone conserves $w$;
the pair channel does not.

For a common Maxwell--Boltzmann kinetic temperature $T_d$, write fugacities
$z_\Phi,z_X$ relative to zero chemical potential. Detailed balance gives
\begin{equation}
 D=\gamma_d^\eq(z_X-z_\Phi^2),\qquad
 P=\gamma_p^\eq(z_\Phi^2-z_X^2),\qquad
 \gamma_p^\eq=(n_X^\eq)^2\langle\sigma_r v_{\rm M\o l}\rangle.
\end{equation}
Fast source decay/inverse decay imposes $z_X=z_\Phi^2$. Full chemical
equilibrium would additionally require $z_\Phi=1$; it is not imposed.
The collision structure, not dimensional analysis or the conserved charge,
determines whether and when that second equilibrium is approached.

The standard invariant reaction measure and flux definitions are checked
against De Romeri et al. equations (21)--(25) \cite{rates}. Our use of
classical distributions is an explicit reduction, not their Bose treatment.
For equal-mass distinguishable incoming sources the MB thermal average is
\begin{equation}
 \langle\sigma_rv_{\rm M\o l}\rangle=
 \frac{1}{8M^4T_d K_2(M/T_d)^2}
 \int_{4M^2}^\infty\sigma_r(s)(s-4M^2)\sqrt{s}
 K_1(\sqrt{s}/T_d)\dd s.
\end{equation}
We integrate with $y=(\sqrt s-2M)/T_d$ and exponentially scaled Bessels.
A generalized Laguerre weight $\sqrt y e^{-y}$ removes the threshold
square root. Across 25 temperatures in $T_d/m\in[10^{-5},10]$, 64- and
128-node rates differ by at most $3.45\,10^{-7}$ relative. The nonrelativistic
limit agrees with the displayed velocity product.

\section{Reciprocal energy closure without a rest-energy cancellation}
The gas is isolated from the visible sector during this comparison. Its
total dark energy and pressure satisfy
\begin{equation}
 \dot\epsilon_d+3H(\epsilon_d+p_d)=0.
\end{equation}
Pair production can increase rest mass only by reducing kinetic energy or
the other stored dark energy; it supplies no new total energy.
The thermodynamics uses an ideal quasiparticle equation of state. Interaction
and virial contributions to that equation of state are not included here.

For numerical stability define $\delta=M-2m$ and the reduced fugacity
$y=z_\Phi e^{-m/T_d}$. Write
\begin{align}
 A&=\frac{m^2T_d}{2\pi^2}e^{m/T_d}K_2(m/T_d),\\
 B&=\frac{M^2T_d}{2\pi^2}e^{M/T_d}K_2(M/T_d)e^{-\delta/T_d},\\
 n&=Ay,\quad b=By^2,\quad
 y=\frac{2w}{A+\sqrt{A^2+8Bw}}.
\end{align}
The last expression is the positive quadratic solution and stays well
conditioned for exponentially small equilibrium densities. With partial
source equilibrium the pair rate is
\begin{equation}
 P=B^2\langle\sigma_rv\rangle y^2
 [e^{-2m/T_d}-y^2].
\end{equation}
It can have either sign along the expanding trajectory.

For $e_i=m_iK_1(m_i/T_d)/K_2(m_i/T_d)+3T_d$, put
\begin{equation}
 k_\Phi=e_\Phi-m,\quad k_X=e_X-2m,\quad
 k=nk_\Phi+bk_X,\quad p=(n+b)T_d.
\end{equation}
Here $\epsilon_d=2(mw+k)$ and $p_d=2p$. Subtracting weighted rest energy
gives the exact moment identity within this closure,
\begin{equation}
 \dot k+3H(k+p)=-mP.
\end{equation}
If $P>0$, increasing the eventual rest count costs the kinetic/internal
budget. $P<0$ releases that rest energy into the remaining gas.

Let $c_i=\dd e_i/\dd T_d>0$. At fixed temperature and fixed weighted count,
respectively,
\begin{align}
 k_w&=\frac{nk_\Phi+2bk_X}{n+4b},\\
 k_T&=nc_\Phi+bc_X+
 \frac{nb(2k_\Phi-k_X)^2}{(n+4b)T_d^2}>0.
\end{align}
With $\ell=\ln a$ the actual two-variable evolution is
\begin{align}
 \frac{\dd\ln Y_w}{\dd\ell}&=\frac{P}{Hw},\\
 \frac{\dd T_d}{\dd\ell}&=
 \frac{-3p+3nb(k_X-2k_\Phi)/(n+4b)-(P/H)(m+k_w)}{k_T},\qquad
 Y_w=\frac{2w}{s_{\rm SM}}.
\end{align}
No rest-mass term is discarded. In the daughter-only elastic limit these
equations reduce to fixed $Y_w$ and $T_d\propto a^{-2}$ when nonrelativistic.
Thus rapid elastic scattering alone need not heat a homogeneous cold gas.
It changes the distribution and transport; pressure evolution follows its
energy moment. Number conversion changes that conclusion and is calculated
here explicitly.

There is also an independent thermodynamic check. The MB entropy density
per sign is
\begin{equation}
 s_d=\frac{k}{T_d}+n+b-w\ln y,\qquad
 \mu_\Phi=m+T_d\ln y,\quad \mu_X=2\mu_\Phi.
\end{equation}
The first law $\dd(mw+k)=T_d\dd s_d+\mu_\Phi\dd w$ and its Euler
identity give
\begin{equation}
 \dot s_d+3Hs_d=-\frac{\mu_\Phi}{T_d}P\ge0.
\end{equation}
Indeed $P=\gamma_p^\eq(z_\Phi^2-z_\Phi^4)$ and
$\mu_\Phi/T_d=\ln z_\Phi$ have opposite signs. Pair-off trajectories
conserve comoving entropy; pair-on trajectories produce it. This is a
conditional MB reaction-network entropy statement, not a claim about
uncomputed quantum kinetics.

\section{The executed expanding-gas comparison}
The initial visible temperature is 10 MeV and the final visible temperature
is 1 keV. The background contains a photon/electron ideal gas and three
massless neutrinos, with instantaneous neutrino decoupling at 2 MeV.
Electron Fermi integrals are included; visible entropy and the neutrino
temperature are evolved consistently. This is an inherited radiation
background, not the framework's native cosmology or precision neutrino
decoupling calculation. Its 500-row entropy table conserves the declared
entropy to the checked tolerance.

The specified initial weighted yield is
\begin{equation}
 Y_{w,i}=\frac{.12}{2.742\,10^8(m/\mathrm{GeV})}
 =2.18818\,10^{-7}.
\end{equation}
It is the old abundance-calibrated input, not a completed prediction from
the Higgs era. Initial source densities follow the assumed fast partial
equilibrium. Three independent initial kinetic temperatures are investigated;
none is claimed to have been produced by the earlier portal calculation.

\begin{center}
\begin{tabular}{r r r l}
\toprule Initial $T_d/m$ & Final $Y_w/Y_{w,i}$ & Final $T_d$ (GeV) & Dilution gate\\\midrule
.3, pair off & 1 & $5.59317\,10^{-12}$ & met\\
.3, pair on & 1.03195 & $9.58651\,10^{-12}$ & met\\
1, pair off & 1 & $4.21977\,10^{-11}$ & met\\
1, pair on & 1.97924 & $1.67032\,10^{-11}$ & met\\
3, pair off & 1 & $3.46622\,10^{-10}$ & met\\
3, pair on & 4.52622 & $3.39091\,10^{-11}$ & failed\\\bottomrule
\end{tabular}
\end{center}
Every number is a result for this reduced initial-value problem. In
particular, $1.97924$ is not an observed relic excess or a physical Higgs
production prediction. Turning the pair reaction off keeps the weighted
count exactly fixed but retains the same expansion/source closure.

The initial-energy ceilings on the yield ratios are 1.57214, 3.36915 and
9.15433, respectively. They follow from nonnegative pressure and
$\epsilon_d\ge2mw$: comoving energy cannot increase in the isolated gas.
The computed yield peaks 1.09428, 2.30842 and 6.02160 remain below those
ceilings. The subsequent decline is reciprocal reverse conversion and
freeze-out, rather than arbitrary deletion of mass.

For pair-on runs the independent expansion-work residual
\begin{equation}
 \frac{a_f^3\epsilon_{d,f}-\epsilon_{d,i}
 +\int_0^{\ell_f}3a^3p_d\dd\ell}{\epsilon_{d,i}}
\end{equation}
is below $5\,10^{-10}$ in magnitude. The weighted-number collision moment
is independently integrated as well. The refined $.3$ run has energy and
number residuals below $7\,10^{-12}$. ODE/rate refinement changes its final
yield by less than $10^{-9}$ relative. These numerical invariants verify
the implemented reduced equations, not the completeness of the collision set.
The pair-on final comoving entropy ratios are 1.12636, 1.62315 and 2.98564.
Their separately integrated entropy-production balances agree to better
than $1.4\,10^{-8}$ relative to initial entropy. Pair-off entropy ratios
equal one within $7\,10^{-10}$.

\section{Validity diagnostics and preserved failures}
The original script, pinned in failed run 001, required maximum MB
occupation $y<.02$, source-decay/evolution rate ratio $>10^5$, and a lower
bound on elastic/evolution rate ratio $>10$. These thresholds are retained.
The pair-on maximum occupations are .0059944, .0114825 and .0229111.
The hot case fails the original dilution gate. It remains an approximate
classical comparison and is not promoted to a result meeting that gate.
The failure means the declared small-occupation criterion is missed; it is
not evidence for a condensate or a universal failure of a classical expansion.

All three pair-on runs have elastic lower bounds at least 500 times their
largest local evolution scale under the chosen normalization, and source
decay/evolution ratios above $7.5\,10^8$. The elastic bound uses the
pole-free daughter opposite-charge inequality
\begin{equation}
 \sigma_{+-}(s)\ge\frac{(\lambda_{\rm tot}-g_c^2/M^2)^2}{\pi s}.
\end{equation}
These are quantitative support for the kinetic and fast-decay reductions,
not a solution of the full momentum-space collision operator. The
occupation/evolution checks are sampled on the stored 401-point trajectories;
no global-in-time certificate is claimed from those samples.

Runs 001 and 002 stopped on the hot-case dilution assertion; run 002 records
the measured diagnostics before stopping. Run 003 preserves the same
threshold, exports all computed trajectories and reports the failed scope
gate separately from numerical execution checks. It does not loosen the
gate to manufacture a pass. The final authoritative run 004 records
585 numerical/algebraic controls, including 500 background-entropy checks.
Run 003 remains the pre-entropy checkpoint. This count is not a
measure of physical completeness.

\section{What the full reaction system still needs}
The explicit pair process and source inverse decays already account for
real intermediate sources in a reaction network. At nonzero daughter
momenta, parts of the previous six-point amplitude can develop on-shell
source propagators and overlap those sequential processes. Adding a raw
finite-width $4\leftrightarrow2$ rate on top would count those histories twice.
The remaining collision hierarchy needs consistent real-intermediate-state
subtraction or an off-shell kinetic derivation. The general issue is
documented by Ala-Mattinen et al. \cite{ris}; their subtraction prescription
has not been transferred to this multiparticle scalar system here.

The preceding exact zero-incoming-momentum six-point result remains valid:
its intermediate source invariants lie below the positive-gap source pole.
It neither supplies the finite-temperature subtracted remainder nor bounds
it at the warmer initial states evolved in this report. Those omitted
rates have not been bounded sufficiently to establish the displayed yield
ratios as predictions of the complete scalar action.

Other remaining obligations are the real Higgs-era source and thermal
history, the resulting initial distributions, Bose factors and medium
dispersion/width corrections, and the full kinetic response where timescale
reductions fail. At sufficiently late nonrelativistic temperatures there
is no reason to add a fictitious visible-temperature dark thermal mass;
it must be calculated from the actual states.

\section{Cold energy exchange is not automatically dark-energy exchange}
Internal conversion here moves energy between particle rest mass and
kinetic energy while the total dark gas obeys its conservation equation.
It does not require the vacuum to provide or absorb energy. In a bound
galaxy, cosmological expansion is not a substitute for a physical local
formation receiver; the prior halo-settling obligation remains open.

The previous cold--vacuum consistency identity assumed a dust component.
For a component with nonzero $w_c=p_c/(\rho_cc^2)$, mass-equivalent total
energy density $\rho_c$, conserved baryons, and $Q$ positive from dark
energy into that component, the correct equations are
\begin{align}
 \dot\rho_c+3H(1+w_c)\rho_c&=Q,\\
 \dot\rho_{\rm DE}+3H(1+w_{\rm DE})\rho_{\rm DE}&=-Q,\\
 a_0^2&=\kappa^2c^2G\rho_{\rm DE}.
\end{align}
With constant $G,c,\kappa$, elimination of $Q$ gives
\begin{equation}
 \boxed{2\frac{\dd\ln a_0}{\dd\ln a}+3(1+w_{\rm DE})
 +\frac{\rho_c}{\rho_{\rm DE}}
 \left[\frac{\dd\ln(\rho_c/\rho_b)}{\dd\ln a}+3w_c\right]=0.}
\end{equation}
The proof is direct substitution of the two continuity equations after
differentiating the scale postulate. For this isolated gas, $Q=0$ even
when its eventual rest count changes. Treating $m$ times that count as
the entire $\rho_c$ during the warm conversion epoch would falsely infer
vacuum exchange. The old dust identity is recovered at $w_c=0$.
This distinction is necessary for a coherent cold/dark-energy story.
It derives neither the scale postulate nor $\kappa$, and fixes no vacuum
energy magnitude.

\section{Next route and completion status}
The executed result changes the next research decision: source production
cannot be converted into final daughter abundance by multiplying by four
and stopping. The pair channel, inverse decay and finite-temperature
irreducible number change must be evolved with reciprocal energy exchange.
The initial energy distribution can affect the final count substantially.

Next derive the real-intermediate-state-consistent irreducible scalar
collision kernel and its error relative to the explicitly evolved network.
Then replace the specified initial states with the stabilized candidate's
actual Higgs-era production and thermal history, and compute perturbations
from that same state. A portal recalibration before those steps would hide
the unresolved abundance mechanism rather than close it.

The galaxy-response action, finite counted halo formation, source partition,
lensing metric, $\kappa$, independent parameter evidence and common-model
empirical tests remain unresolved. The unique/novel-path ambition remains
subject to literature comparison: reaction-network reprocessing and dark
temperature evolution are known mechanisms. This report establishes their
specific consequences for the present chosen scalar comparison, not global
novelty or full closure.

\begin{thebibliography}{9}
\bibitem{rates} V. De Romeri, D. Karamitros, O. Lebedev and T. Toma,
\emph{Neutrino dark matter and the Higgs portal: improved freeze-in analysis},
arXiv:2003.12606v2, equations (21)--(25) and the discussion of thermal masses.
\url{https://arxiv.org/html/2003.12606v2}.
Only invariant rate/flux conventions and methodological overlap are used;
the quoted model's sterile-neutrino results are not imported.
\bibitem{ris} K. Ala-Mattinen, M. Heikinheimo, K. Kainulainen and K. Tuominen,
\emph{Anatomy of real intermediate state-subtraction scheme},
arXiv:2309.16615v2, sections I--III, especially the on-shell overlap in
the introduction and equations (1)--(4).
\url{https://arxiv.org/html/2309.16615v2}.
The source authenticates the double-counting issue; it does not supply a
completed multiparticle subtraction for this model.
\end{thebibliography}

\clearpage

\part{Source-assisted channels and pole matching}

\begin{abstract}
The previous pair/source abundance evolution omitted a pole-free
three-to-two reaction already present in the same scalar action.
We derive its complete tree amplitude and finite-temperature MB reaction
density. At the selected 2-MeV point and $T_d=m$, its two conjugate channels
add approximately 4.95 times the earlier pair-conversion coefficient.
Exact single, double and nested source-pole residues identify the reaction
network needed before subtracting the six-point remainder. A positive
one-pole interior matching construction has a proved smooth-profile error
bound; threshold-boundary and overlapping-pole matching remain unresolved.
The old abundance curves remain correct for their stated reduced model,
not verified predictions of the complete scalar theory. Full closure,
the galaxy-response connection and the vacuum normalization remain open.
\end{abstract}

\section{Current state and exact scope}
The previous goal turn made progress by deriving and executing a reduced
pair-conversion gas model. The current checkout is now
\texttt{bd5dbb23a2fa98f11f4bd9c13f1ef29970a06458}; its intervening commits
concern independent empirical lanes. The actual local kinetic files and
raw results were inspected before continuing. No same-route live job was
found in this checkpoint, and no earlier simulation was restarted.
All new computation handles have finished. Frozen model criteria and
other lanes are unchanged.

Use natural units and the scalar vertex conventions
\begin{equation}
 \mathcal L\supset-g(X^*\Phi^2+X\Phi^{*2})-\lambda|\Phi|^4,
 \quad q_\Phi=1,\quad q_X=2,
 \quad V_3=-2ig,\quad V_4=-4i\lambda.
\end{equation}
Here $\lambda=\lambda_{\rm tot}=\lambda_{\rm rep}+g^2/\Delta$,
$\Delta=M^2-4m^2>0$. The rate benchmark is the separately specified
$m=2$ MeV, $u_*=.02$, $r_\lambda=2$ mass-family candidate.
The equations below do not identify it with the static gravitational
polarization field. The source is treated as a narrow on-shell quasiparticle;
medium and Bose corrections are not included.

\section{The single-source cuts require more than a pair reaction}
The same action permits the two physically distinct crossings
\begin{align}
 X+\Phi^*+\Phi^*&\rightleftarrows\Phi+\Phi^*,\label{eq:SAI}\\
 X+\Phi+\Phi^*&\rightleftarrows\Phi+\Phi.\label{eq:SAJ}
\end{align}
Their conjugates are separate charge-related reactions. The first has
two identical initial daughters; the second has two identical final
daughters. Each needs one factorial $2!$ in its invariant event rate.
Neither was included in the prior pair/source evolution.

At the tree-amplitude level, cut an incoming same-charge pair from the
six-daughter reaction $2\Phi+2\Phi^*\to\Phi+\Phi^*$. With
$D_+=(p_1+p_2)^2-M^2$,
\begin{equation}
 \operatorname*{Res}_{D_+=0}\mathcal M_6=2g\mathcal M_5,
\end{equation}
where the right side is the physical source-assisted amplitude in
equation~\eqref{eq:SAI}. Cutting both incoming same-charge pairs gives
\begin{equation}
 \operatorname*{Res}_{D_+=D_-=0}\mathcal M_6
 =4g^2\mathcal M_{XX^*\to\Phi\Phi^*}.
\end{equation}
The nested pole of $\mathcal M_5$ itself has residue
$2g\mathcal M_{XX^*\to\Phi\Phi^*}$. A single real-source cut of the crossed
six-daughter reaction $3\Phi+\Phi^*\to2\Phi$ instead gives
equation~\eqref{eq:SAJ}. These identities follow by removing the cut tree
propagator and its cubic end vertex, with the stated amplitude sign convention.

The exact checks retain every labeled diagram. There are 36 six-point
trees and six one-source/four-daughter five-point trees. Single, double,
nested and crossed residues agree as exact polynomials in $g/m$ at nine
rational parameter points. External momenta are exactly on shell by using
rational spatial squared weights, rather than floating square roots.
The evidence is \texttt{source\_assisted\_matching\_run\_002}; its 72
exact controls support the implementation. The general factorization
argument is the tree cut, not a conclusion extrapolated from nine samples.
This is self-review.

The consequence for subtraction is precise: deleting only the double-source
pole and donating it to $XX^*\leftrightarrow\Phi\Phi^*$ leaves single-source
histories unmatched. The source-assisted collision channels must be supplied
before the residual six-point rate can be called fully matched.

\section{Complete five-point tree amplitude}
Take all external momenta incoming. Let $Q$ be the charge-two source,
$p$ the unique positive-charge daughter momentum, and
$r_0,r_1,r_2$ the three negative-charge daughter momenta. Then
\begin{equation}
 Q+p+r_0+r_1+r_2=0.
\end{equation}
Three trees have one quartic and one cubic vertex; three have three cubic
vertices. Summing them gives the compact exact expression
\begin{equation}
 \boxed{\mathcal M_5=-8g\sum_{k=0}^2
 \frac{\lambda+g^2/[(r_i+r_j)^2-M^2]}
 {(Q+r_k)^2-m^2},\qquad \{i,j,k\}=\{0,1,2\}.}
\label{eq:SAamp}
\end{equation}
For crossing~\eqref{eq:SAI}, one source denominator can be on shell:
the invariant mass of the two incoming negative daughters can equal $M$.
Its residue is the previously stated pair-scattering amplitude. A full
thermal rate for this crossing is therefore not obtained by simply squaring
the zero-width expression.

\section{The other crossing has no internal real-particle pole}
For equation~\eqref{eq:SAJ}, let $Q,p,r$ be future-directed incoming
source, positive daughter and negative daughter; let $k,l$ be the two
future-directed outgoing positive daughters. In equation~\eqref{eq:SAamp},
the negative momenta are $r,-k,-l$.

The source denominators obey
\begin{align}
 (k+l)^2-M^2&=s-M^2\ge4m(M+m)>0,\\
 (r-k)^2-M^2&\le-M^2,\qquad
 (r-l)^2-M^2\le-M^2.
\end{align}
For the daughter denominators,
\begin{equation}
 (Q+r)^2-m^2\ge M(M+2m)>0.
\end{equation}
The vector $K=Q-k=l-(p+r)$ is strictly spacelike. If $K$ were future
causal, $l=(p+r)+K$ would have mass at least $2m$, impossible. If $K$
were past causal, $k=Q-K$ would have mass at least $M>m$, also impossible.
$K=0$ would identify unequal on-shell masses. Thus
\begin{equation}
 (Q-k)^2-m^2<-m^2,\qquad (Q-l)^2-m^2<-m^2.
\end{equation}
The whole tree is pole-free throughout the physical region. This is a
uniform kinematic statement, not just a numerical absence of poles.

Since $\lambda>g^2/M^2$, an explicit uniform amplitude bound is
\begin{equation}
 |\mathcal M_5|\le C_5\equiv8|g|
 \left[\frac{\lambda+g^2/[4m(M+m)]}{M(M+2m)}
 +\frac{2\lambda}{m^2}\right].
\end{equation}
It also gives a finite positive analytic bound on the MB reaction density:
\begin{equation}
 \gamma_3^\eq\le\frac{C_5^2}{16\pi}I_X I_\Phi^2,
 \qquad I_i=\int\frac{\dd^3p}{(2\pi)^3 2E_i}e^{-E_i/T_d}
 =\frac{m_iT_d}{4\pi^2}K_1(m_i/T_d).
\end{equation}
The factor $1/(16\pi)$ is the maximal two-body phase space with its
identical-final $1/2!$. The computed rates respect this bound. It is loose
and does not replace quadrature-error assessment.

\section{Threshold amplitude and event normalization}
For all three incoming particles at rest in crossing~\eqref{eq:SAJ},
the final energies are $(M+2m)/2$. Evaluation of equation~\eqref{eq:SAamp}
gives
\begin{equation}
 \mathcal M_{5,0}=
 \frac{8g\lambda(M+m)}{Mm(M+2m)}
 -\frac{2g^3(5M+8m)}{mM^2(M+m)(M+2m)}.
\end{equation}
It is nonzero for the target range $2m<M<4m$, $g>0$,
$\lambda_{\rm rep}>0$. To see the sign, factor out
$8g/[Mm(M+2m)]$ and substitute
$\lambda=\lambda_{\rm rep}+g^2/(M^2-4m^2)$.
The remaining $g^2$ numerator is proportional to
\begin{equation}
 -x^3+24x+32=64+(4-x)(x^2+4x-8)>0,\qquad x=M/m\in(2,4).
\end{equation}
The $\lambda_{\rm rep}$ part is positive too. This establishes the
nonvanishing vacuum-tree threshold amplitude on that stated domain.

For a dilute nonrelativistic uncorrelated gas, the forward event rate of
one crossing~\eqref{eq:SAJ} is
\begin{equation}
 R_3=K_{3,0}n_Xn_\Phi n_{\Phi^*},\qquad
 K_{3,0}=\frac{\sqrt{1-4m^2/(M+2m)^2}\,|\mathcal M_{5,0}|^2}
 {128\pi M m^2}.
\end{equation}
The initial invariant measures give $1/[(2M)(2m)^2]$, and the final phase
space with the identical factor gives $\beta_f/(16\pi)$. Finite-temperature
amplitudes cannot be replaced by this threshold coefficient in a warm gas.

\section{A resolved finite-temperature kernel}
With classical equilibrium distributions and the final $1/2!$ included,
\begin{align}
 \gamma_3^\eq(T_d)&=\frac{T_d}{16\pi^3}
 \int_{(M+2m)^2}^\infty\sqrt{s}K_1(\sqrt{s}/T_d)\mathcal I_5(s)\dd s,\\
 \mathcal I_5(s)&=\frac{\beta_f}{1024\pi^4s}
 \int_{4m^2}^{(\sqrt{s}-M)^2}\!
 \sqrt{\mathcal K(s,M^2,u)}\sqrt{1-4m^2/u}
 \langle|\mathcal M_5|^2\rangle_{\rm angles}\dd u.
\end{align}
Here $\mathcal K$ is the K\"all\'en polynomial, $\beta_f=\sqrt{1-4m^2/s}$,
and the angular brackets are a normalized average over the initial relative
polar angle and final polar/relative azimuthal angles. The three-body phase
space follows by inserting the pair invariant $u$ and two two-body phase
spaces. The thermal total-momentum integral gives
$T_d\sqrt{s}K_1(\sqrt{s}/T_d)/(8\pi^3)$; the additional factor $1/2!$
produces the displayed coefficient. The invariant rate conventions agree
with the primary-source discussion in \cite{SAthermal}.

The code uses generalized Laguerre quadrature for total kinetic energy,
Jacobi quadrature for the initial Dalitz invariant, and polar Gaussian and
periodic azimuthal quadratures. Exponentially scaled Bessels avoid
underflow in $\widehat\gamma_3=\gamma_3^\eq e^{4m/T_d}$.
All internal denominators retain their proved physical signs.

The original coarse run failed its 2-percent refinement gate at $T_d=m$.
A separate resolution experiment found that refining energy/Dalitz nodes
alone changed the warm values negligibly, while the angular grid mattered.
The failed run remains intact. The authoritative refined run increases
angular resolution and retains the original gate. At $T_d=3m$, its two
last angular grids differ by .261 percent. These differences are finite
numerical evidence, not rigorous error bars. The low-$T_d/m=.003$ rate
approaches the threshold-density normalization within 2 percent, with
remaining finite-temperature corrections retained.

For source partial equilibrium $z_X=z_\Phi^2$ and equal charge populations,
the two conjugate source-assisted reactions add
\begin{equation}
 P_3=2\gamma_3^\eq(z_\Phi^2-z_\Phi^4)
\end{equation}
to $\dot w+3Hw$. The earlier pair contribution is
$P_2=\gamma_2^\eq(z_\Phi^2-z_\Phi^4)$. Consequently their coefficient
ratio is independent of the shared fugacity in this MB reduction:
\begin{center}
\begin{tabular}{r r r}
\toprule $T_d/m$ & $2\gamma_3^\eq/\gamma_2^\eq$ & $\gamma_3/\gamma_{3,\rm NR}$\\\midrule
.003 & .0017634 & .98001\\
.01 & .0096558 & .93529\\
.03 & .047251 & .82095\\
.1 & .25802 & .53614\\
.3 & 1.10193 & .19128\\
1 & 4.95265 & .016137\\
3 & 20.7964 & .00039732\\\bottomrule
\end{tabular}
\end{center}
In particular this omitted reaction is not justified as negligible at the
warmer temperatures used in the prior reduced abundance comparison. Its
threshold approximation would greatly overestimate the warm reaction.
These are vacuum-tree, common-MB coefficients; they are not observed rates
or a new primordial abundance calculation.

The exact residue and new kernel computations also explain why the prior
yield ratio 1.97924 must remain a result for its pair-only network. No
updated abundance ratio is asserted until the source-assisted channels,
remaining subtractions and applicable quantum/medium corrections have
been treated in the same evolution.

\section{A positive interior pole matching lemma}
The general real-intermediate-state problem is known \cite{SARIS}. The
following is a scoped construction for a single pole strictly inside the
physical integration domain. It is not a completed thermal prescription.

Write $D=s-M^2$, $a=M\Gamma>0$ and
\begin{equation}
 \mathcal A(D)=B(D)+\frac{R(D)}{D+ia},\qquad -W\le D\le W.
\end{equation}
Assume real $B,R$, nonnegative integration weight $v$, and define
$F=vR^2\in C^2$, $G=vRB\in C^1$ on the symmetric window. Let
$|F''|\le F_2$, $|G'|\le G_1$. Keep the full nonnegative squared
amplitude outside the window. Inside it keep the background $vB^2$ and
replace the resonant contribution by the positive local pole term
\begin{equation}
 \frac{2F(0)}a\arctan(W/a)
 =\frac{\pi F(0)}a Z_W,\qquad
 Z_W=\frac2\pi\arctan(W/a)\in(0,1).
\end{equation}
Equivalently the scattering amplitude uses the pole propagator only
outside the window, with the retained real-state spectral weight counted
separately. All retained squared-amplitude and pole terms are nonnegative.

The difference from the original full window integral is
\begin{equation}
 E_W=\int_{-W}^W\frac{F(D)-F(0)}{D^2+a^2}\dd D
 +2\int_{-W}^W\frac{G(D)D}{D^2+a^2}\dd D.
\end{equation}
Odd constant/linear pieces cancel, so Taylor's theorem and the mean-value
bound prove
\begin{equation}
 \boxed{|E_W|\le(F_2+4G_1)[W-a\arctan(W/a)].}
\end{equation}
This bounds the omitted regular interference as well as variation of the
resonant coefficient. If one instead uses the full unweighted pole term
$\pi F(0)/a$, its additional known difference
$[\pi-2\arctan(W/a)]F(0)/a$ must be counted. $Z_W$ is a matching
coefficient of the window, not an independently measured matter fraction.

The numerical benchmark uses $v(D)=e^{-D/L}$ and constant real $R,B$,
with $W=\sqrt{aL}$. Its exact bound has
$F_2=R^2 e^{W/L}/L^2$, $G_1=|RB|e^{W/L}/L$.
Twenty width/background cases satisfy the bound; a direct full-window
integration supplies an independent control at a moderate width.
At the selected vacuum source point,
\begin{equation}
 \Gamma=1.05750\,10^{-11}\ \mathrm{GeV},\quad
 a=4.23084\,10^{-14}\ \mathrm{GeV}^2.
\end{equation}
Assigning $L=\Delta/4$ solely as a scale diagnostic gives
$W=8.22760\,10^{-12}$ GeV$^2$,
$W/a=194.47$, $W/\Delta=.001286$, and $Z_W=.996726$.
The width/window/threshold-gap hierarchy is available in vacuum; the
actual thermal $F_2,G_1$ have not been bounded here.

The lemma fails to cover a pole window truncated by a phase-space endpoint,
complex backgrounds with absorptive interference, and overlapping poles
without further matching. In particular the upper initial-pair invariant
can approach $M^2$ near the source-pair threshold. A fixed symmetric window
cannot simply be used there. Those physical boundary and double-pole
patches remain open, rather than being hidden in the finite benchmark.

\section{What is established and what remains}
The current evidence consists of exact source-pole factorization,
a complete pole-free source-assisted amplitude, a bounded/resolved thermal
kernel, and the interior matching lemma. It supplies a concrete way to
advance the previously unspecified collision remainder. It does not supply
the whole remainder or make the earlier abundance curves complete.

Next add this pole-free kernel to the energy/entropy-conserving evolution,
with a sufficiently resolved temperature table and its error carried
through the result. In parallel as mathematical dependencies, match
crossing~\eqref{eq:SAI} and the six-point remainder using actual thermal
profile bounds, threshold endpoints and the exact pole residues. Do not
add an unsubtracted raw six-point rate to an explicit real-source network.
The additional reaction may change the dilution/source-equilibrium gates;
their old outcomes cannot be inherited without re-evolving.

This is a collision calculation for chosen massive fields. It neither
derives the Higgs-era initial state nor identifies the fields with the
galactic polarization medium. The vacuum offset, $\kappa$, the acceleration
scale postulate, physical halo settling, finite mass supply, source
selection, relativistic lensing and native cosmology remain unresolved.
Internal rest/kinetic conversion still is not automatically transfer to
dark energy. The pressure-aware exchange identity from the previous
checkpoint remains the applicable bookkeeping condition.

Novelty boundary: source-assisted number change, tree factorization,
thermal phase-space integration and real-state matching have established
prior art. This report gives their explicit application to the current
carrier and a scoped mathematical error estimate. No global novelty,
observational identification or independent-referee claim is made.

\begin{thebibliography}{9}
\bibitem{SAthermal} V. De Romeri et al.,
\emph{Neutrino dark matter and the Higgs portal: improved freeze-in analysis},
arXiv:2003.12606v2, invariant reaction measure and flux, equations (21)--(25).
\url{https://arxiv.org/html/2003.12606v2}.
The common-MB specialization and five-point phase space here are derived
for this scalar model; its sterile-neutrino results are not imported.
\bibitem{SARIS} K. Ala-Mattinen et al.,
\emph{Anatomy of real intermediate state-subtraction scheme},
arXiv:2309.16615v2, sections I--V.
\url{https://arxiv.org/html/2309.16615v2}.
The primary source was inspected for the on-shell overlap and cut matching
framework. It does not provide our physical boundary profiles or complete
multi-source subtraction.
\end{thebibliography}

\clearpage

\part{Enlarged evolution and the missing galaxy bridge}

\section{What this checkpoint establishes}
The physical theory is open. This checkpoint evolves the previously derived
pole-free source-assisted channel in the same expanding gas used for the pair
comparison. It also tests whether the chosen Higgs production portal could
itself supply the baryon-dependent galaxy response. The latter test finds
an extremely small direct local effect in this particular candidate. It
does not exclude a separate collective state or ordinary gravitating halos.
All calculations below are conditional on the selected 2-MeV scalar model;
they do not replace the frozen theory or fix its parameter origins.

\section{The enlarged reciprocal collision term}
Use $m=0.002\,\mathrm{GeV}$, $M=0.004000799920016\,\mathrm{GeV}$,
$g=7.29225381605\,10^{-6}\,\mathrm{GeV}$,
$\lambda_{\rm rep}=0.0166178017868$ and
$\lambda_{\rm tot}=0.0249267026802$. Let $n$ and $b$ denote the
per-sign daughter and source densities, and $w=n+2b$. Fast
$X\leftrightarrow\Phi\Phi$ and its conjugate imply
$z_X=z_\Phi^2$, not $z_\Phi=1$. The two conjugate pole-free reactions
\begin{equation}
 X+\Phi+\Phi^*\leftrightarrow 2\Phi,
 \qquad X^*+\Phi+\Phi^*\leftrightarrow 2\Phi^*
\end{equation}
each change the total weighted count by one per forward event. A forward
pair event changes it by two. Therefore the per-sign source in
$\dot w+3Hw=P$ is
\begin{equation}
 P=(\gamma_{2,\rm eq}+2\gamma_{3,\rm eq})(z_\Phi^2-z_\Phi^4).
 \label{ag:collision}
\end{equation}
The factor of two on the three-body density includes the conjugate channel
and the per-sign convention; no raw resonant six-point rate is added.

For numerical evaluation let $y=z_\Phi e^{-m/T_d}$ and
\begin{align}
 A&={m^2T_d\over2\pi^2}e^{m/T_d}K_2(m/T_d),\\
 B&={M^2T_d\over2\pi^2}e^{M/T_d}K_2(M/T_d)
       e^{-(M-2m)/T_d},\\
 n&=Ay,\quad b=By^2,\qquad
 y={2w\over A+\sqrt{A^2+8Bw}}.
\end{align}
If $K_{3,0}$ is the previously derived threshold event coefficient and
$h(\theta)=\gamma_{3,\rm eq}e^{4m/T_d}/(K_{3,0}BA^2)$,
then the cancellation-safe form is
\begin{equation}
 P=\left[B^2\langle\sigma_{XX^*\to\Phi\Phi^*}v\rangle
             +2K_{3,0}BA^2h(T_d/m)\right]
             y^2(e^{-2m/T_d}-y^2).
 \label{ag:scaled}
\end{equation}
In particular one must not absorb the second term into an effective
pair cross section proportional to $1/B$ at arbitrarily low temperature.

The thermal table has 47 points over $10^{-4}\leq T_d/m\leq3$,
including seven prior refined anchors; its nested coarse grid has 27 points.
Logarithmic PCHIP interpolation preserves positive rates. Below the table,
$h=\exp[(T_d/m)/10^{-4}\,\ln h(10^{-4})]$ tends to one. Replacing
that tail by one is a separate tested control. The earlier table run
calculated all points but failed JSON serialization of a NumPy Boolean;
the corrected run preserves the same mathematical calculation. This is
a numerical quadrature table, not a rigorous error enclosure.

\section{Energy, abundance and entropy}
Define $k_i=e_i-q_i m$, where $q_\Phi=1,q_X=2$ and
$e_i=m_iK_1(m_i/T_d)/K_2(m_i/T_d)+3T_d$.
Then $k=nk_\Phi+bk_X$, $p=(n+b)T_d$, and the total energy density
is $\epsilon=2(mw+k)$. With
\begin{align}
 k_w&={nk_\Phi+2bk_X\over n+4b},\\
 k_T&=nC_{V,\Phi}+bC_{V,X}
       +{nb(2k_\Phi-k_X)^2\over(n+4b)T_d^2}>0,
\end{align}
the evolution equations are
\begin{align}
 {\dd\ln Y\over\dd\ln a}&={P\over Hw},\\
 {\dd T_d\over\dd\ln a}&={-3p+3nb(k_X-2k_\Phi)/(n+4b)
                      -(P/H)(m+k_w)\over k_T}.
\end{align}
The per-sign entropy density satisfies
\begin{equation}
 s={k\over T_d}+n+b-w\ln y,\qquad
 \dot s+3Hs=-\left({m\over T_d}+\ln y\right)P\geq0.
\end{equation}
The number moment, expansion-work energy identity, and entropy production
are integrated independently as numerical controls.

Initial visible temperature is 10 MeV, final visible temperature is 1 keV.
The background is the specified ideal photon/electron gas with three
massless neutrinos decoupling instantaneously at 2 MeV. Initial weighted
yield $Y_i=0.12/(2.742\,10^8m)=2.18818\,10^{-7}$ is prescribed,
as is the initial dark temperature. Neither comes from a solved production
history. The outputs are:
\begin{center}
\begin{tabular}{rrrrr}
\toprule
$T_{d,i}/m$ & Pair-only $Y_f/Y_i$ & Enlarged $Y_f/Y_i$ & $\max y$ & MB gate\\
\midrule
0.3 & 1.031953 & 1.027186 & 0.006587 & pass\\
1.0 & 1.979244 & 1.958101 & 0.012369 & pass\\
3.0 & 4.526224 & 4.493585 & 0.024112 & fail\\
\bottomrule
\end{tabular}
\end{center}
The retained dilution gate is $y<0.02$. A numerical solver success does
not override its failure for the hottest state. For the three enlarged
trajectories the sampled elastic/evolution ratio exceeds 499 and the
source-decay/evolution ratio exceeds $6.7\,10^7$. These are sampled
timescale diagnostics, not a full momentum-space kinetic proof.

The nested temperature-grid change in final yield is below
$1.8\,10^{-5}$ fractionally in these comparisons. The $T_{d,i}/m=0.3$
joint ODE/pair-rate refinement changes final yield by less than $10^{-8}$;
the low-temperature-tail control changes it below $10^{-12}$. The run
contains 618 bounded controls, including 500 inherited background rows.
These controls do not bound omitted physical channels.
The added channel raises the peak count but reduces the eventual count
slightly: the reversible reaction alters the thermal and chemical history.
Its large instantaneous coefficient alone did not determine the final
abundance. Source-pole-containing assisted reactions, irreducible
six-point remainders, in-medium and Bose corrections remain absent.

\section{Why a production portal is not yet the galaxy bridge}
The tree-level Higgs portal is
$-\lambda_{HX}|H|^2|X|^2$, with
$\lambda_{HX}=1.56375810276\,10^{-10}$. At the chosen matching scale
$\lambda_{H\Phi}=0$ is an input, not an established all-orders property.
Writing $H=(0,(v+h)/\sqrt2)$, a slowly varying, nonrelativistic nucleon
medium has the Higgs source $-f_N\rho_b h/v$, where $\rho_b$ here is
rest-energy density in natural units. We use the historical benchmark
$f_N=0.308(18)$ from Ref.~\cite{ag:hoferichter}, not a claim to the latest
determination. At tree level and wavelengths much longer than $m_h^{-1}$,
\begin{equation}
 h\simeq-{f_N\rho_b\over v m_h^2},\qquad
 \delta M_X^2=\lambda_{HX}vh
            \simeq-{\lambda_{HX}f_N\rho_b\over m_h^2}.
 \label{ag:higgsresponse}
\end{equation}
This first changes the source mass and its exchange contact; it does not
linearly source a daughter field. The threshold same-charge contact is
\begin{align}
 \lambda_{\rm eff}(\rho_b)&=\lambda_{\rm tot}
                    -{g^2\over\Delta+\delta M_X^2},
       \qquad \Delta=M^2-4m^2,\\
 {\delta\lambda_{\rm eff}\over\lambda_{\rm rep}}
   &= {g^2\delta M_X^2\over
            \lambda_{\rm rep}\Delta(\Delta+\delta M_X^2)}
     \simeq {\delta M_X^2\over r_\lambda\Delta},
       \qquad r_\lambda=2.
\end{align}
For the chosen comparison density $0.1\,M_\odot\,\mathrm{pc}^{-3}$,
\begin{align}
 \rho_b&=2.91714\,10^{-41}\,\mathrm{GeV}^4,\\
 \delta M_X^2&=-8.99204\,10^{-56}\,\mathrm{GeV}^2,\\
 {\delta\lambda_{\rm eff}\over\lambda_{\rm rep}}
   &=-7.02503\,10^{-48}.
\end{align}
Even the chosen dense comparison $10^3\,M_\odot\,\mathrm{pc}^{-3}$
gives a fractional contact shift only $-7.02503\,10^{-44}$.
This excludes identifying this direct perturbative local contact response
with an order-one galactic response at these inputs. A putative collective
amplification would require an independently derived susceptibility, its
stable operating range and its relaxation time; naming a medium does not
supply that derivation. The required inverse local shift is of order
$10^{47}$ at the middle comparison density, if amplification of this
particular contact is the intended mechanism.

The nongravitational microscopic Compton lengths are
\begin{equation}
 {\hbar c\over m_h}=1.579\,10^{-18}\,\mathrm m,\qquad
 {\hbar c\over m}=9.866\,10^{-14}\,\mathrm m.
\end{equation}
A new linear Yukawa mediator reaching 1 kpc would require
$m_{\rm med}c^2\lesssim6.395\,10^{-27}\,\mathrm{eV}$ by the
range criterion. This statement concerns massive microscopic exchange;
it does not remove the long-range metric field or constrain a derived
gapless collective mode by its constituent mass.

On the empty dark-vacuum branch $X=\Phi=0$, the dark classical equations
have no linear baryonic source from this quadratic portal. With ordinary
Einstein gravity and a vacuum constant, the spherical weak-field solution
has acceleration
\begin{equation}
 g_r=-{GM_b\over r^2}+{\Lambda c^2 r\over3}.
\end{equation}
It does not yield $|g|\simeq\sqrt{GM_ba_0}/r$. This is a statement
about this empty branch, not a theorem forbidding occupied halos,
different curvature couplings or a collective phase.

\section{A condensate cannot simply be assumed for this halo gas}
As a separate ideal-gas comparison choose total rest-energy density
$0.3\,\mathrm{GeV\,cm}^{-3}$, an equal neutral mixture, and
one-dimensional velocity dispersion $\sigma_v=150\,\mathrm{km\,s}^{-1}$.
If these are equilibrium particles with $T=m\sigma_v^2$,
\begin{equation}
 n_\Phi={\rho_{\rm halo}\over2m},\qquad
 \lambda_T={\sqrt{2\pi}\hbar\over m\sigma_v},\qquad
 n_\Phi\lambda_T^3=9.05709\,10^{-21}.
\end{equation}
The ideal three-dimensional Bose critical occupation is
$\zeta(3/2)=2.61238$. With those fixed density and velocity inputs the
mass giving this occupation would be about $15.35\,\mathrm{eV}$,
because the occupation scales as $m^{-4}$. This is neither an observational
mass bound nor a condensate formation prediction. It shows that the chosen
MeV equilibrium halo gas cannot be assigned a Bose-condensate phonon
without an additional physical state or a different history. Non-equilibrium
condensates, a separate light component and a noncondensed response medium
need their own equations rather than this ideal-gas assumption.

\section{The next discriminating routes}
The viable work now branches at a concrete physical distinction.
First, close the particle abundance as a particle problem: match the
source-containing cuts and finite thermal endpoint profiles, then evolve
the actual Higgs-production initial state with reciprocal energy flow.
This can establish whether the stable MeV carrier is viable cold matter;
it cannot by itself establish the galaxy law.

Second, derive a response state from an action and compute its retarded
baryon-to-force susceptibility. Its static limit must give the galaxy
law, its temporal poles must show stable relaxation, and its metric stress
must give lensing. Equations~\eqref{ag:higgsresponse} and the occupation
test make a direct Higgs contact and an assumed MeV condensate poor
identifications for that state. A distinct collective sector is an open
route, with additional parameters counted explicitly. No novel or unique
mechanism has yet been demonstrated.

Third, make any vacuum link pass the pressure-aware exchange identity.
For constant $G,c,\kappa$, conserved baryons and $a_0=\kappa c
\sqrt{G\rho_{\rm DE}}$, use energy-equivalent densities and
$p_i=w_i\rho_i c^2$:
\begin{align}
 \dot\rho_c+3H(1+w_c)\rho_c&=Q,\\
 \dot\rho_{\rm DE}+3H(1+w_{\rm DE})\rho_{\rm DE}&=-Q,\\
 2{\dd\ln a_0\over\dd\ln a}+3(1+w_{\rm DE})
 +{\rho_c\over\rho_{\rm DE}}
 \left[{\dd\ln(\rho_c/\rho_b)\over\dd\ln a}+3w_c\right]&=0.
\end{align}
Here $\rho_c$ includes the carrier's total rest and kinetic energy.
The reactions above conserve that total internally and have $Q=0$;
their changing weighted count is not conversion into vacuum energy.
For a constant vacuum, $w_{\rm DE}=-1$, $\rho_{\rm DE}$ and $a_0$
are constant. The additive vacuum term and fitted $\kappa$ remain
unexplained in the scalar model. A single action that fixes them, produces
the carrier, selects the response phase and satisfies metric/cosmological
tests remains the actual closure target.

\section{Evidence and limitations}
The deterministic scripts and input-hashed manifests are in
\texttt{source\_assisted\_rate\_table\_run\_002},
\texttt{assisted\_abundance\_run\_001}, and
\texttt{portal\_response\_run\_001}. There are respectively 47, 618
and 9 bounded controls. Self-review verifies the reaction-count convention,
contact-shift sign, density units and conditional scope. It does not supply
independent referee certification or global novelty. Historical report
sections retain their previous scope; this checkpoint supersedes their
request to run the pole-free source-assisted evolution.

\begin{thebibliography}{9}
\bibitem{ag:hoferichter} M. Hoferichter, P. Klos, J. Men\'endez and
A. Schwenk, Improved limits for Higgs-portal dark matter from LHC searches,
Phys. Rev. Lett. 119 (2017) 181803,
\href{https://arxiv.org/abs/1708.02245v2}{arXiv:1708.02245v2}.
The abstract authenticates the quoted $f_N$ benchmark; the local-medium
reduction above is derived here with its stated approximations.
\end{thebibliography}

\clearpage

\part{Finite-time collective response and soft particle matching}

\section{The question and existing obstructions}
The preceding checkpoint found a tiny direct Higgs-induced response at
chosen galactic densities. Could a stable medium nevertheless amplify it
through near-critical susceptibility? This report supplies a finite-time
test of that proposal, and calculates a previously omitted soft-Higgs loop
contribution to the daughter coupling. Neither supplies a completed galaxy
mechanism or a microscopic origin for the vacuum energy.

Earlier project work already derives the protected carrier's dilute coherent
phase pressure proportional to the square of its chemical-potential variable,
and shows why simply assigning it a MOND phonon is unjustified. A separate
barotropic-pressure argument already excludes using a single density-only
equation of state to support different exactly flat speeds over overlapping
density intervals. These results are retained, not rediscovered here.
The superfluid dark-matter route is also established literature:
Ref.~\cite{pr:superfluid} explicitly assumes its phonon action and a baryon
coupling; it does not derive those ingredients from this scalar carrier.

\section{A finite-time bound for a stable quadratic sector}
Consider physical coordinates $q$ with constant real positive-definite
kinetic matrix $M$ and positive-semidefinite stiffness $K$:
\begin{equation}
 L={1\over2}\dot q^TM\dot q-{1\over2}q^TKq
                +J(t)b^Tq,\qquad q(0)=\dot q(0)=0.
 \label{pr:system}
\end{equation}
Let $U$ orthogonally diagonalize $M^{-1/2}KM^{-1/2}$, with eigenvalues
$\Omega_j^2\geq0$. For a unit constant source the reciprocal observable
$b^Tq$ has step response
\begin{equation}
 R_b(t)=\sum_j f_j{1-\cos(\Omega_jt)\over\Omega_j^2},\qquad
 f_j=\bigl(U_j^TM^{-1/2}b\bigr)^2\geq0,\quad
 \sum_j f_j=b^TM^{-1}b.
 \label{pr:step}
\end{equation}
The zero-frequency summand is defined by its limit $f_jt^2/2$.
Since $0\leq1-\cos x\leq x^2/2$,
\begin{equation}
 0\leq R_b(t)\leq {t^2\over2}b^TM^{-1}b.
 \label{pr:bound}
\end{equation}
This proves the bound for arbitrary finite dimension, rather than inferring
it from numerical sampling. The static susceptibility, when it exists, is
$\chi_b(0)=b^TK^{-1}b=\sum_j f_j/\Omega_j^2$. It can diverge as a
mode softens while the finite-time bound remains unchanged.
Reaching a fraction $f>0$ of that static response requires
\begin{equation}
 t\geq\sqrt{{2f\chi_b(0)\over b^TM^{-1}b}}.
 \label{pr:minimumtime}
\end{equation}
This is necessary, not sufficient, and does not establish settling.
Undamped modes may overshoot their static response and need never settle.

The result also has a useful cross-observable version. For an observable
$a^Tq$ and any source satisfying $|J(s)|\leq J_{\max}$ on $0\leq s\leq t$,
the retarded mode kernel is $\sin(\Omega_j(t-s))/\Omega_j$.
Its absolute value is at most $t-s$. Cauchy--Schwarz on the mode
coefficients therefore gives
\begin{equation}
 |a^Tq(t)|\leq {J_{\max}t^2\over2}
       \sqrt{(a^TM^{-1}a)(b^TM^{-1}b)}.
 \label{pr:cross}
\end{equation}
This proof does not require reciprocal positive spectral weights for the
cross response. A regulated positive oscillator bath can be included among
the coordinates before elimination. A continuum extension requires finite
inertial sums and control of its limit; an unregulated effective kernel
does not automatically meet these hypotheses.

The theorem assumes zero perturbation initial data, a fixed stable
quadratic expansion, canonical second-order dynamics and a bounded source.
Preloaded excitations, a changing background, nonlinear phase transitions,
active gain or unstable modes are not covered. Their energy source and
evolution must instead be calculated. Neither a universal halo exclusion
nor a new general theorem of mathematical physics is claimed.

\section{Application to the existing scalar--polarization construction}
For a single real Fourier component, use the previously proposed static
energy and bare kinetic matrices
\begin{equation}
 K_k=\begin{pmatrix}Ck^2&-k\\-k&C^{-1}+\Pi_k\end{pmatrix},\qquad
 M=\begin{pmatrix}Z_\varphi&0\\0&I_P\end{pmatrix},\qquad
 C,Z_\varphi,I_P>0,\quad k>0,\quad\Pi_k>0.
\end{equation}
The determinant is $Ck^2\Pi_k>0$. A source on $\varphi$ has
\begin{equation}
 \chi_{\varphi\varphi}(0,k)
 ={C^{-1}+\Pi_k\over Ck^2\Pi_k}
 ={1\over Ck^2}+{1\over C^2k^2\Pi_k},\qquad
 R_{\varphi\varphi}(t,k)\leq {t^2\over2Z_\varphi}.
 \label{pr:twofield}
\end{equation}
Reducing $\Pi_k$ increases the static response but does not remove the
necessary time in Eq.~\eqref{pr:minimumtime}. For the dimensionless example
$C=Z_\varphi=I_P=k=1$, $\Pi_k=10^{-6}$, the static susceptibility is
$1000001$, while the exact step response at $t=10$ is only $25.25115$.
Reaching half the static value cannot occur before $t=1000.0005$ by the
bound. These are dimensionless controls, not galaxy times in years.

This two-coordinate numerical example freezes $\Pi_k$ as an instantaneous
stiffness. The actual nodal-fermion construction has a frequency-dependent
retarded kernel, positive absorption and induced inertia. Its earlier
positive-bath representation allows the theorem to be applied to a
regulated full quadratic realization when the stated positivity and finite
inertial-sum assumptions hold. The frozen two-coordinate curve is not
promoted to that fermionic medium's actual formation history.

The prior isotropic gapless branch has $\Pi_k\simeq\eta k$ statically
and $\Omega_{\rm slow}^2\simeq C^2\eta k^3/Z_\varphi$ at small $k$.
The static susceptibility scales as $k^{-3}$ in that regime; its response
time grows correspondingly as $k^{-3/2}$ at least. Actual damping and
settling are further requirements. None of the unfixed dimensional kinetic
coefficients has been calibrated by the dimensionless calculation here.

There is a dimensional necessary condition if this same bare scalar sector
is required to have propagation speed $C/Z_\varphi\leq c^2$. Define its
static enhancement over its own uncoupled scalar response by
$\mathcal A=Ck^2\chi_{\varphi\varphi}(0,k)
=1+(C\Pi_k)^{-1}$. Equations~\eqref{pr:minimumtime} and
\eqref{pr:twofield}, with $L=1/k$, imply
\begin{equation}
 t\geq {L\over c}\sqrt{2f\mathcal A},\qquad
 \mathcal A\leq {1\over2f}\left({ct\over L}\right)^2
 \quad\hbox{if fraction $f$ is reached by time $t$}.
 \label{pr:causalbound}
\end{equation}
For half-response in 10 Gyr at $L=1$ kpc, the necessary enhancement ceiling
is about $9.4\,10^{12}$. A putative $10^{47}$ enhancement in this same
channel would take at least $1.03\,10^{27}$ years to reach half-response
under these assumptions. This does not identify the earlier contact shift
with $\mathcal A$; matching their source and observable is still required.
It closes the specific shortcut of freely increasing static susceptibility
while retaining arbitrarily fast, initially unperturbed subluminal response.
It does not exclude a nonlinear or preloaded state, or a different source
normalization. The additional bare-speed assumption is not proved for the
current preferred-frame band model or for an unspecified ultraviolet theory.

The deterministic run checks 48 two-field parameter combinations and six
three-coordinate cross responses. Generalized-eigenmode responses agree
with independently integrated ODEs in the latter cases. The resulting 648
controls test implementations and examples; the displayed inequalities are
proved by the spectral argument above.

\section{A normal-gas comparison with an explicit time normalization}
As a separate controlled surrogate, take a homogeneous mean-subtracted
Newtonian fluid, fixed density $\rho_0$, adiabatic sound speed $c_s$, and
no expansion. Linear continuity, Euler and Poisson equations give
\begin{equation}
 \ddot{\delta\rho_c}+(c_s^2k^2-\omega_J^2)\delta\rho_c
       =\omega_J^2\delta\rho_b,
 \qquad \omega_J^2=4\pi G\rho_0.
\end{equation}
For $k\geq k_J=\omega_J/c_s$, zero initial perturbation and a constant
baryon-density Fourier source, define $a=k/k_J$ and $\tau=\omega_Jt$.
The density ratio is
\begin{equation}
 {\delta\rho_c(t)\over\delta\rho_b}
 ={1-\cos(\sqrt{a^2-1}\,\tau)\over a^2-1}
 \leq {\tau^2\over2}.
 \label{pr:gas}
\end{equation}
At $a=1$ the limiting equality holds; there is no finite static solution.
At $a>1$ the static ratio is $(a^2-1)^{-1}$. Below $k_J$ the system is
Jeans unstable and this stable-mode argument is not applicable.

For the selected comparison $\rho_0=0.3\,\mathrm{GeV\,cm}^{-3}/c^2$
and one-dimensional dispersion $150\,\mathrm{km\,s}^{-1}$, let
$c_s^2=(5/3)\sigma_v^2$. Then
\begin{equation}
 \omega_J^{-1}=47.3144\,\mathrm{Myr},\qquad
 k_J^{-1}=9.37049\,\mathrm{kpc},\qquad
 {1\over2}(\omega_J\,10\,\mathrm{Gyr})^2=22334.9.
\end{equation}
Here $k_J^{-1}$ is an inverse wavenumber, not the wavelength $2\pi/k_J$.
For example $a=1.000001$ has static ratio about $499999.75$ but a
10-Gyr step ratio about $22169.10$. Critical static response does not
yield arbitrarily rapid amplification.

This is a linear fluid comparison, not the earlier collisionless Maxwellian
Volterra calculation, nor a physical finite halo equilibrium. Large density
contrasts would invalidate linearization; the inequalities can be tested
with arbitrarily small source amplitudes. These density-response units
must not be equated to the earlier dimensionless contact shift: connecting
the contact to the actual source vector $b$ and force observable $a$ is a
missing matching calculation. The value 22334.9 is therefore not a
universal bound on the theory's galactic amplification.

\section{The soft-Higgs daughter coupling is not exactly zero}
The selected scalar action has $g(X^*\Phi^2+X\Phi^{*2})$ and
$\lambda_{HX}|H|^2|X|^2$ in its potential. Even with a zero daughter
portal at tree matching, a cubic loop with an internal daughter and source
produces a finite soft-Higgs vertex. Let the external Higgs momentum vanish
and the external daughter be on shell, $p^2=m^2$.

The normalization and sign can be checked directly in the Euclidean
functional integral. Integrating out the quadratic complex $X$ field gives
the nonlocal action term $-g^2\Phi^{*2}G_X\Phi^2$. Contracting one
daughter pair gives a quadratic correction $-4g^2G_XG_\Phi$; the factor
four comes from the two choices at each end. Its mass-squared correction,
up to an $M$-independent subtraction at fixed other parameters, is
\begin{equation}
 \Sigma_{\rm mass}(p^2)={g^2\over4\pi^2}
 \int_0^1\dd x\,\ln{D(x,p^2)\over\mu_R^2}+\text{constant},\quad
 D=xM^2+(1-x)m^2-x(1-x)p^2.
\end{equation}
The divergent subtraction in the cubic mass diagram does not affect its
finite derivative with respect to $M^2$. Inserting the soft Higgs through
$\delta M^2=\lambda_{HX}vh$ therefore gives the leading coefficient
of $-vh|\Phi|^2$ in the Lagrangian:
\begin{equation}
 C_{H\Phi}^{\rm soft,loop}
 =\lambda_{HX}{g^2\over4\pi^2}
 \int_0^1{ x\,\dd x\over xM^2+(1-x)^2m^2}>0.
 \label{pr:loop}
\end{equation}
External-leg corrections enter beyond this leading order when the tree
daughter portal is zero. A symmetry-allowed finite daughter-portal matching
term can add to this result; the loop contribution alone is not a
renormalization-condition-independent prediction for the full coupling.
The statement is that the specified tree zero is not an exact vanishing
of this perturbative soft vertex.
This vertex respects the dark global $U(1)$. That symmetry statement alone
does not extend a classical coldness bound to the full renormalized action.

As an analytic normalization control, at $M=2m$ the integral is
\begin{equation}
 {1\over m^2}\int_0^1{x\over(1+x)^2}\dd x
 ={\ln2-1/2\over m^2}.
\end{equation}
Independent quadratures and a mass-derivative check are performed at six
mass ratios. The selected 2-MeV candidate gives approximately
\begin{equation}
 {C_{H\Phi}^{\rm soft,loop}\over\lambda_{HX}}
 \simeq6.50\,10^{-8},\qquad
 C_{H\Phi}^{\rm soft,loop}\simeq1.02\,10^{-17}.
\end{equation}
With the same slow nonrelativistic nucleon medium as the previous
checkpoint, $vh=-f_N\rho_b/m_h^2$. Thus this loop contribution gives
\begin{equation}
 \delta m_\Phi^2=-C_{H\Phi}^{\rm soft,loop}{f_N\rho_b\over m_h^2}.
\end{equation}
At the chosen $0.1\,M_\odot\,\mathrm{pc}^{-3}$ density its magnitude is
of order $6\,10^{-63}\,\mathrm{GeV}^2$, a fractional daughter
mass-squared shift of order $1.5\,10^{-57}$. This perturbative soft
effect does not repair the galaxy-response mechanism. Its smallness is
not a complete proof about collective effects or a different finite
matching coupling. Equation~\eqref{pr:loop} also cannot be used as the
on-shell $125$-GeV Higgs decay form factor: the soft and hard kinematics
are different. No new freeze-in abundance is inferred from it.

\section{What changed and what remains}
The static critical construction now has a necessary finite-time response
condition independent of how small its stiffness is made. The particle
bridge now includes an explicit leading radiative soft coupling, with its
normalization, sign and matching caveat. These are constraints on the
route to closure, not closure itself.

A successful continuation must calculate the physical source vector and
kinetic normalization of one occupied response state, apply the time bound
with those units, and solve the nonlinear formation and energy receiver.
The nodal medium remains a separately postulated occupied sector; the
protected scalar carrier is not automatically that sector. Their bridge,
source selection, finite mass supply, actual metric/lensing and cosmological
tests remain required. The additive vacuum magnitude and the normalization
of $a_0=\kappa c\sqrt{G\rho_{\rm DE}}$ remain unfixed. Internal scalar
number-changing reactions still have zero vacuum energy-transfer $Q$.

Evidence is in \texttt{passive\_response\_run\_001} (648 bounded controls)
and \texttt{soft\_higgs\_matching\_run\_001} (19 controls), with
\texttt{causal\_response\_run\_001} (three unit controls).
The mathematical and physical review is self-review. Neither a globally
novel mechanism nor independent referee
certification is claimed.

\begin{thebibliography}{9}
\bibitem{pr:superfluid} L. Berezhiani and J. Khoury,
Theory of Dark Matter Superfluidity,
\href{https://arxiv.org/html/1507.01019v2}{arXiv:1507.01019v2},
Sec.~1.3 and Sec.~4. Cited only as established context for a distinct
phonon construction; no current observational verdict is imported.
\end{thebibliography}

\clearpage

\part{A reciprocal carrier--medium density candidate}

\section{Candidate and scope}
This is a separately specified, nonrelativistic effective density-energy
candidate linking a cold carrier density to the existing nodal-fermion
response medium. It is not a change to the frozen theory, a derived
relativistic completion, or a matching calculation from the MeV scalar
action. Its new density coefficients are inputs. The result is constructive
at the local thermodynamic level: a reciprocal coupling with a sufficient
boundedness/convexity condition and a solved finite-charge benchmark.
The simplest strong-coupling branch also exposes a new obstruction to
preserving the desired galaxy law.

Density-mediated attraction and collapse in Bose--Fermi mixtures are
established subjects; Ref.~\cite{rb:mixture} is cited for that context,
not as a theorem about this candidate. The calculations below are derived
for the project's nodal energy and chosen contact terms. No global novelty
claim is made.

\section{Keep the occupation-dependent energy}
Let $\Delta=h|\mathbf P|\geq0$ be the nodal gap, and let $q$ denote
signed fermion number density relative to the specified filled negative
bands. This $q$ is not an electric charge density or the cold particle
count. In the zero-temperature low-energy patch, with the same quadratic
subtraction as earlier reports, the grand-potential density is
\begin{equation}
 \Omega_F(\Delta,\mu)=
 \begin{cases}
 A\Delta^3/3,& |\mu|\leq\Delta,\\
 A|\mu|\Delta^2/2-A|\mu|^3/6,& |\mu|>\Delta,
 \end{cases}\qquad A>0.
 \label{rb:grand}
\end{equation}
The $-A|\mu|^3/6$ term cannot be discarded when $\mu$ depends on
the carrier: it contributes to reciprocal density feedback. Differentiating,
\begin{equation}
 q=-\partial_\mu\Omega_F
 ={A\over2}\operatorname{sgn}\mu\,(\mu^2-\Delta^2)
                  \Theta(|\mu|-\Delta).
\end{equation}
Legendre transformation gives the particularly simple canonical energy
\begin{equation}
 F_F(\Delta,q)={A\over3}
                 \left(\Delta^2+{2|q|\over A}\right)^{3/2}.
 \label{rb:canonical}
\end{equation}
For $q\ne0$ define $E=\sqrt{\Delta^2+2|q|/A}$; then
\begin{equation}
 \partial_qF_F=\operatorname{sgn}q\,E,\qquad
 \partial_\Delta F_F=A\Delta E.
\end{equation}
At $q=0$ the charge subgradient is the interval $[-\Delta,\Delta]$.
The canonical formulation thus retains a neutral gapped region rather than
forcing a spurious smooth compressibility there.

These expressions require $\Delta,E,T\ll\Lambda$ and use the already
chosen critical subtraction. At finite cutoff there are analytic
corrections. They do not determine the filled band's absolute gravitational
energy or vacuum constant. A reference filled sea is additional matter
content; counting its signed excitations does not count its complete
cosmological inventory.

\section{An explicit reciprocal local interaction}
Let $n\geq0$ be the cold number density in this low-energy density model.
Choose
\begin{equation}
 \mathcal E_{\rm dens}=m_cn+{K_c\over2}n^2
       +\zeta nq+{u\over2}q^2+F_F(\Delta,q),
 \quad K_c>0,\quad u\geq{\zeta^2\over K_c}.
 \label{rb:energy}
\end{equation}
In natural units $[n]=[q]=E^3$, $[A]=E$, and
$[K_c]=[\zeta]=[u]=E^{-2}$. An operator-level motivation is a regulated
preferred-frame interaction
$\zeta\,\hat n_c:\!\Psi^\dagger\Psi\!:
+(u/2)(:\!\Psi^\dagger\Psi\!:)^2$, with carrier contact repulsion.
It preserves both number symmetries in this specified NR sector. The
density functional is a chosen Hartree-level description, not a rigorous
many-body bound for those regulated operators. Matching its $n,K_c,\zeta,u$
to the original $X,\Phi$ fields, including their neutral production stock
and number-changing channels, remains undone. In particular no condensate
or dilute coherent phase is assumed established for the MeV halo.

The sufficient stability condition follows from the exact identity
\begin{equation}
 {K_c\over2}n^2+\zeta nq+{u\over2}q^2
 ={K_c\over2}\left(n+{\zeta q\over K_c}\right)^2
  +{1\over2}\left(u-{\zeta^2\over K_c}\right)q^2.
 \label{rb:square}
\end{equation}
Also $F_F$ is jointly convex in $(\boldsymbol\Delta,q)$:
$|\boldsymbol\Delta|^2+2|q|/A$ is nonnegative convex and its $3/2$
power is increasing convex. On the smooth $q>0$ radial sector,
\begin{equation}
 F_{\Delta\Delta}=A(E+\Delta^2/E),\quad
 F_{qq}={1\over AE},\quad F_{\Delta q}={\Delta\over E},
 \qquad \det\operatorname{Hess}_{\Delta,q}F_F=1.
\end{equation}
The nonsmooth charge origin is covered by convexity/subgradients.
Consequently Eq.~\eqref{rb:energy} is nonnegative above the stated rest
term and has no negative local density-Hessian direction under the
sufficient condition. This is a local thermodynamic result for the
chosen patch energy, not full quantum, gradient, gravitational or dynamical
stability. With $u=0$ the coupled density Hessian can lose positivity;
a bare density coupling cannot simply inherit the old coldness proof.

Reciprocity is explicit. The chemical-potential shifts are
$\delta\mu_c=\zeta q$ and $\delta\mu_F=\zeta n+u q$.
For number currents $\mathbf j_c,\mathbf j_F$ with closed boundaries,
\begin{equation}
 {\dd\over\dd t}\int\zeta nq\,\dd^3x
 =\int\zeta(\mathbf j_c\cdot\nabla q+
                       \mathbf j_F\cdot\nabla n)\,\dd^3x.
\end{equation}
The carrier and fermion work terms are respectively
$-\zeta\mathbf j_c\cdot\nabla q$ and
$-\zeta\mathbf j_F\cdot\nabla n$, cancelling this change.
The $u q^2/2$ term has the corresponding internal pressure/work balance.
Thus the bridge transfers energy and momentum through an interaction;
it does not supply free cooling or transfer energy to a vacuum constant.
An actual occupation/transport evolution is still required.

\section{Reciprocal occupation and local force}
Use the existing positive-square field energy, once only,
\begin{equation}
 \mathcal E_{\rm field}={C\over2}
       |\nabla\varphi-\mathbf P/C|^2
       +F_F(h|\mathbf P|,q),\qquad C>0.
\end{equation}
Here $F_F$ is the same term as in Eq.~\eqref{rb:energy}; it is not added
twice and no extra bare cubic is included. A prescribed baryon source
has energy $q_b\rho_b\varphi$. Neglect gradients of $\mathbf P$ only
for this local-density static benchmark. In the spherical aligned branch,
\begin{equation}
 D=C\varphi'-P={q_bM_b(<r)\over4\pi r^2},\qquad
 {D\over C}=A h^2P E,
 \qquad \Delta E=S(r):={D\over CAh}.
 \label{rb:field}
\end{equation}
At fixed cold density, the fermion chemical equilibrium is
\begin{equation}
 \mu_F-\zeta n=\operatorname{sgn}q\,E+u q.
\end{equation}
When $|\mu_F-\zeta n|\leq\sqrt S$, the neutral solution has
$q=0$, $\Delta=E=\sqrt S$. Otherwise, with $s$ the sign of that
chemical potential, there is a unique $E>\sqrt S$ satisfying
\begin{equation}
 |\mu_F-\zeta n|=E+{uA\over2}(E^2-S^2/E^2),\quad
 \Delta={S\over E},\quad q={sA\over2}(E^2-\Delta^2).
\end{equation}
The right-hand side is strictly increasing. Charged regions satisfy
$\Delta\leq\sqrt S$: at fixed baryon flux the bridge suppresses this
local field response relative to the original neutral branch.

One first benchmark prescribed $n=1+2/(1+r^2)$ and solved
$\int q\,\dd^3x=0$ on $0.1\leq r\leq R$. Setting the background
chemical potential by hand to give only negative interior $q$ would violate
that inventory condition. The solution instead has inner negative charge,
a neutral cubic interval, and positive compensating outer charge.
Nine cases were checked by two Gauss grids and independently by adaptive
integration split at the occupation edges. This benchmark does not derive
the cold density. The next calculation removes that prescription.

\section{Vary both densities at fixed inventories}
Take a specified external baryonic potential energy per carrier $U(r)$.
Vary $n$, $q$ and the local scalar response with fixed cold number $N_c$
and fixed total fermion charge zero. On a positive-$n$ branch,
\begin{align}
 K_cn+\zeta q+U&=\nu,\qquad \nu=\mu_c-m_c,\\
 \operatorname{sgn}q\,E+u q+\zeta n&=\mu_F.
\end{align}
Eliminating $n$ gives
\begin{equation}
 \operatorname{sgn}q\,E+u_{\rm eff}q
       =\widetilde\mu+{\zeta U\over K_c},\quad
 u_{\rm eff}=u-{\zeta^2\over K_c}\geq0,\quad
 \widetilde\mu=\mu_F-{\zeta\nu\over K_c}.
 \label{rb:eliminate}
\end{equation}
This feedback is absent if the cold density is frozen. If the solution
crosses $n=0$, a constrained active-set problem replaces these formulas;
the executed examples explicitly check $n>0$.

Choose dimensionless $A=C=h=K_c=1$, $u=\zeta^2+1$,
\begin{equation}
 U(r)=-{2\over\sqrt{1+r^2}},\qquad
 D(r)={r\over(1+r^2)^{3/2}},\qquad 0.1\leq r\leq R.
\end{equation}
These are a prescribed external Plummer trap and scalar baryon flux, not
a self-consistent gravitational halo. The cold inventory is fixed at the
uncoupled value $N_c=\int(1-U)\,\dd^3x$. Because $\int q=0$,
the integrated first equation gives $\nu=1$ exactly and
\begin{equation}
 n(r)=1-U(r)-\zeta q(r),\qquad
 \int[n-(1-U)]\,\dd^3x=-\zeta\int q\,\dd^3x=0.
\end{equation}
The local cold profile is now an output of the reciprocal equilibrium.
For $\zeta=1,u=2$, adaptive results are:
\begin{center}
\begin{tabular}{rrrrr}
\toprule
$R$ & $\widetilde\mu$ & $Q_F^+=|Q_F^-|$ & neutral interval
 & maximum $\delta n/n_{\rm uncoupled}$\\
\midrule
10 & .344375 & 6.03679 & $[2.8165,8.6441]$ & .1913\\
30 & .122857 & 41.4106 & $[8.1088,24.3947]$ & .2307\\
100 & .0384464 & 208.539 & $[26.0006,78.0232]$ & .2461\\
\bottomrule
\end{tabular}
\end{center}
The full runs include $\zeta=0.3,1,3$. Both inventories and all three
stationarity equations are checked. Every sample has positive cold density.
There is also a simple uniform check for the selected roots:
$0<\widetilde\mu<1$, $0<\zeta\leq3$ and $u_{\rm eff}=1$.
In positive-charge regions $q+\sqrt{2q}\leq\widetilde\mu<1$, so
$q<(\sqrt3-1)^2/2$. Negative-charge regions increase $n$.
Consequently $n>1-3(\sqrt3-1)^2/2>0.196$ throughout the annulus,
not just at quadrature nodes. The numerical roots satisfy the assumed
chemical-potential interval with a finite margin.
Furthermore,
the sampled excitation energies remain below one percent of the chosen
cutoff $\Lambda=1000$. The cutoff test does not bound omitted nonlocal
gradient terms or finite-temperature effects. The container and each
inventory are finite for every example. Reflecting boundaries are external
constraints here; a physical wall/interface stress and energy are not
included. The cold inventory $N_c$ changes across the
different-$R$ family. None is a formed galaxy or a physical dimensional fit.

\section{A proved boundary-sensitive limit}
There is a discriminating asymptotic for this benchmark. Write the large-$r$
inputs as $U\sim-V_0/r$ and $S\sim D_0^2/r^2$, and define
$B=\zeta V_0/K_c>D_0$. Consider the family of larger annuli with the
same background density and fixed zero fermion charge. Then
\begin{equation}
 \widetilde\mu R\longrightarrow2B,\qquad
 {r_{\rm neutral,in}\over R}\longrightarrow{1\over2}-{D_0\over2B},
 \qquad
 {r_{\rm neutral,out}\over R}\longrightarrow{1\over2}+{D_0\over2B}.
 \label{rb:boundary}
\end{equation}
To prove this, set $r=Rx$, $\widetilde\mu=a/R$ and $w=ax-B$.
At fixed $x>0$, the $u_{\rm eff}q$ term is subleading. In charged regions,
\begin{equation}
 E\sim{|w|\over Rx},\quad
 \Delta\sim{D_0^2\over Rx|w|},\quad
 q\sim{A\,\operatorname{sgn}w\over2R^2x^2}
          (w^2-D_0^4/w^2),\qquad |w|>D_0.
\end{equation}
The neutral interval is $|w|\leq D_0$. After multiplying by the spherical
volume element, $Q_F/R$ tends to
\begin{equation}
 {2\pi A\over a}\{J(a-B)-J(B)\},\qquad
 J(t)={t^3\over3}+{D_0^4\over t}-{4D_0^3\over3},\quad t>D_0,
\end{equation}
in a neighborhood of its root $a=2B$.
The derivative $J'(t)=t^2-D_0^4/t^2$ is strictly positive.
Thus neutrality uniquely selects $a=2B$, giving Eq.~\eqref{rb:boundary}.
The limit is controlled: $E\leq|\widetilde\mu+\zeta U/K_c|$ in
charged regions and $|U(r)|\leq V_0/r$ give
$R^2x^2|q(Rx,a/R)|\leq A(|a|x+B)^2/2$ on the annulus.
Extend the integrand by zero below its inner endpoint. This is an
integrable bound on $0\leq x\leq1$, uniform in $R$ for fixed $a$.
Dominated convergence and brackets on either side of
the strictly monotone charge root give the stated convergence, rather
than presuming a limit of the roots.

For $D_0=1$, $\zeta=1$ has $B=2$. At $R=10000$ the calculation gives
$\widetilde\mu R=3.99499$ versus limit $4$, and neutral edge fractions
$(.250314,.750941)$ versus $(.25,.75)$. For $\zeta=3$, $B=6$,
the corresponding values are $11.9660$ versus $12$, and
$(.417849,.584989)$ versus $(5/12,7/12)$.
The $R=1000$ comparisons approach those limits less closely. These
finite checks support the implementation; the asymptotic follows from the
displayed scaling and charge integral.

This is an increasingly large background-bath family, not a way to give a
single isolated galaxy infinite cold stock. It proves that the neutral
cubic interval can track the boundary rather than converge to a fixed
interior region. For the ordered limit $R\to\infty$ followed by a large
fixed radius with $B>D_0$, the interior charged branch has
$E\sim B/r$ and $\Delta\sim D_0^2/(Br)$.
If both the prescribed gravitational trap strength $V_0$ and scalar source
amplitude $D_0^2$ are proportional to $M_b$ at fixed couplings, this
coefficient $D_0^2/B$ becomes independent of $M_b$. It does not give the
required $\sqrt{M_b}$ scaling. This is a conditional obstruction to the
strong minimal density bridge as a universal mass--velocity selector,
not a theorem excluding weaker couplings, different states, self-gravity,
nonlocal response or other bridge operators.

\section{Mediator matching: a constructive option and a charge gate}
An ordinary massive vector provides a useful microscopic test of the
coefficient matrix. Let it couple to conserved currents as
$V_\mu(g_cJ_c^\mu+g_FJ_F^\mu)$, with a healthy positive mass $m_V$.
The static Gauss operator is $-\nabla^2+m_V^2$ and its inverse is
positive. At wavelengths much larger than $m_V^{-1}$, the density energy is
\begin{equation}
 \delta\mathcal E_V={1\over2m_V^2}(g_cn_Q+g_Fq)^2.
\end{equation}
Here $n_Q=J_c^0$ is the signed cold charge density. With bare repulsions
$K_0,u_0\geq0$, the matched coefficients are
\begin{align}
 K_c&=K_0+g_c^2/m_V^2,&
 u&=u_0+g_F^2/m_V^2,&
 \zeta&=g_cg_F/m_V^2,\\
 K_cu-\zeta^2&=K_0u_0+
              (K_0g_F^2+u_0g_c^2)/m_V^2\geq0.
\end{align}
The chosen $K_c=1,u=\zeta^2+1$ matrix is realizable in this contact
sector with $u_0=0$, $K_0=1/(\zeta^2+1)$,
$g_c^2/m_V^2=\zeta^2/(\zeta^2+1)$ and
$g_F^2/m_V^2=\zeta^2+1$. Thus the positive coefficient relation need
not be an unexplained cancellation of unrelated negative terms.
This statement does not supply a causal relativistic completion of the
preferred-frame nodal bands or fix physical mediator masses/couplings.

The charge identification is decisive. Gauge invariance of
$X^*\Phi^2$ requires $Q_X=2Q_\Phi$. Its leading NR current is
\begin{equation}
 n_Q=n_{\Phi+}-n_{\Phi-}+2(n_{X+}-n_{X-}),
\end{equation}
whereas the positive cold stock includes the sum of particle and antiparticle
densities. The Higgs process $h\to XX^*\to2\Phi+2\Phi^*$ creates
zero total cold charge. On its equal-sign-density branch $n_Q=0$ even
with substantial positive rest mass. One therefore cannot replace $n_Q$
by the positive $n$ in the preceding benchmark and call the vector a
completed bridge to that neutral stock. A charged subcomponent/asymmetry
or a separately solved charge-polarization dynamics is required.
Neutrality does not forbid charge fluctuations or particle scattering,
but it removes this
identification of the mean Hartree source. No charged production history
has been supplied here.

A charge-even alternative is a real scalar mediator with an energy term
\begin{equation}
 {m_\sigma^2\over2}\sigma^2+\sigma(a_cn+a_Fq).
\end{equation}
Eliminating its static homogeneous value gives
$\delta\mathcal E_\sigma=-(a_cn+a_Fq)^2/(2m_\sigma^2)$.
The attractive sign differs from the vector's sign. With bare repulsions,
\begin{align}
 K_c&=K_0-a_c^2/m_\sigma^2,&
 u&=u_0-a_F^2/m_\sigma^2,&
 \zeta&=-a_ca_F/m_\sigma^2,\\
 K_cu-\zeta^2&=K_0u_0-
              (K_0a_F^2+u_0a_c^2)/m_\sigma^2.
\end{align}
Positive final diagonal coefficients and nonnegative determinant are
needed; scalar exchange alone does not generate the stabilization.
A coupling $\sigma|\Phi|^2$ is charge-even and respects the original
dark $U(1)$. Coupling the same scalar to signed medium density uses
$\sigma u_\mu J_F^\mu$, where $u^\mu$ is the medium rest-frame vector.
This is an allowed preferred-frame EFT possibility, with additional bare
contacts and a medium-state assumption, not an established covariant
particle theory. Its masses, repulsions, matching, finite-temperature
response and energy receiver remain to be calculated. Both mediator
options preserve the distinction between neutral cold mass and signed
fermion/cold charge rather than hiding it in a density symbol.

\section{What this supplies and the next obligation}
The previously unspecified bridge is now an explicit effective density
interaction with reciprocal chemical forces, a sufficient convexity
condition, a counted signed fermion inventory, and a jointly responding
cold population in a specified external potential. These are stronger
results than keeping a prescribed cold density or a fixed chemical-potential
reservoir. The added $u$ term and all bridge coefficients remain new inputs;
they are not explained by the original particle model or dark energy.

The next decisive calculation is to match a charge-even bridge to the
neutral particle action, or supply and count an actual charged subcomponent,
then determine its state/temperature and transport,
then solve the nonlocal occupied-band problem with live gravitational
sources and a physical boundary. A vacuum response term, an energy receiver
and absolute metric stress must be counted once. The present energy transfer
is internal to cold/fermion/field sectors and has no vacuum $Q$ by itself.
The source-assisted Higgs-production hierarchy remains independently open.
No galaxy fit, lensing prediction, origin of $\kappa$ or vacuum magnitude,
or whole-theory closure follows from these examples.

\section{Evidence}
The deterministic scripts and manifests are
\texttt{reciprocal\_density\_bridge\_run\_001} (151 controls),
\texttt{bridge\_charge\_resolution\_run\_001} (88),
\texttt{joint\_cold\_bridge\_run\_001} (97), and
\texttt{bridge\_large\_radius\_run\_002} (46), and
\texttt{bridge\_mediator\_matching\_run\_001} (16 symbolic identities).
Counts include inherited symbolic and local checks and are not mutually
independent physics validations. The first large-radius run completed the
mathematics but failed to collect the nested namespace's check list for
output; the corrected second run changes that bookkeeping only. No gate
was weakened. Proof and scope review are self-review, not an independent
referee or novelty certificate.

\begin{thebibliography}{9}
\bibitem{rb:mixture} K. Manabe and Y. Ohashi,
Thermodynamic stability, compressibility matrix, and effects of mediated
interactions in a strongly-interacting Bose--Fermi mixture,
Phys. Rev. A 103 (2021) 063317,
\href{https://arxiv.org/abs/2103.16843v1}{arXiv:2103.16843v1}.
The primary abstract authenticates the established density-collapse and
repulsive-stabilization context. No quantitative result for the present
nodal model is imported from it.
\end{thebibliography}

\clearpage

\part{Charge-even particle matching and the density vertex}

\section{The new result and its boundary}
The preceding density bridge is now connected explicitly to the original
complex carrier field by a real scalar mediator. The tree potential has
an exact positive-square representation with specified stabilizing terms.
This removes the earlier vector-current identification problem for neutral
particle--antiparticle stock at leading slow-density order. It does not
derive the mediator, its coefficients, the occupied fermion medium, a
relativistic band completion, or the galaxy law.

A separate calculation gives the density-vertex response of the neutral
gapped nodal patch. It differs from the mass/polarization-vertex kernel
computed earlier. It has zero uniform dynamical response, a computed
finite-momentum absorption, and a conservative static stability bound for
the chosen mediator contacts. Physical production, thermal history and
formation remain required. All claims here are self-reviewed and scoped;
global novelty is unestablished.

\section{A particle-level charge-even potential}
Use $\hbar=c=1$, $I=|\Phi|^2$ and the original potential
\begin{equation}
 V_c=m^2I+4m^2|X|^2+
       \left|\sqrt{\Delta_c}X+{g\over\sqrt{\Delta_c}}\Phi^2\right|^2
       +\lambda_{\rm rep}I^2,
 \qquad \Delta_c>0,\quad\lambda_{\rm rep}>0.
\end{equation}
Let $q_F$ be the signed fermion density relative to the chosen filled
bands, and let $F_F(\mathbf P,q_F)\geq0$ be the previously specified
canonical local medium energy. Introduce a real scalar $\sigma$ with
canonical positive kinetic term, mass $M_\sigma>0$, and
\begin{align}
 \delta V={}&{M_\sigma^2\over2}\sigma^2
       +a\sigma I+b\sigma q_F
       +\delta\lambda I^2+{u_0\over2}q_F^2,\\
 \delta\lambda&={a^2\over2M_\sigma^2},\qquad
 \zeta_I=-{ab\over M_\sigma^2},\qquad
 u_0={b^2\over M_\sigma^2}+u_R,\qquad
 u_R\geq{\zeta_I^2\over2\lambda_{\rm rep}}.
 \label{ns:parameters}
\end{align}
Here $[a]=E$, $[b]=1$, $[\zeta_I]=E^{-1}$ and $[u_0]=[u_R]=E^{-2}$.
The added quartic and repulsive fermion contact are chosen matching
conditions, not symmetry-protected relations or predictions of the old
action. A preferred-frame realization of the density vertex is
$b\sigma u_\mu J_F^\mu$. The medium rest frame $u^\mu$ is still an
assumption requiring a covariant physical origin. A Lorentz-scalar
$\sigma\bar\Psi\Psi$ would be a different, mass-type vertex.

There is an exact algebraic identity:
\begin{align}
 V_c+\delta V+F_F={}&m^2I+4m^2|X|^2+
 \left|\sqrt{\Delta_c}X+{g\over\sqrt{\Delta_c}}\Phi^2\right|^2\\
 &+{M_\sigma^2\over2}
       \left[\sigma+{aI+bq_F\over M_\sigma^2}\right]^2\nonumber\\
 &+\lambda_{\rm rep}
       \left[I+{\zeta_Iq_F\over2\lambda_{\rm rep}}\right]^2
 +{1\over2}\left[u_R-{\zeta_I^2\over2\lambda_{\rm rep}}\right]q_F^2
 +F_F.\nonumber
 \label{ns:positive}
\end{align}
Thus the stated density-functional extension preserves the pointwise lower
bound $V\geq m^2I+4m^2|X|^2$. The already completed Higgs-sector square
can be retained without changing this statement. The added scalar may have
nontrivial local mass shifts, but the whole potential, not an isolated
effective mass, determines this lower bound.

The dark transformation $(\Phi,X)\mapsto(e^{i\alpha}\Phi,e^{2i\alpha}X)$
still leaves all new terms invariant. With the original current convention
$j_\Phi^0=i(\Phi^*\dot\Phi-\Phi\dot\Phi^*)$, the kinetic identity is
\begin{equation}
 |\dot\Phi+im\Phi|^2
       =|\dot\Phi|^2+m^2I-mj_\Phi^0.
\end{equation}
The source has the analogous $\dot X+2imX$ identity. Therefore an extension
whose remaining Hamiltonian density is the specified nonnegative medium
energy and positive kinetic/gradient terms obeys $E\geq m|Q_c|$ relative
to its chosen vacuum. This is a conditional classical/density-functional
result. It is not a bound for an unspecified regulated fermion operator,
arbitrary finite-temperature state, gravity or the full quantum ultraviolet
theory. As before, $Q_c=0$ supplies no nonzero neutral-stock abundance or
coldness protection. No such physical closure is inferred from the square.

The exact compensating quartic in Eq.~\eqref{ns:parameters} is a convenient
scattering prescription, not a necessary equality for boundedness. More
generally set $\lambda_*=\lambda_{\rm rep}+\delta\lambda
-a^2/(2M_\sigma^2)$. The same argument works on the open parameter region
\begin{equation}
 \lambda_*>0,\qquad
 u_0-{b^2\over M_\sigma^2}>{\zeta_I^2\over2\lambda_*}.
\end{equation}
The algebraic energy condition therefore is not restricted to an exact
cancellation surface. Radiative control and actual physical matching of
that open region are still unproved.

The fermion repulsion can also have an explicit mediator origin: a healthy
massive vector coupled only to $J_F^\mu$ gives $u_0=g_F^2/M_W^2$ at
frequencies and momenta well below $M_W$. It does not require assigning a
vector charge to neutral cold mass, since the cold coupling is through
$\sigma I$. For the example below, $M_W=0.2$ GeV would give
$g_F\simeq0.1012$. This is an optional EFT realization of the chosen
repulsion, with another physical degree of freedom and mass mechanism to
count. It is not a completion of the preferred-frame band Hamiltonian or
an observationally validated dark sector.

\section{What matches to the neutral slow density}
For slowly changing, incoherent NR particle and antiparticle populations,
\begin{equation}
 \langle I\rangle_{\rm fast}\simeq {n_++n_-\over2m},
 \qquad J_c^0\simeq n_+-n_-+2(n_{X+}-n_{X-}).
\end{equation}
Unlike a vector-current source, $a\sigma I$ need not vanish for neutral
stock. Solving the soft static scalar equation gives
\begin{equation}
 \sigma\simeq-{aI+bq_F\over M_\sigma^2},\qquad
 V_{\rm eff}=V_c+\zeta_I Iq_F+{u_R\over2}q_F^2+F_F.
\end{equation}
Thus the slow charge-even bridge coefficient is
\begin{equation}
 \zeta_{\rm NR}={\zeta_I\over2m},\qquad
 K_{\rm NR,bound}={\lambda_{\rm rep}\over2m^2},\qquad
 u_R\geq{\zeta_{\rm NR}^2\over K_{\rm NR,bound}}.
\end{equation}
The $K$ expression translates the quartic lower bound using
$\langle I^2\rangle\geq\langle I\rangle^2$; it is not the actual
neutral-mixture equation of state. Occupation correlations, the opposite-charge
channel and the source population must be included to derive that EOS.
Fast coherent $2m$ terms and a mediator which cannot follow them also
prevent an unrestricted instantaneous density replacement. The fundamental
potential uses $I$ and retains those distinctions.

\section{The compensation is not an unchanged scattering theory}
The added quartic has four-point amplitude $-4\delta\lambda$. The scalar
has vertex $-ia$ and contributes two exchange trees in each elastic channel:
\begin{align}
 \delta\mathcal M_{++}&=-4\delta\lambda
       -{a^2\over t-M_\sigma^2}-{a^2\over u-M_\sigma^2},\\
 \delta\mathcal M_{+-}&=-4\delta\lambda
       -{a^2\over t-M_\sigma^2}-{a^2\over s-M_\sigma^2}.
\end{align}
For identical charges at threshold $t=u=0$, compensation is exact:
$\delta\mathcal M_{++}=0$. For a particle--antiparticle pair at threshold,
$s=4m^2,t=0$, however,
\begin{equation}
 \delta\mathcal M_{+-}={4m^2a^2\over
             M_\sigma^2(M_\sigma^2-4m^2)}.
 \label{ns:opposite}
\end{equation}
This finite annihilation-frequency remainder is precisely why soft density
matching cannot be used as all-channel vacuum matching. The formula excludes
the mediator pole; the executed family has $M_\sigma>2m$. A real scalar
also adds the physical vacuum process
\begin{equation}
 \sigma\to\Phi\Phi^*,\qquad
 \Gamma_\sigma={a^2\over16\pi M_\sigma}
                 \sqrt{1-{4m^2\over M_\sigma^2}}.
\end{equation}
The inverse process and additional scatterings must be counted if the
mediator is populated. This is an enlarged particle theory, so the old
freeze-in history and collision network cannot simply be inherited.

For the separately selected $m=2$ MeV carrier choose
$M_\sigma=20$ MeV, $a/M_\sigma=.02$, $b=-.01$ and
$u_R=\zeta_I^2/\lambda_{\rm rep}$, twice the sufficient minimum.
The tree matching gives
\begin{align}
 \delta\lambda&=2.0\,10^{-4},&
 \zeta_I&=.01\,\mathrm{GeV}^{-1},&
 \zeta_{\rm NR}&=2.5\,\mathrm{GeV}^{-2},\\
 u_R&=.00601764\,\mathrm{GeV}^{-2},&
 u_0&=.256018\,\mathrm{GeV}^{-2},&
 K_{\rm NR,bound}&=2077.23\,\mathrm{GeV}^{-2}.
\end{align}
The old same-charge threshold cross section is retained at this tree level;
the opposite-charge threshold cross section changes by about
$-3.3433\,10^{-4}$ fractionally. The vacuum scalar width is
$1.55939\,10^{-7}$ GeV and its rest lifetime is $4.22095\,10^{-18}$ s.
Neither is a medium equilibration time. Four mediator masses and three
CM speeds are tested, with exchange channels independently enumerated by
their allowed charge partitions. There is no observational fit or preferred
physical mediator mass inferred from this chosen family.

\section{Compute the density vertex, not the old mass vertex}
Take a homogeneous neutral gapped nodal patch at zero temperature.
For axial energy momentum $\mathcal Q=v k_z$ and positive imaginary
frequency $\Omega$, let $s=\sqrt{\mathcal Q^2+\Omega^2}$.
The response convention is $\delta q_F=\chi_F\delta\mu_F$ for a
Hamiltonian source $-\delta\mu_F q_F$. With nodal DOS coefficient $A$,
\begin{align}
 {\chi_F(i\Omega,\mathcal Q)\over A}
 &=\mathcal Q^2\int_0^1
       {x(1-x)\,\dd x\over\sqrt{\Delta^2+x(1-x)s^2}}\\
 &=\mathcal Q^2\left[{\Delta\over2s^2}
       +{s^2-4\Delta^2\over4s^3}
                       \arctan{s\over2\Delta}\right].
 \label{ns:density}
\end{align}
The continuous limits are understood. In particular
\begin{equation}
 \chi_F(i\Omega,0)=0,\qquad
 \chi_F(0,\mathcal Q)\simeq {A\mathcal Q^2\over6\Delta}
 \quad(\mathcal Q\ll\Delta),\qquad
 \chi_F\big|_{\Delta=0}={A\pi\mathcal Q^2\over8s}.
\end{equation}
The uniform zero follows because the total fermion number commutes with
the unperturbed band Hamiltonian. It is not the nonzero uniform mass-vertex
response from earlier calculations. An oscillating uniform $\sigma q_F$
source cannot automatically be assigned that earlier interband damping.

An independent band-projector evaluation checks Eq.~\eqref{ns:density}.
With $\mathbf p_\pm=\mathbf p\pm\boldsymbol{\mathcal Q}/2$,
$E_\pm=\sqrt{p_\pm^2+\Delta^2}$ and $S_E=E_++E_-$,
\begin{equation}
 {\chi_F\over A}={1\over2\pi}\int\dd^2p\,
 \left[1-{\mathbf p_+\cdot\mathbf p_-+\Delta^2\over E_+E_-}\right]
 {S_E\over S_E^2+\Omega^2}.
\end{equation}
The overlap, not a mass-vertex trace, fixes the normalization. Radial and
angular refinements agree with both the closed form and Feynman-parameter
integral in six selected cases. This continuum axial-patch result requires
all external energies well below the band cutoff; it is not an isotropic,
finite-band, doped or finite-temperature medium calculation.

The functional form agrees with the established gapped Dirac density
response in Ref.~\cite{ns:pyatkovskiy}, Eqs.~(A.13) and (A.16), after
removing its electromagnetic prefactors. In its two-dimensional convention
$A_{2D}=N_f/(\pi v_F^2)$ and
$\chi_F=\epsilon_0\Pi_0/(2\pi e^2)$, with $\hbar=1$ and
$\mathcal Q=v_Fk$. Our nodal coefficient $A$ has its previously derived
three-dimensional patch measure instead. This comparison authenticates
the source-dependent normalization; it does not identify dark matter with
graphene or establish novelty.

For $e^{-i\omega t}$ continue
$s=\sqrt{\mathcal Q^2-(\omega+i0)^2}$. For $\omega>0$ and
$t=\sqrt{\omega^2-\mathcal Q^2}$ the absorption is
\begin{equation}
 {\operatorname{Im}\chi_F\over A}
 ={\pi\mathcal Q^2(t^2+4\Delta^2)\over8t^3}
                 \Theta(t-2\Delta)\geq0.
\end{equation}
The sign is positive for response to a chemical source. The induced
mediator stiffness contains $-b^2\chi_F$ and hence has the passive
negative imaginary part. Finite-momentum fermion excitations can receive
energy; their actual production and heat transport are still uncomputed.

For a weak external mode $\sigma=\operatorname{Re}[\sigma_k
e^{-i\omega t+i\mathbf k\cdot\mathbf x}]$ the leading mean absorbed power
density in the repulsively dressed density channel is
\begin{equation}
 \overline{\mathcal P}_{\sigma\to F}
 ={b^2\omega\over2}{\operatorname{Im}\chi_F\over
                    |1+u_0\chi_F|^2}|\sigma_k|^2\geq0.
\end{equation}
For an approximately free narrow scalar mode of energy density
$\omega^2|\sigma_k|^2/2$, its corresponding leading energy-loss rate is
$b^2\operatorname{Im}\chi_F/(\omega|1+u_0\chi_F|^2)$.
Large dispersive self-energies require the actual mode residue instead.
Both formulas vanish at $\mathcal Q=0$ in this neutral patch. They identify
an energy receiver without supplying a temperature, equilibration time,
nonlinear occupation history or vacuum energy transfer.

\section{A conservative static mediator bound}
Retain the bare repulsion $u_0q_F^2/2$ used in the particle potential.
In the neutral density channel its random-phase response is
\begin{equation}
 \chi_{F,u_0}={\chi_F\over1+u_0\chi_F},\qquad
 0\leq\chi_{F,u_0}(0,k)\leq u_0^{-1}.
\end{equation}
For the reduced homogeneous density model with bare $I$-curvature at least
$2(\lambda_{\rm rep}+\delta\lambda)$, the soft cold susceptibility obeys
$\chi_I\leq[2(\lambda_{\rm rep}+\delta\lambda)]^{-1}$.
Additional positive cold gradients, source-square curvature or pressure
only reduce this susceptibility within that comparison. The actual neutral
particle gas has not been proved to obey that density-curvature model.
Self-gravity and negative-compressibility backgrounds are excluded.
The scalar inverse static kernel then obeys
\begin{align}
 \mathcal K_\sigma(0,k)
 &=M_\sigma^2+k^2-a^2\chi_I(0,k)
                 -b^2\chi_{F,u_0}(0,k)\\
 &\geq M_\sigma^2+k^2
       -{a^2\over2(\lambda_{\rm rep}+\delta\lambda)}-{b^2\over u_0}.
 \label{ns:kernelbound}
\end{align}
The zero-$k$ lower margin is positive when the contact inequality is strict.
Indeed the algebraic $(I,q_F,\sigma)$ Hessian with the nonnegative medium
curvature omitted has determinant
\begin{equation}
 M_\sigma^2(2\lambda_{\rm rep}u_R-\zeta_I^2)>0.
\end{equation}
Its scalar Schur complement is exactly the lower margin in
Eq.~\eqref{ns:kernelbound}. For the selected family it exceeds
$0.0116126\,M_\sigma^2$. This establishes a conservative static density-channel
test under the stated susceptibility assumptions, not all coupled
inhomogeneous modes or global dynamical stability. Doped density/mass mixing,
other spatial directions, finite temperature, counterterms and full cold
momentum response are not certified by this bound.

\section{Energy exchange and the vacuum distinction}
The canonical mediator equation and its separately defined stress obey
\begin{align}
 (\Box+M_\sigma^2)\sigma&=-aI-bq_F,\\
 T_\sigma^{\mu\nu}&=\partial^\mu\sigma\partial^\nu\sigma
  -\eta^{\mu\nu}\left[{1\over2}(\partial\sigma)^2
                         -{M_\sigma^2\over2}\sigma^2\right],\\
 \partial_\mu T_\sigma^{\mu\nu}
     &=-(aI+bq_F)\partial^\nu\sigma.
\end{align}
Thus its energy loss is reciprocal work on the other sectors when their
equations and interaction energy are included. Assigning the interaction
to one component changes the split, not total energy conservation.
The absolute metric stress of the preferred-frame medium and its rest
vector remains an additional requirement.

The canonical mediator density and isotropically averaged pressure are
\begin{equation}
 \rho_\sigma={\dot\sigma^2+|\nabla\sigma|^2+M_\sigma^2\sigma^2\over2},
 \quad p_\sigma={\dot\sigma^2\over2}
                -{|\nabla\sigma|^2\over6}
                -{M_\sigma^2\sigma^2\over2}.
\end{equation}
An isolated rapidly oscillating homogeneous massive mode has average
$p_\sigma\simeq0$, whereas a homogeneous static potential term alone has
$p_\sigma=-\rho_\sigma$. Neither limiting expression makes the
density-sourced configuration into an independently conserved vacuum fluid.
Its source, interaction stress, gradients and decay must be counted.
At $I=q_F=0$ the selected tree mediator vacuum is $\sigma=0$.
An additive constant $V_0$ remains free and supplies
$T_{V_0}^{\mu\nu}=V_0\eta^{\mu\nu}$ with pressure $-V_0$.
Nothing here fixes $V_0$ or derives the acceleration normalization from it.
Particle/medium work is not a homogeneous cold-to-vacuum transfer by itself.

\section{Remaining requirements}
This continuation supplies an explicit charge-even mediator tied to
$|\Phi|^2$, an exact stabilized potential, its tree scattering consequences,
and the correct linear neutral-medium vertex. The mediator exchanges energy
internally. Adding a vacuum constant still leaves these local equations
unchanged, and no vacuum $Q$, vacuum magnitude or $\kappa$ is derived.

The next decisive obligations are radiative matching of the chosen contact
relations, the neutral mixture's actual correlations/EOS, mediator population
and the enlarged production network, and finite-momentum transport with the
actual medium state. These must precede applying the old external-trap
density profiles to this particle candidate. A physical occupied medium,
finite stock and boundaries, full metric/lensing, native cosmology and joint
empirical validation remain necessary. The theory remains open.

\section{Evidence}
The input-hashed executions are
\texttt{neutral\_scalar\_bridge\_run\_001} (55 bounded controls) and
\texttt{scalar\_density\_response\_run\_001} (140).
The analytic square and susceptibility arguments specify the claim scope;
finite controls are not proofs of the missing physical completion.

\begin{thebibliography}{9}
\bibitem{ns:pyatkovskiy} P. K. Pyatkovskiy,
Dynamical polarization, screening, and plasmons in gapped graphene,
\href{https://arxiv.org/pdf/0808.0931v2}{arXiv:0808.0931v2},
J. Phys.: Condens. Matter 21 (2009) 025506, Eqs.~(A.13), (A.16).
\end{thebibliography}

\clearpage

\part{Radiative thresholds and vacuum matching}

\section{Result and conventions}
The new charge-even mediator has a calculable one-loop tadpole and carrier
mass correction. The small fractional carrier correction can exceed the
original narrow source detuning. A separately specified weaker-coupling
comparison keeps that decay open without finite pole-mass readjustment in
the retained calculation. Neither comparison is a complete renormalized
particle/medium theory or a new observational exclusion.

The vacuum calculation also makes a distinction essential to the full
cold/dark-energy story: fixing a field's vacuum expectation value and mass
does not fix the additive gravitational vacuum energy. A known sequestering
gravity option is examined as a separate route, with its hypotheses and a
nonzero feedback from a proposed acceleration coefficient. No complete new
gravity action, predicted vacuum magnitude or novel mechanism is claimed.

The loop sector uses the same scalar normalization and parameters as the
previous report. The reference renormalization scale is $\mu_R=m$.
The symbol $\Delta_c=M^2-4m^2$ denotes source detuning; $\Delta_f$
continues to mean a nodal gap where needed. They are different quantities.

\section{The scalar one-loop potential and required matching}
On the homogeneous $\Phi=X=0$ axis, a constant mediator background gives
the two real daughter fluctuations the common squared mass
$z(\sigma)=m^2+a\sigma$. For $z>0$, their
$\overline{\rm MS}$ contribution is
\begin{equation}
 V_1^\Phi(\sigma)={z(\sigma)^2\over32\pi^2}
       \left[\ln{z(\sigma)\over\mu_R^2}-{3\over2}\right].
 \label{rv:CW}
\end{equation}
The normalization follows from the standard real-scalar formula in
Ref.~\cite{rv:martin}, Eqs.~(1.1) and (3.2), with two real degrees of
freedom. This is the scalar-only contribution on this axis, not the full
Higgs, medium, gauge or two-loop effective potential. A region with
$z\leq0$ is outside this real-axis expansion and cannot be declared stable
by extrapolating its Taylor series.

The derivatives at the origin are
\begin{align}
 T_\sigma&={a m^2\over16\pi^2}
                         \left[\ln{m^2\over\mu_R^2}-1\right],\\
 \delta M_{\sigma,0}^2&={a^2\over16\pi^2}\ln{m^2\over\mu_R^2},\\
 V_1'''(0)&={a^3\over16\pi^2m^2},\qquad
 V_1''''(0)=-{a^4\over16\pi^2m^4}.
\end{align}
The last two are finite zero-momentum vertices in this calculation, not
evidence that independent cubic/quartic UV counterterms are required by
this diagram. The negative fourth Taylor coefficient alone does not prove
an instability of the exact multiscalar potential.

For an independent UV-operator check write
$\Phi=(f_1+if_2)/\sqrt2$, $X=(x_1+ix_2)/\sqrt2$ and form the real
five-field potential Hessian $\mathcal H$. With
$\lambda_\Phi=\lambda_{\rm tot}+\delta\lambda$, its
$\sigma$-dependent squared trace is exactly
\begin{equation}
 \operatorname{Tr}\mathcal H^2(\sigma)
 -\operatorname{Tr}\mathcal H^2(0)
 =2a^2\sigma^2+4am^2\sigma+16a\lambda_\Phi\sigma|\Phi|^2.
\end{equation}
Thus a tadpole counterterm is needed even though the tree construction
chose $\sigma=0$ at empty vacuum. The same normalization gives the retained
one-loop contributions
\begin{equation}
 \beta_{t_\sigma}={am^2\over8\pi^2},\quad
 \beta_{M_\sigma^2}={a^2\over8\pi^2},\quad
 \beta_a={a\lambda_\Phi\over2\pi^2},\quad
 \left.\beta_{m^2}\right|_a={a^2\over8\pi^2}.
 \label{rv:betas}
\end{equation}
These are selected scalar contributions, not the complete beta-function
system. At this order the explicit scale derivative of
Eq.~\eqref{rv:CW} is cancelled on the axis by
\begin{equation}
 \beta_{V_0}^\Phi={m^4\over16\pi^2},\qquad
 \beta_{t_\sigma}\sigma+{\beta_{M_\sigma^2}\over2}\sigma^2.
\end{equation}

At $\mu_R=m$, keeping the renormalized tree linear coefficient $t_R=0$
without an additional tadpole-matching condition shifts the minimum at leading order by
$\delta\sigma=-T_\sigma/M_\sigma^2$.
For $a=.0004$ GeV and $M_\sigma=.02$ GeV this is
$2.53303\,10^{-8}$ GeV. Its additional fractional carrier mass-squared
shift is $a\delta\sigma/m^2=2.53303\,10^{-6}$. An alternative, used in
the next pole comparison, fixes $\langle\sigma\rangle=0$ by a finite
tadpole condition. These are different matching choices; neither fixes
the vacuum constant or all other masses.

\section{The narrow source threshold is the stronger test}
With the tadpole separately fixed to zero, the new mediator's one-loop
contribution to the daughter pole mass at $p^2=m^2$ is
\begin{equation}
 \delta m_{\Phi,\rm pole}^2
 ={a^2\over16\pi^2}\int_0^1\dd x\,
       \ln{xM_\sigma^2+(1-x)^2m^2\over\mu_R^2}.
 \label{rv:pole}
\end{equation}
This is the finite $\overline{\rm MS}$-parameter comparison at the stated
scale, not a prediction obtained while simultaneously fixing the pole mass
to its old numerical value. The new $a$ interaction does not shift the
$X$ pole at this retained one-loop order. Consequently
\begin{equation}
 \delta\Delta_{c,\rm pole}=-4\delta m_{\Phi,\rm pole}^2,
 \qquad
 \left.\beta_{\Delta_c}\right|_a=-{a^2\over2\pi^2}.
\end{equation}
For the original $m=2$ MeV, $\Delta_c=6.4\,10^{-9}$ GeV$^2$ family:
\begin{center}
\begin{tabular}{rrrrr}
\toprule
$M_\sigma$ (MeV)&$a$ (MeV)&$\delta m_\Phi^2/m^2$
 &$\Delta_{c,\rm pole}$ (GeV$^2$)&decay open\\
\midrule
6 & .12 & $3.2631\,10^{-5}$ & $5.87790\,10^{-9}$ & yes\\
10 & .20 & $1.4793\,10^{-4}$ & $4.03307\,10^{-9}$ & yes\\
20 & .40 & $9.2379\,10^{-4}$ & $-8.38071\,10^{-9}$ & no\\
200 & 4.00 & $.20799$ & $-3.32148\,10^{-6}$ & no$^*$\\
20 & .10 & $5.7737\,10^{-5}$ & $5.47621\,10^{-9}$ & yes\\
\bottomrule
\end{tabular}
\end{center}
All entries retain only the new $a^2$ correction with the old numerical
masses treated as renormalized inputs at $\mu_R=m$. The starred case also
fails the declared $|\delta m^2|/m^2<0.1$ loop-size diagnostic; its
quantitative extrapolation is therefore not certified. A successful
numerical integration does not override that scope failure.

The threshold diagnostic is scale-consistent. Running the input
$\Delta_c$ with Eq.~\eqref{rv:betas} cancels the explicit logarithm in
Eq.~\eqref{rv:pole} at the retained order. Changing the scale while
holding all those renormalized numerical inputs fixed would instead define
a different physical matching choice. Alternatively one may impose the
original chosen pole masses as inputs and choose the finite counterterms
accordingly. The negative-gap entry does not exclude that on-shell model;
it demonstrates the threshold matching it needs.
The retained cancellation is not a complete RG flow or a certificate that
the tree positive-square parameterization stays in its positive cone at
every scale.

At $M_\sigma=20$ MeV the no-readjustment criterion
$4\delta m_\Phi^2<\Delta_c$ gives $a<.263210$ MeV in this comparison.
The separately specified $a=.10$ MeV continuation lies below it. Its
compensated tree coefficients are
\begin{equation}
 \delta\lambda=1.25\,10^{-5},\quad
 \zeta_{\rm NR}=.625\,\mathrm{GeV}^{-2},\quad
 u_R=.000376103\,\mathrm{GeV}^{-2},\quad
 u_0=.250376\,\mathrm{GeV}^{-2}.
\end{equation}
It retains the symbolic tree positive-square condition and passes the
stated new-loop mass/threshold comparisons. Its thermodynamics, actual
abundance and full renormalized amplitudes are not inherited from the
older stronger candidate.

The mediator's own vacuum bubble supplies an independent normalization
check:
\begin{equation}
 \delta\mathcal K_\sigma(p^2)
 ={a^2\over16\pi^2}\int_0^1\dd x\,
 \ln{m^2-x(1-x)(p^2+i0)\over\mu_R^2}.
\end{equation}
At $p^2=M_\sigma^2>4m^2$, its imaginary part obeys
$-\operatorname{Im}\delta\mathcal K_\sigma/M_\sigma
=\Gamma_{\sigma\to\Phi\Phi^*}$, reproducing the preceding tree width.
At low $p^2$ its induced kinetic coefficient is
$\delta Z_\sigma=a^2/(96\pi^2m^2)>0$.
A zero curvature correction at $\mu_R=m$ therefore does not imply a zero
on-shell mediator self-energy.
This vacuum calculation is different from absorption into the fermion
medium, whose density vertex has the uniform Ward zero found earlier.

\section{Vacuum energy is an independent matching condition}
At the chosen reference scale the free scalar fluctuation pieces are
\begin{align}
 V_1^\Phi(0)&=-{3m^4\over64\pi^2},\\
 V_1^X(0)&={M^4\over32\pi^2}
                   \left[\ln{M^2\over m^2}-{3\over2}\right],\\
 V_1^\sigma(0)&={M_\sigma^4\over64\pi^2}
                   \left[\ln{M_\sigma^2\over m^2}-{3\over2}\right].
\end{align}
For the 20-MeV mediator these are respectively
$-7.59909\,10^{-14}$, $-9.19156\,10^{-14}$ and
$7.86549\,10^{-10}$ GeV$^4$. A declared comparison reference
$\epsilon_{\rm ref}=(2.25\,\mathrm{meV})^4=2.56289\,10^{-47}$ GeV$^4$
is smaller than their sum by a factor about $3.07\,10^{37}$.
This reference is an input scale, not a new cosmological fit.

Those partial, scheme-dependent terms are not a prediction that the
Universe has that large or signed cosmological constant. The full theory
contains a renormalized vacuum parameter:
\begin{equation}
 \epsilon_{\rm vac,phys}=V_{0,R}(\mu_R)
               +V_{\rm loops}(\hbox{renormalized vacuum};\mu_R).
\end{equation}
Its scale dependence cancels when all retained terms are matched
consistently. Fixing the tadpole, fixing pole masses and fixing this value
are independent conditions. The massive mediator and the nodal medium
have not provided a symmetry or constraint fixing the last one.
Other matter, counterterms, the occupied-state stress and gravity are
absent from this partial scalar sum. No cancellation by the filled bands
can be assumed from a chosen zero of their homogeneous energy.

\section{A distinct known gravity route and its compatibility gate}
Vacuum-energy sequestering is a known possible gravity modification,
not a mechanism invented here. In the matter-loop, semiclassical-gravity
setting of Ref.~\cite{rv:sequester}, nonlinear four-form couplings enforce
constant gravitational auxiliaries. Its Eq.~(10), in that source's gravity
conventions, is
\begin{equation}
 M_G^2G^\mu{}_{\nu}
 =\mathcal P[T]^\mu{}_{\nu}-\Delta\Lambda\,\delta^\mu{}_{\nu},
 \qquad
 \mathcal P[T]^\mu{}_{\nu}
 =T^\mu{}_{\nu}-{1\over4}\langle T^\alpha{}_{\alpha}\rangle
                                      \delta^\mu{}_{\nu}.
\end{equation}
The normalized historic spacetime average makes $\mathcal P$ annihilate
a constant multiple of the identity. Residual flux/history data are still
needed, and graviton loops and ill-defined infinite-volume averages are
outside the imported statement. The source requires a covariant matter
action; our preferred-frame band/frame sector has not supplied one.
This remains a conditional route, not an adopted completion of the model.

Three compatibility conclusions can already be checked. First, cancelling
constant vacuum stress does not cancel the local $\sigma$ tadpole or
threshold derivatives. Second, the positive physical residual dark-energy
coefficient $B_{\rm eff}$ need not equal the auxiliary bare $\Lambda$:
it also contains the nonvacuum historic trace contribution after the
vacuum split. Third, a target $a_0^2=\kappa_a^2G B_{\rm eff}$ in natural
units cannot be replaced by a function of the shifting bare counterterm
without deriving the consequences. The source's $M_G$ is not the fitted
acceleration coefficient $\kappa_a$.

For an explicit check, suppose instead that a positive auxiliary vacuum
parameter $B$ is put directly into a toy deep-force energy,
\begin{equation}
 a_0(B)=\kappa_a\sqrt{GB},\qquad
 W(g;B)={g^3\over12\pi G a_0(B)}.
\end{equation}
This is a chosen local deep-force functional, not the solved full nodal
action. Its unavoidable derivative is
\begin{equation}
 \partial_B W=-{W\over2B},\qquad
 {W\over B}={\kappa_a^2\over12\pi}(g/a_0)^3.
\end{equation}
The usual four-form variation of a volume term $-B$ now receives
$-\partial_BW$ as well: its volume integrand is proportional to
$1+\partial_BW$, not simply one. Thus the original constraint and its
vacuum-cancellation argument cannot be carried over unchanged after this
force coupling. For $\kappa_a=1/2$ the correction magnitude is
$(g/a_0)^3/(96\pi)$; it is small in the deep regime but not zero.
This does not disprove a completed residual-flux coupling. It identifies
the extra equation such a construction must solve. A bare-counterterm
identification, or an unchanged sequestering equation with a new
$B$-dependent matter action inserted afterwards, is insufficient.

\section{What is established and what remains}
The scalar bridge has concrete radiative matching obligations and a
weaker-coupling continuation which survives the retained threshold test.
The reference scalar vacuum calculation exposes the independent vacuum
condition rather than predicting it. A known gravity route can address
constant matter-loop stress under additional hypotheses; the acceleration
coupling changes those hypotheses unless its auxiliary variation is counted.
This is progress toward a joined action, not whole-theory closure.

Next are the complete matching of tadpoles, masses, interactions and the
open positive region, the neutral mixture EOS and populated-mediator
kinetics, and an explicit covariant frame/residual-vacuum action if the
gravity route is pursued. Finite supply, physical boundaries, formation,
source selection, metric/lensing, native cosmology and joint empirical
gates remain open. No unique or globally novel mechanism is established.

Evidence: \texttt{scalar\_radiative\_audit\_run\_003} has 33 controls,
including the retained RG gap check and a weaker comparison;
\texttt{vacuum\_bridge\_compatibility\_run\_001} has seven algebraic
controls. The first radiative run failed its original quadrature gate for
$M_\sigma/m=100$. Refinement from 128 to 512 Gauss points fixes the
endpoint resolution without relaxing the $10^{-10}$ tolerance. Run 002
completed the calculation; run 003 adds the scale-consistency checks.
Physical scope failures remain in the output. Review is self-review.

\begin{thebibliography}{9}
\bibitem{rv:martin} S. P. Martin, Two-loop effective potential for a general
renormalizable theory and softly broken supersymmetry,
\href{https://arxiv.org/pdf/hep-ph/0111209v2}{hep-ph/0111209v2},
Phys. Rev. D 65 (2002) 116003, Eqs.~(1.1), (3.2).
\bibitem{rv:sequester} N. Kaloper, A. Padilla, D. Stefanyszyn and
G. Zahariade, A Manifestly Local Theory of Vacuum Energy Sequestering,
\href{https://arxiv.org/pdf/1505.01492v2}{arXiv:1505.01492v2},
Phys. Rev. Lett. 116 (2016) 051302, Eqs.~(7)--(10).
\end{thebibliography}

\clearpage

\part{A Planck-scaled flux vacuum--force candidate}

\section{What this construction tests}
The preceding checkpoint showed why inserting a force coefficient directly
as a function of a sequestered vacuum counterterm changes the constraint.
Here the force instead depends on a distinct fixed conjugate flux $Q$.
A specified Planck-dependent flux energy is varied together with the
gravity auxiliaries. This produces an explicit conditional classical link:
one flux can control both a retained vacuum term and the local cubic-force
coefficient. The ordinary unscaled flux energy is a control case which
fails to survive the projected stress equation.

This is a new candidate action family for this programme, not an amendment
to the frozen theory, a derived microscopic medium or whole-theory closure.
The functions, couplings, boundary ensemble and quantum assumptions are
stated below. The absolute flux value and fitted normalization are not
predicted. Four-forms, variable gravitational couplings and self-tuning
are established research areas \cite{pf:sequester,pf:qtheory,pf:planck};
the bounded source comparison does not establish global novelty.

\section{One action and its independent variables}
Use gravity signature $(-+++)$ and natural units in this report. Write
$s=M_G^2>0$ for the gravitational multiplier, $s_*>0$ for a reference
squared Planck mass, $\Lambda$ for a vacuum counterterm auxiliary, and
$Q>0$ for a separate flux variable. The action family is
\begin{align}
 S={}&\int\dd^4x\sqrt{-g}\left[{sR\over2}-\Lambda
               -V_r(s,Q)+\mathcal L_m(g,\psi,Q)\right]\nonumber\\
 &+\int f(\Lambda/\mu^4)F_\Lambda
       +\int\widehat f(s/s_*)F_s+\int QF_Q+S_{\partial,Q},\\
 V_r(s,Q)&={Q^2\over2}(s/s_*)^r.
 \label{pf:action}
\end{align}
The three four-forms are field strengths of independent three-form gauge
potentials. Smooth nonlinear $f,\widehat f$ have nonzero derivatives on
the chosen branch. They implement the usual local constraint mechanism
of Ref.~\cite{pf:sequester}; the added $V_r$ and $Q$ sector are the
candidate extension derived here.

The boundary term $S_{\partial,Q}$ selects fixed conjugate $Q$ for the
third three-form. In an electric description it is the Legendre boundary
term, so the reduced potential is $+V_r$, rather than the sign obtained by
substituting an on-shell field strength into an unrestricted action.
The integrated $F_Q$ is not independently fixed as another boundary datum
in this ensemble. Membranes and flux jumps are absent.

The matter hypothesis is explicit: $\mathcal L_m$ is a generally covariant,
regulated matter/frame action with no explicit dependence on $s$ or
$\Lambda$, while dependence on $Q$ is allowed and varied. Its fields
would have to include the carrier, response medium and physical rest frame.
The current preferred-frame nodal Hamiltonian has not yet supplied that
complete action. The gravity calculation below is conditional on this
matter hypothesis. Renormalization of curvature operators and graviton
loops are not established by it.

Dimensions are $[s]=[s_*]=[Q]=E^2$ and $[V_r]=[\Lambda]=E^4$.
Neither $s$ nor the flux variable $Q$ is the fitted acceleration coefficient
$\kappa_a$. The flux is also unrelated to the carrier's Noether charge.

\section{Vary first; eliminate afterwards}
Gauge-potential variation gives
\begin{equation}
 \partial_\mu\Lambda=\partial_\mu s=\partial_\mu Q=0
\end{equation}
on a connected source-free branch. With
$T_{\mu\nu}=-(2/\sqrt{-g})\delta S_m/\delta g^{\mu\nu}$, the
on-shell local metric equation is
\begin{equation}
 sG_{\mu\nu}=T_{\mu\nu}-(\Lambda+V_r)g_{\mu\nu}.
 \label{pf:metric}
\end{equation}
Before imposing constancy, the usual
$(\nabla_\mu\nabla_\nu-g_{\mu\nu}\Box)s$ term is also present.
The other variations include
\begin{align}
 {f'\over\mu^4}F_\Lambda&=\mathrm{vol}_4,\\
 {\widehat f'\over s_*}F_s
      &=\left(\partial_s V_r-{R\over2}\right)\mathrm{vol}_4,\\
 F_Q&=\left(\partial_QV_r-\partial_Q\mathcal L_m\right)
                                                  \mathrm{vol}_4.
 \label{pf:auxiliary}
\end{align}
In particular a $Q$-dependent force or matter vacuum term contributes to
the last equation. It is not dropped or treated as an unchanged field
strength.

Let $\langle\cdot\rangle$ be a well-defined normalized spacetime average,
with $c_\Lambda=\int F_\Lambda$ and $c_s=\int F_s$. The first two
constraints imply
\begin{align}
 \langle R\rangle&=2\partial_sV_r+\mathcal R_F,\\
 \mathcal R_F&=-{2\mu^4\widehat f'\over s_*f'}
                                      {c_s\over c_\Lambda},\qquad
 \Delta_F={s\mathcal R_F\over4}.
\end{align}
Taking the trace of Eq.~\eqref{pf:metric}, averaging, and using this
constraint yields
\begin{align}
 \Lambda+V_r&={\langle T\rangle\over4}
                   +{s\partial_sV_r\over2}+\Delta_F,\\
 sG_{\mu\nu}&=T_{\mu\nu}-{\langle T\rangle\over4}g_{\mu\nu}
                -\left[{s\partial_sV_r\over2}+\Delta_F\right]g_{\mu\nu}.
 \label{pf:projected}
\end{align}
These equations are the result of the new action's variations; the old
projected equation with an extra flux stress appended afterwards would
not give the same answer.

For the selected power family, $s\partial_sV_r=rV_r$. Thus the retained
flux contribution to the last coefficient is $rV_r/2$.
An unscaled flux, $r=0$, disappears from that coefficient. If the design
requirement is to retain its entire positive fixed-flux energy, the power
is selected as $r=2$. This conditional selection is not a symmetry-based
derivation of that functional dependence.

Split $T_{\mu\nu}=-\epsilon_{\rm vac}(Q)g_{\mu\nu}
+\tau_{\mu\nu}$. Because $Q$ is constant, its constant matter vacuum
contribution cancels algebraically in Eq.~\eqref{pf:projected}, leaving
\begin{equation}
 sG_{\mu\nu}=\tau_{\mu\nu}-B_{\rm eff}g_{\mu\nu},\qquad
 B_{\rm eff}={rV_r\over2}+\Delta_F+{\langle\tau\rangle\over4}.
 \label{pf:Beff}
\end{equation}
The same $\epsilon_{\rm vac}'(Q)$ remains in the $F_Q$ constitutive
equation. At fixed conjugate $Q$ it changes that field strength, not the
chosen flux value by itself. This algebraic cancellation is not a proof
that a full global solution with unchanged auxiliary flux data exists
for every loop correction. The residual, boundary and matter hypotheses
remain part of the physical problem.

There is a useful check on what this construction explains. The invertible
field redefinition $\Lambda'=\Lambda+V_r(s,Q)$ removes $V_r$ from the
bulk vacuum term but changes the topological function to
$f([\Lambda'-V_r(s,Q)]/\mu^4)F_\Lambda$. Its $s,Q$ derivatives carry
exactly the same new constraints. Thus the selected power can equivalently
be regarded as a chosen cross-coupling of the four-form constraints; it
is not a parameter-free explanation hidden by a choice of variables.

\section{A fixed-flux force sector with the required scaling}
For a local static test, choose the matter response energy at fixed $Q$:
\begin{equation}
 \mathcal E_{\rm resp}={C\over2}
       |\nabla\varphi-\mathbf P/C|^2
       +W(Q,p)+q_b\rho_b\varphi,\qquad
 W(Q,p)={\alpha p^3\over Q},\quad p=|\mathbf P|,
 \quad C,\alpha>0.
 \label{pf:force}
\end{equation}
Here $[\varphi]=E$, $[p]=E^2$, $[q_b]=E^{-1}$ and $C,\alpha$ are
dimensionless. This inverse-flux functional is a new model input. It has
not been obtained from the original scalar particle action or a protective
symmetry. In particular it is not a consequence of the earlier affine
four-form mass portal, whose failure to select a halo law remains intact.

The action can retain this dependence while keeping $\mathcal L_m$
independent of $\Lambda,s$. A possible leading nodal matching would
require $A h(Q)^3=3\alpha/Q$, or
$h(Q)=[3\alpha/(AQ)]^{1/3}$. That specifies another coupling function;
it does not derive it. If the fermions are retained, their generated cubic
must not also be added as an independent $W$. Their nonlocal gradients,
temperature, occupancy and counterterms still have to be treated.

For $r=2$ the classical fixed-flux energy at fixed $s$ is
\begin{equation}
 \mathcal U(Q,p)=V_2(s,Q)+\alpha p^3/Q.
\end{equation}
Its conjugate electric field strength, before additional matter terms, is
\begin{equation}
 f_Q=Q(s/s_*)^2-\alpha p^3/Q^2,\qquad
 \partial_Qf_Q=(s/s_*)^2+2\alpha p^3/Q^3>0.
\end{equation}
The electric Legendre map is therefore monotone on the positive-$Q$
branch. The joint radial Hessian obeys
\begin{equation}
 \det\operatorname{Hess}_{Q,p}\mathcal U
 ={6\alpha p\over Q}(s/s_*)^2
             +{3\alpha^2p^4\over Q^4}>0\quad(p>0).
\end{equation}
The vector tangential curvature is $\partial_pW/p=3\alpha p/Q\geq0$.
These are classical auxiliary/static convexity results. They do not prove
full matter/metric dynamics, and an arbitrary $Q$-dependent loop vacuum
function could alter the electric Legendre curvature. The branch is
singular at $Q=0$; no smooth zero-vacuum continuation is claimed.

Variation of Eq.~\eqref{pf:force} gives
\begin{equation}
 \nabla\cdot\mathbf D=q_b\rho_b,\qquad
 \mathbf D=C\nabla\varphi-\mathbf P,\qquad
 {\mathbf D\over C}={3\alpha p\over Q}\mathbf P.
\end{equation}
Outside a spherical source, in an aligned local branch,
\begin{equation}
 D={q_bM_b\over4\pi r^2},\qquad
 p={B_P\over r},\qquad
 B_P=\sqrt{{q_bM_bQ\over12\pi C\alpha}},\qquad
 \varphi'={D+p\over C}.
\end{equation}
The scalar acceleration on baryons has magnitude $q_b\varphi'$.
The $p$ term gives a flat-speed coefficient on an annulus where it dominates:
\begin{equation}
 {v_{\rm flat}^4\over M_b}={q_b^3Q\over12\pi C^3\alpha}.
 \label{pf:TF}
\end{equation}
An infinite isolated halo is not established by this local solution.

\section{An executed one-scale nodal realization of the leading function}
There is a less arbitrary way to obtain the required leading dependence
than prescribing $h(Q)$ alone: postulate that all dimensionful parameters
of the nodal medium scale with $\mu_Q=\sqrt Q$,
\begin{equation}
 m_*=c_m\mu_Q,\quad k_0=c_k\mu_Q,\quad
 \Lambda_F=c_\Lambda\mu_Q,\quad h={c_h\over\mu_Q},\quad
 A=A_*\mu_Q,\quad A_*={Nc_m\over2\pi v}.
\end{equation}
Here $\Lambda_F$ is a band cutoff, not the gravity counterterm. This
common-scaling ansatz is additional physics whose origin remains open.
Renormalization and the other sectors can add independent scales or
logarithms; the finite-cutoff scaling test is not a proof of protection
under those corrections.
It yields
\begin{equation}
 {A h^3\over3}p^3={A_*c_h^3\over3}{p^3\over Q},\qquad
 \alpha={A_*c_h^3\over3}.
\end{equation}
The previously calculated finite-cutoff determinant, with the same critical
quadratic subtraction, scales more generally as
\begin{equation}
 V_{F,\rm ren}(Q,p)=Q^2\mathcal F(p/Q),
\end{equation}
not exactly as a cubic at every field amplitude. This follows directly
from its radial expression
\begin{equation}
 V_{F,\rm ren}=A\int_0^{\Lambda_F}
 {\Delta_f^4\,\dd t\over2[\sqrt{t^2+\Delta_f^2}+t]^2},\qquad
 \Delta_f=hp,
\end{equation}
by scaling all energies with $\mu_Q$. Its small-gap limit is the required
$\alpha p^3/Q$. The critical quadratic counterterm coefficient
$A\Lambda_Fh^2/2=A_*c_\Lambda c_h^2/2$ is independent of $Q$;
it is still a chosen tuning, not a protected symmetry result.

For a dimensionless illustration choose
$N=2$, $v=c_h=c_\Lambda=1/2$, $c_m=2\sqrt2$, $c_k=2$.
Then $A_*=4\sqrt2/\pi$ and
$\alpha=\sqrt2/(6\pi)\simeq.0750264$. These constants are chosen, not
derived; their normalization matches the calibrated example below.
The full circular nodal tube lies within its radial domain because
$c_\Lambda<c_k^2/(2c_m)=1/\sqrt2$.
The maximum radial group speed in this tube is
$\sqrt{c_k^2+2c_mc_\Lambda}/c_m=.923880$ and the axial speed is $1/2$.
This is a bounded-band group-speed check, not a proof of causal ultraviolet
front propagation or of coupled clock/metric stability.

The finite determinant was computed by independent physical radial
quadrature and a 60-digit closed expression for three $Q$ values and
$p/Q=.001,.01,.1$. After dividing by $Q^2$ the curves coincide.
Their ratios to the cubic approximation are respectively
$.999625$, $.996250$, $.962562$. Thus the claimed leading law is controlled
only in the small-gap regime; the exact finite-cutoff force still needs
its own solution.

It is also necessary to keep track of the filled reference sea. Before
the vacuum subtraction, its contribution is
\begin{equation}
 \epsilon_{\rm sea}(Q)=-{A\Lambda_F^3\over3}
       =-{A_*c_\Lambda^3\over3}Q^2.
\end{equation}
Any part identified as a genuine constant scalar-potential term cancels
in the projected metric equation, while its $Q$ derivative remains in
Eq.~\eqref{pf:auxiliary}. The homogeneous band energy alone does not prove
a vacuum-form stress: its pressure, reference inventory, frame and regulator
must be varied in a covariant completion. Other sea stresses remain in
$\tau_{\mu\nu}$ and its historic trace. With only this
additional term at $s=s_*$, empty electric curvature is
$1-2A_*c_\Lambda^3/3$. It is $.849947>0$ for the illustration.
Other cutoffs can fail that check, and other UV/matter terms are not
certified. The metric cancellation therefore must not be used to discard
the constitutive derivative or assume electric convexity automatically.

At the declared reference $V_2=(2.25\,\mathrm{meV})^4$, one has
$\mu_Q=2.67572$ meV. The original 2-MeV carrier is about
$7.47\,10^8$ times heavier. Coupling that sector to the medium introduces
a serious scale and state-matching obligation. Neither the common-scaling
ansatz, the required cold fermion state, nor its protection from the
carrier/mediator sector has been established.

\section{Which Newton constant and which vacuum energy?}
The $D$ term is an additional inverse-square force. For an unscreened
static baryon coupling it must be included once:
\begin{equation}
 G_E={1\over8\pi s},\qquad
 G_\varphi={q_b^2\over4\pi C},\qquad
 G_{\rm lab}=G_E+G_\varphi.
\end{equation}
A physical screened or relativistic matter coupling could change this
dictionary and would have to be derived. Our dimensionless example has
$G_{\rm lab}=C=q_b=1$, hence $G_E=1-1/(4\pi)>0$ and
$s=s_*=[8\pi G_E]^{-1}$. This is a bookkeeping example, not a Solar
System, fifth-force, cosmological or lensing fit. Inferring vacuum density
from a Hubble rate would also require this theory's gravity dictionary.

At the reference Planck sector $s=s_*$, the empty $r=2$ flux has
$V_2=Q^2/2$. Defining the annular acceleration scale by
$G_{\rm lab}a_0=v_{\rm flat}^4/M_b$, one obtains
\begin{align}
 a_0&={q_b^3Q\over12\pi C^3\alpha G_{\rm lab}},\\
 a_0&=\kappa_{\rm eff}\sqrt{G_{\rm lab}V_2},\qquad
 \kappa_{\rm eff}={q_b^3\over
                6\sqrt2\pi C^3\alpha G_{\rm lab}^{3/2}}.
 \label{pf:link}
\end{align}
In SI units the second line is
$a_0=\kappa_{\rm eff}c\sqrt{G_{\rm lab}\rho_4}$ with
$\rho_4=V_2/c^2$. The square-root dependence on the flux energy follows
from the one action at fixed measured Planck sector. The value
$\kappa_{\rm eff}=1/2$ would still be a calibration of the displayed
coupling combination, not a prediction. Keeping it fixed across other
$s$ sectors would impose further coupling relations and could change
the auxiliary variations.

Equation~\eqref{pf:Beff} identifies the remaining condition for replacing
$V_2$ by the actual dark-energy coefficient:
\begin{equation}
 B_{\rm eff}=V_2+\delta B,\quad
 \delta B=\Delta_F+\langle\tau\rangle/4,\qquad
 {a_0^2\over\kappa_{\rm eff}^2G_{\rm lab}B_{\rm eff}}
                   ={V_2\over V_2+\delta B}.
\end{equation}
Thus a flux normalization, auxiliary flux/history data and the full
nonvacuum historic trace remain to be fixed. The target is not obtained
by silently setting these terms to zero.
Indeed, in dimensionless reference-sector units with $\langle\tau\rangle=0$,
the allowed choices $(Q,\Delta_F)=(1,5/2)$ and $(2,1)$ both have
$B_{\rm eff}=3$, while their force scales differ by a factor two.
This is an explicit counterexample to claiming that the action with
unrestricted residual data predicts $a_0$ from the measured dark energy.
The additional boundary/history condition is a real obligation.

As a conditional example, zero $c_s$ would set $\Delta_F=0$.
For a regulated future de Sitter region with a conserved dust density,
$a(t)=e^{Ht}$ and $\tau=-\rho_{m0}e^{-3Ht}$ give
\begin{equation}
 \langle\tau\rangle_{[0,T]}
       =-{3H\rho_{m0}T\over e^{3HT}-1}\longrightarrow0.
\end{equation}
This only shows how a normalized average can vanish under stated late-time
and boundary assumptions. It is not a solution for the actual cosmology
or fixed four-form flux constraints. Infinite-volume limits, earlier
history and the occupied medium's stress require their own audit.

An empty local constant-curvature branch can be exhibited without that
averaging example. For $r=2$, $\tau_{\mu\nu}=0$, $F_s=0$, fixed
$Q,s>0$ and $\Lambda=-\epsilon_{\rm vac}(Q)$, take
\begin{equation}
 H^2={V_2\over3s},\qquad R={4V_2\over s},\qquad
 F_Q=[\partial_QV_2+\epsilon_{\rm vac}'(Q)]\mathrm{vol}_4.
\end{equation}
The metric equation and $R/2=\partial_sV_2$ are then both satisfied;
$F_\Lambda$ is fixed by its own equation. This proves consistency of
these local empty-background equations under the stated conditions. It
does not solve a global fixed-boundary problem, choose the fluxes, or
establish perturbative health of an occupied cosmological solution.
On this empty branch the SI consistency relation is
\begin{equation}
 {a_0\over cH_\Lambda}
   =\kappa_{\rm eff}\sqrt{{3G_{\rm lab}\over8\pi G_E}}.
\end{equation}
Here $H_\Lambda$ is the vacuum curvature scale, not the matter-dependent
$H(z)$. Constant $Q$ does not imply a time-varying $a_0$ proportional to
the full expansion rate through a matter era.

\section{Finite stock, energy and remaining physical tests}
The response has positive excess energy but an infinite-tail problem:
\begin{equation}
 E_W(R_1,R_2)={4\pi\alpha B_P^3\over Q}
                         \ln{R_2\over R_1}.
\end{equation}
It diverges if the local $1/r$ branch is taken to infinite radius.
A physical finite response region and matching phase are therefore
mandatory, and could change the force and projected lensing signal.
The flat coefficient above is an annular result, not a complete halo.
The field $\mathbf P$ is also not thereby identified with a conserved
displacement of the cold mass. The earlier finite-supply and mass-transport
obligations are not removed.

At fixed $Q,s$, the pure flux energy is constant. Changes in $W$ exchange
energy with the response/matter fields and sources through the common
action; they do not establish a homogeneous conversion of vacuum energy
into cold particles. The original pressure-aware exchange null relation
remains a separate consistency condition. No membrane-driven $Q$ history,
particle abundance or charge asymmetry has been generated here.

This candidate supplies a conditional classical gravity/force link and
retains the $Q$-dependent matter variation instead of confusing $Q$ with
the vacuum counterterm. It still needs a microscopic origin and radiative
control of $\alpha p^3/Q$, a covariant occupied-band/frame action, actual
cold production/transport, a physical interface and energy receiver,
metric/lensing and local-gravity tests, native cosmology and joint
empirical validation. The absolute flux, residual boundary data and
$\kappa$ are not derived. No whole-theory or global-novelty claim is made.

\section{Evidence and literature scope}
\texttt{planck\_flux\_link\_run\_002} records 63 bounded controls:
exact trace/auxiliary algebra, monotone electric inversions, sample static
Hessians, force roots, calibrated normalization, logarithmic energy and a
late-average example. The general identities are proved in the equations
above, not by finite sampling. Run 001 is the earlier completed calculation;
run 002 explicitly adds the Newtonian constant dictionary. This is self-review.
The finite-cutoff one-scale test is
\texttt{one\_scale\_nodal\_flux\_run\_002}, with 17 controls.
Its preceding completed run used the same mathematics; run 002 adds the
explicit mpmath version to the execution metadata.

The local sequestering ancestor was checked at its precise field equations.
Primary abstracts for Refs.~\cite{pf:qtheory,pf:planck} establish related
four-form/gravitational-coupling and Planck-variation research. Their detailed
solutions are not imported here. The searches are bounded context/overlap
checks; neither absent matches nor a different notation establishes novelty.

\begin{thebibliography}{9}
\bibitem{pf:sequester} N. Kaloper, A. Padilla, D. Stefanyszyn and
G. Zahariade, A Manifestly Local Theory of Vacuum Energy Sequestering,
\href{https://arxiv.org/pdf/1505.01492v2}{arXiv:1505.01492v2},
Eqs.~(7)--(10). The new $V_r$ and third flux must be varied separately.
\bibitem{pf:qtheory} F. R. Klinkhamer and G. E. Volovik,
Dynamic vacuum variable and equilibrium approach in cosmology,
\href{https://arxiv.org/abs/0806.2805v7}{arXiv:0806.2805v7}.
Its four-form amplitude is not identified with our conjugate $Q$.
\bibitem{pf:planck} D. Sobral-Blanco and L. Lombriser,
Exploring the self-tuning of the cosmological constant from Planck mass
variation, \href{https://arxiv.org/abs/2012.01838v3}{arXiv:2012.01838v3}.
\end{thebibliography}

\clearpage

\part{Covariant host stress, density scaling and finite-stock tails}

\section{Result and scope}
The missing frame stress is physically consequential. A band sea whose scales
are fixed by the auxiliary flux can be a vacuum potential in a chosen
covariant completion. A sea whose scales follow material compression instead
has radiation-like stress; its negative energy cannot by itself support a
healthy perfect fluid. An explicit positive host repairs the elementary
longitudinal signs, but supplies new equation-of-state data. Neither choice
identifies the host with the original MeV particles.

A density compensation can preserve the leading cubic response. It transfers
density dependence into the quadratic critical subtraction and, in the
regulated local toy, worsens the infinite-tail energy. A finite material
inventory alone therefore does not complete the halo. These are conditional
calculations on specified effective actions, not a full covariant fermion,
metric, formation or quantum completion. Units are $\hbar=c=1$, metric
signature $(-+++)$, fixed positive flux $Q$ of mass dimension two.

\section{A covariant barotropic host}
Introduce three scalar material labels $X^I$ and define
\begin{align}
 B^{IJ}&=g^{\mu\nu}\partial_\mu X^I\partial_\nu X^J,&
 b&=\sqrt{\det B}>0,\\
 J^\mu&=\frac{\epsilon^{\mu\nu\rho\sigma}\epsilon_{IJK}}
 {3!\sqrt{-g}}\partial_\nu X^I\partial_\rho X^J\partial_\sigma X^K,
 &u^\mu&=J^\mu/b,\qquad u^\mu u_\mu=-1.
\end{align}
Orientation selects a future-directed current. Antisymmetry gives
$\nabla_\mu J^\mu=0$ identically. Choose
\begin{equation}
 S_h=-\int d^4x\sqrt{-g}\,\rho(b,Q).
\end{equation}
The standard nondissipative fluid dictionary is derived in
Ref.~\cite{MFfluid}, especially its equations (11)--(13), (25)--(34).
Its charged-fluid thermodynamic convention identifies $b$ with entropy
density. Here a separately normalized barotropic material current is used;
a relation to number requires fixed-composition, isentropic microscopic
matching. It is not the signed fermion density $q_F$, the carrier charge
$Q_c$, or an established cold-particle number.

Varying the metric at fixed $Q$ gives
\begin{equation}
 T_{\mu\nu}^{h}=(\rho+p)u_\mu u_\nu+p g_{\mu\nu},\qquad
 p=b\rho_b-\rho,\qquad \rho+p=b\rho_b.
 \label{mf:stress}
\end{equation}
The pressure also follows from $-\partial_V[V\rho(N/V,Q)]_N$.
For homogeneous fixed-background perturbations, up to a boundary term,
\begin{equation}
 \mathcal L_h^{(2)}=\frac{\rho+p}{2}
 [\dot{\boldsymbol\pi}^{,2}-c_s^2(\boldsymbol\nabla\cdot
 \boldsymbol\pi)^2],\qquad c_s^2=\frac{b\rho_{bb}}{\rho_b}.
 \label{mf:phonon}
\end{equation}
A direct one-dimensional check uses
$b=b_0\sqrt{(1+\epsilon\partial_x\pi)^2-\epsilon^2\dot\pi^2}$.
Thus $\rho+p>0$ and $\rho_{bb}>0$ are sufficient for positive temporal
and longitudinal gradient terms. Transverse gradients vanish at this
perfect-fluid order. This is not a transverse quantum, full fermion,
metric or ultraviolet stability certificate. At $b=0$ the material frame
is undefined; a vanishing dilute enthalpy does not establish a uniformly
controlled vacuum phase.

\section{The filled sea has two different stress interpretations}
Let $\mu=\sqrt Q$ and use the previous example's positive constants
$A_*=1.8006326323142123$, $c_\Lambda=c_h=1/2$. Define
\begin{equation}
 y=b/Q^{3/2},\qquad d=A_*c_\Lambda^3/3
 =0.07502635967975885.
\end{equation}
If $A=A_*\mu$ and $\Lambda_F=c_\Lambda\mu$ are independent of $b$,
the sea term $\epsilon_{\rm sea}=-A\Lambda_F^3/3=-dQ^2$ is a
$b$-independent potential in this chosen completion. Equation
\eqref{mf:stress} gives $p=-\epsilon_{\rm sea}$. Its vacuum-form
metric stress may enter the previous sequestering projection, while its
$Q$ derivative must remain in the electric constraint. This observation
does not prove that an actual occupied band and its frame have only that
stress.

If all momentum scales instead compress as $b^{1/3}$ with fixed velocity,
then $A=A_*\mu y^{1/3}$ and $\Lambda_F=c_\Lambda\mu y^{1/3}$ imply
\begin{equation}
 \epsilon_{\rm sea}=-d b^{4/3},\qquad
 p_{\rm sea}=\epsilon_{\rm sea}/3,\qquad
 T^{\mu}_{\ \mu,\rm sea}=0,\qquad
 \rho_{\rm sea}+p_{\rm sea}=-\tfrac43d b^{4/3}<0.
\end{equation}
This term survives subtraction of a vacuum-form constant: its trace is
zero and its stress is nonzero. Adding a constant vacuum potential cannot
repair its kinetic sign. Adding a rest term $M(Q)b$ can make the enthalpy
positive at low density, but $\rho_{bb}=-4d b^{-2/3}/9<0$ then gives a
longitudinal instability. A positive finite $b^2$ term fails to dominate
that curvature as $b\to0$.

For example, writing $\rho/Q^2=c_0+c_1 y+c_2 y^2/2-d y^{4/3}$,
its curvature changes sign at
\begin{equation}
 y_{\rm crit}=(4d/(9c_2))^{3/2}.
\end{equation}
For $c_2=1$ this is $0.00608901487444$. Positivity above this threshold
still needs the enthalpy test. It is not a theorem excluding other hosts.

An explicit sufficient repair is to supply a positive host $y^{4/3}$
coefficient $2d$, leaving net $\lambda=d>0$ after the sea:
\begin{equation}
 \rho=Q^2 f(y),\qquad
 f=c_0+c_1 y+\lambda y^{4/3}+c_2y^2/2,
 \qquad c_1,\lambda,c_2>0.
 \label{mf:healthy}
\end{equation}
For every $y>0$, $f'>0$, $f''>0$, and
$f'-yf''=c_1+(8\lambda/9)y^{1/3}>0$, proving
$\rho+p>0$ and $0<c_s^2<1$. Its trace is
\begin{equation}
 -\rho+3p=Q^2[-4c_0-c_1 y+c_2 y^2].
\end{equation}
The extra coefficients are inputs requiring particle/state matching.
The illustrative $c_0=0,c_1=c_2=1,\lambda=d$ was checked at 61 densities
from $10^{-9}$ to $10^3$ in addition to these exact identities.

\section{Cosmological roles and spectral compensation}
In a homogeneous comoving phase, current conservation fixes
$b\propto a^{-3}$ at constant $Q$. Consequently the four terms of
\eqref{mf:healthy} redshift as $a^0,a^{-3},a^{-4},a^{-6}$ respectively,
and obey $\dot\rho+3H(\rho+p)=0$. Only the constant has vacuum stress.
No vacuum-to-material production follows from this identity. Heating,
entropy production and non-equilibrium particle creation require additional
kinetics beyond this nondissipative action.

The compensated band excess has the exact local regulated expression
\begin{align}
 V_F&=-\frac A3[(\Lambda_F^2+\Delta_f^2)^{3/2}
             -\Lambda_F^3-\Delta_f^3]
             +\frac{A\Lambda_F\Delta_f^2}{2},\\
 \Delta_f&=hp,\quad p=|\boldsymbol P|,\qquad
 V_F=\frac\beta3p^3+O(p^4),\quad \beta=Ah^3.
 \label{mf:det}
\end{align}
Its quadratic critical counterterm is $K_{\rm ct}p^2$ with
$K_{\rm ct}=A\Lambda_Fh^2/2$. Three explicitly chosen assignments are
\begin{center}
\begin{tabular}{lccccc}
\toprule
Choice & $A/\mu$ & $\Lambda_F/\mu$ & $h\mu$ & $\beta$ scaling
 & $K_{\rm ct}$ scaling\\
\midrule
Flux & $A_*$ & $c_\Lambda$ & $c_h$ & $Q^{-1}$ & $1$\\
Material & $A_*y^{1/3}$ & $c_\Lambda y^{1/3}$ & $c_hy^{-1/3}$
 & $Q^{-1}y^{-2/3}$ & $1$\\
Compensated cubic & $A_*y^{1/3}$ & $c_\Lambda y^{1/3}$ & $c_hy^{-1/9}$
 & $Q^{-1}$ & $y^{4/9}$\\
\bottomrule
\end{tabular}
\end{center}
For the previous fixed-coupling annular force dictionary, $a_0\propto
\beta^{-1}$. Thus the material assignment gives $a_0\propto a^{-2}$
in a homogeneous expanding phase, rather than a constant vacuum-linked
acceleration. A bound host can have another density history; this is a
conditional mismatch, not a general exclusion of the theory.

At fixed velocity, $A,\Lambda_F\propto b^{1/3}$ and $h\propto b^\gamma$
give exponents $1/3+3\gamma$ for $\beta$ and $2/3+2\gamma$ for
$K_{\rm ct}$. Their zeros are incompatible. Changing other spectral data
can evade this restricted result. The compensated cubic assignment selects
$\gamma=-1/9$ and requires the critical subtraction itself to vary with
material density. In particular
\begin{equation}
 b\partial_b(K_{\rm ct}p^2)=\tfrac49 K_{\rm ct}p^2.
\end{equation}
Omitting this contribution after matching the force changes the host
Euler/chemical equation. The total finite-cutoff excess has a small
nonzero rearrangement even when its leading cubic coefficient is constant.
For $Q=1,p=10^{-4},y=1$, it is
$b\partial_bV_F/V_F=1.66672916\times10^{-5}$, whereas the separate
counterterm contribution is $5.00175731\times10^{-10}$ and
$V_F=7.50235462\times10^{-14}$. The cancellation must be varied as a
whole. At nonzero $\boldsymbol P$, this derivative is a material
rearrangement at fixed $p$, not the full isotropic pressure: the norm,
frame and polarization also vary with the metric. Only the $P=0$ host
stress has been completed here.

\section{Finite inventory does not by itself cure the tail}
Prescribe, for diagnosis only, a spherical host
\begin{equation}
 y(r)=(1+r^2)^{-2},\quad \int d^3x\,y=\pi^2,\qquad
 D(r)=r/(1+r^2)^{3/2},\qquad C=Q=1.
\end{equation}
This is a finite-count profile with a chosen unit Plummer source, not a
solution of the coupled host or gravitational equations. The exact local
force equation from \eqref{mf:det} is
\begin{equation}
 D=\frac{\partial V_F}{\partial p}
 =\frac{2Ah\Lambda_F\Delta_f^2}
 {\Delta_f+\Lambda_F+\sqrt{\Delta_f^2+\Lambda_F^2}}.
 \label{mf:tail}
\end{equation}
It was solved on 161 radii $1\le r\le10^8$ per choice, with relative
force tolerance $10^{-9}$ and 60-digit energy evaluation. Fits over the
last 20 points give the asymptotic regulated-toy exponents
\begin{center}
\begin{tabular}{lccc}
\toprule
Choice & $p(r)$ & $dE_F/d\ln r=4\pi r^3V_F$ & low-gap tail retained?\\
\midrule
Flux & $r^{-1}$ & $r^0$ & yes\\
Material & $r^{-2}$ & $r^{-1}$ & no\\
Compensated cubic & $r^{-2/9}$ & $r^{7/9}$ & no\\
\bottomrule
\end{tabular}
\end{center}
The fitted exponents agree to better than $10^{-6}$ (gate $0.01$).
The low-gap flag means $\Delta_f/\Lambda_F<0.1$ at the last sample;
false flags are retained. The high-gap extensions describe this regulated
local expression only, not a controlled microscopic band theory.

For the flux choice the persistent cubic $p\propto r^{-1}$ gives a
logarithmically divergent excess energy despite finite host count. For
the material choice the high-gap expression becomes
$V_F\sim A\Lambda_Fh^2p^2/2$ with constant $A\Lambda_Fh^2$;
the toy returns to an inverse-square response and its tail energy converges.
Its singular coupling and vanishing cutoff as $b\to0$ require a new phase
completion before claiming a physical finite halo.

For the compensated choice on a general finite-mass power tail
$y\sim r^{-k}$, $k>3$, a putative cubic branch has
$\Delta_f/\Lambda_F\sim r^{-1+4k/9}$ and necessarily leaves the
small-gap regime. Continuing the regulated toy gives
\begin{equation}
 A\Lambda_Fh^2\sim r^{-4k/9},\quad
 p\sim r^{-2+4k/9},\quad
 \frac{dE_F}{d\ln r}\sim r^{-1+4k/9},
\end{equation}
which diverges. This is a conditional asymptotic obstruction for the
specified positive local band energy; it does not include a solved
interface, host relaxation, nonlocal band gradients or gravitational
binding energy and is not a universal no-go theorem.

\section{Research consequence and provenance}
A useful next route is a microscopically selected finite-density phase
with an actual interface and a controlled outside phase. A density plateau
could hold coefficients fixed over an annulus, but its density, charge,
inventory, stresses, interface energy and source-dependent formation must
all follow from one action. Merely prescribing a cutoff or profile does
not supply these conditions. The material-frame action here is a concrete
piece of that task, not the missing full fermion action.

The earlier scalar carrier, neutral production network and charge-even
portal remain distinct candidates requiring identification with this
host. Their abundance, mediator populations, narrow radiative thresholds
and transport obligations remain open. The auxiliary residual flux/history
condition, vacuum magnitude, fitted $\kappa$, finite-stock galaxy dynamics,
relativistic lensing and cosmological data tests also remain open. No
full physical closure or globally novel mechanism is claimed.

Evidence: \texttt{material\_frame\_audit\_run\_001} contains 156 controls;
\texttt{material\_tail\_probe\_run\_002} contains 972 controls, including
483 force roots and positive-energy checks. Both manifests passed freshness
validation. Tail run 001 is retained as failed: a NumPy boolean prevented
JSON serialization after calculation. Run 002 converts that output flag
to a built-in boolean; no physical gates or tolerances changed. Exact
identities, finite numerical ranges and out-of-regime flags are separated
above. This is self-review, without an independent referee or global
novelty search. Frozen theory criteria were not changed.

\begin{thebibliography}{9}
\bibitem{MFfluid} S. Dubovsky, L. Hui, A. Nicolis and D. T. Son,
\emph{Effective field theory for hydrodynamics: thermodynamics, and the
derivative expansion}, arXiv:1107.0731v1 (4 July 2011),
\url{https://arxiv.org/html/1107.0731v1}. Version identified by arXiv's
header; the generated internal typesetting date is not used as the
publication date. The primary fluid dictionary is established literature,
not an originality claim for this programme.
\end{thebibliography}

\clearpage

\part{A finite charged scalar phase and its reciprocal exterior}

\section{A separate candidate, not inherited closure}
A chosen charged scalar phase can provide a selected interior density,
a computed surface tension, and finite charge, while a reciprocal local
response changes from a cubic-dominated interior to a quadratic exterior.
This addresses the previous prescribed-tail obstruction by changing the
phase physics. Fixed-charge stationary spherical solutions are executed
below. The exterior also changes the inverse-square force coefficient;
there is no empirical, dynamic, microscopic or whole-theory closure.

This is a separately dated action candidate. The original positive-square
MeV scalar potential has $E\ge m|Q_c|$ and a repulsive remainder. It does
not supply an energetically protected lump with $E<m|Q_c|$. The attractive
interaction introduced here cannot inherit that bound, the old neutral
production calibration, or its abundance evidence. Metastable or
gravitationally confined bodies are not excluded by this energy observation.
The global host charge $N$ below is distinct from both $Q_c$ and the
auxiliary flux $Q$.

The charged-droplet mechanism is the established Q-ball construction;
Ref.~\cite{PIqball}, section II, gives its scalar conventions, charge and
radial equations. Its field is written $\phi_0 f_{\rm src}/\sqrt2$;
our $f=|\Psi|$ is dimensionful. Its sextic analysis and interface techniques
are prior work, not a novelty claim for this programme. The following
normalizations and the chosen response coupling are derived explicitly.

\section{Particle action, phase selection, and surface}
With signature $(-+++)$ choose
\begin{equation}
 \mathcal L_\Psi=-g^{\mu\nu}\partial_\mu\Psi^*\partial_\nu\Psi-V(f),
 \quad V=m_h^2 f^2-\lambda f^4+g_6 f^6,
 \quad \lambda,g_6>0,\quad m_h^2>\lambda^2/(4g_6).
 \label{pi:action}
\end{equation}
Here $[g_6]={\rm mass}^{-2}$; a sextic EFT needs ultraviolet control.
Define
\begin{equation}
 f_0^2=\lambda/(2g_6),\qquad
 \omega_0^2=m_h^2-\lambda^2/(4g_6)>0.
\end{equation}
The identity
\begin{equation}
 V=\omega_0^2 f^2+g_6 f^2(f^2-f_0^2)^2
\end{equation}
proves nonnegative potential and a minimum of $V/f^2$ at $f_0$ with
$\omega_0<m_h$. For $\Psi=f(r)e^{i\omega t}$ the positive charge convention
is $n=2\omega f^2$, $N=\int d^3x\,n$. In a homogeneous phase
\begin{equation}
 \rho_h=\omega^2 f^2+V,\qquad p_h=\omega^2f^2-V.
\end{equation}
At coexistence $\omega=\omega_0$, $f=f_0$, $p_h=0$ and the selected
charge density is $n_0=2\omega_0 f_0^2$. This is matter with a temporal
phase current, not vacuum energy with $w=-1$.

The exact planar wall grand potential at coexistence is
$U_0=V-\omega_0^2f^2=g_6 f^2(f^2-f_0^2)^2$.
The first integral $f'^2=U_0$ gives
\begin{equation}
 \frac{f^2(z)}{f_0^2}=\frac{1}{1+e^{2z/\ell}},\qquad
 \ell=\frac{1}{\sqrt{g_6}f_0^2},\qquad
 \sigma=\int_0^{f_0}2\sqrt{U_0}\,df
       =\frac{\sqrt{g_6}f_0^4}{2}.
 \label{pi:wall}
\end{equation}
These are planar coexistence results, not exact finite-radius profiles.
A thin spherical droplet has $p_h\simeq2\sigma/R$ and
$\omega-\omega_0\simeq2\sigma/(n_0R)$; its radius scales as
$(3N/(4\pi n_0))^{1/3}$. Charge fixes size, not the baryonic source by itself.

For the larger homogeneous root, $u=f^2$ obeys
$\omega^2=m_h^2-2\lambda u+3g_6u^2$. With $a=-\lambda+3g_6u>0$,
\begin{equation}
 \frac{dp_h}{d\omega}=n,\qquad
 c_s^2=\frac{n}{\omega\,dn/d\omega}
       =\frac{ua}{ua+\omega^2}.
\end{equation}
This is positive and below unity. At the example's coexistence point
it is $1/4$. The scalar action supplies its own material stress and
current; no thermodynamic identification with the previous label density
$b$ or occupied nodal fermions is assumed.

\section{One reciprocal static phase-response energy}
In a local rest-frame static model choose
\begin{align}
 E_{\rm stat}&=\int d^3x\left[
 |\nabla f|^2+\omega^2 f^2+V(f)
 +\frac C2|\nabla\varphi-\boldsymbol P/C|^2
 +W(f,p)+q_b\rho_b\varphi\right],\\
 W(f,p)&=\frac{K_0}{2}(1-s)^2p^2+\frac\beta3s p^3,
 \quad s=f^2/f_0^2,\quad p=|\boldsymbol P|,
 \quad C,K_0,\beta>0.
 \label{pi:energy}
\end{align}
In four-dimensional mass units $[f]={\rm mass}$, $[p]={\rm mass}^2$,
$[C]=[K_0]=1$, $[\beta]={\rm mass}^{-2}$ and $[q_b]={\rm mass}^{-1}$
if $\rho_b$ is a mass-energy density and $[\varphi]={\rm mass}$.
The switching functions are new choices; their critical zero and coefficients
have not been matched or radiatively protected. This is a static model,
not yet the covariant polarization/metric action.

Let $\boldsymbol D=C\nabla\varphi-\boldsymbol P$ with
$\nabla\cdot\boldsymbol D=q_b\rho_b$. For a spherical source,
$D=q_b M_b(r)/(4\pi r^2)$ and
\begin{equation}
 W_p=D/C\equiv d_r,\quad
 Kp+\beta s p^2=d_r,\quad K=K_0(1-s)^2,\qquad
 p=\frac{2d_r}{K+\sqrt{K^2+4\beta s d_r}}.
 \label{pi:root}
\end{equation}
The root is unique at positive $D$ and extends continuously to $s=0$.
At $s=1$, $p=\sqrt{d_r/\beta}$ is cubic; at $s=0$,
$p=d_r/K_0$ is quadratic. The source-centered radial field is regular
through the finite wall in the executed solutions.

At fixed source the correct reduced energy contribution is
$W-d_rp$, not $W$ alone. Indeed, integrate the source term by parts and
use $\nabla\varphi=(\boldsymbol D+\boldsymbol P)/C$ to obtain
$-D^2/(2C)+W-Dp/C$. The first term is independent of the host and omitted
from the reported reduced host energy. Varying $W-d_rp$ at its stationary
$p$ retains $W_f$ and supplies the reciprocal equation
\begin{equation}
 f''+\frac2r f'=\frac12 V_f-\omega^2f+\frac12 W_f,
 \quad
 W_f=\left[K_0(s-1)p^2+\frac\beta3p^3\right]\frac{2f}{f_0^2}.
 \label{pi:bvp}
\end{equation}
Thus the phase is varied with its response backreaction; the interior
is not imposed as a sharp prescribed density step.

\section{Executed finite-charge spherical branches}
Choose dimensionless units $m_h=\lambda=g_6=1$,
$f_0^2=1/2$, $\omega_0=\sqrt{3/4}$, $\ell=2$, $\sigma=1/8$,
$C=1$, $K_0=0.1$, $\beta=1$. Prescribe only the baryons through
$d_r=S r/(1+r^2)^{3/2}$ with $S=0,0.01,0.05$.
For $R_*=20,40$, fix the charge
$N=(4\pi/3)n_0R_*^3=29020.78983,232166.31862$.
Solve \eqref{pi:bvp} together with $N'(r)=8\pi\omega r^2f^2$,
regular $f'(0)=0$, $N(0)=0$, $f(R_{\rm box})=0$ and the specified
$N(R_{\rm box})$. The unknown $\omega$ is solved, not imposed.
$R_{\rm box}=R_*+50$. The singular radial term is handled explicitly
by the collocation solver. The box boundary approximates the infinite
localized scalar tail; amplitudes five units inside it are below $10^{-10}$.

Two meshes/tolerances and independent 8001-point charge/energy quadrature
were used. Representative refined outputs are
\begin{center}
\begin{tabular}{rrrrr}
\toprule
$R_*$ & $S$ & $\omega$ & $E_{\rm red}/(m_h N)$ & max interior $K/(\beta s p)$\\
\midrule
20 & 0 & 0.881021309 & 0.889219746 & ---\\
20 & 0.01 & 0.881021210 & 0.889219080 & 0.029182\\
20 & 0.05 & 0.881020297 & 0.889212062 & 0.012941\\
40 & 0 & 0.873382644 & 0.877254306 & ---\\
40 & 0.01 & 0.873382632 & 0.877254196 & 0.017948\\
40 & 0.05 & 0.873382510 & 0.877253070 & 0.007986\\
\bottomrule
\end{tabular}
\end{center}
Here $E_{\rm red}=E_h+\int(W-d_rp)$; the omitted source-only energy
is finite. $E_h/(m_hN)<1$ also holds separately, but that comparison
alone is not a proof of full sourced stability. Half-coexistence-density
radii are approximately $19.504,39.533$. The tested interior is
$5<r<0.75R_{1/2}$, not the wall. Its quadratic/cubic force ratio remains
below $0.1$ in all four sourced refined solutions, and the resulting
annulus is finite. The exterior $p\sim r^{-2}$ gives
$4\pi r^3W\sim r^{-1}$; both scalar and response tails have finite energy
in this local candidate. The exterior is an analytic asymptotic matching
after the negligible host tail, not a microscopic nodal-band calculation.

Charge errors are below $4\times10^{-10}$ and mesh changes in
$\omega,E_{\rm red}/N$ below $4\times10^{-9}$. A separate fixed-charge
scaling identity,
\begin{equation}
 2T-3\lambda I_4+6g_6 I_6+
 \int d^3x\,W_f(\tfrac32 f+r f')=0,
 \quad I_j=\int d^3x f^j,\quad T=\int d^3x f'^2,
\end{equation}
checks stationarity without relying only on the solver's residual.
All normalized residuals are below $5\times10^{-9}$.

\section{Finite radial curvature test}
Eliminate $\omega$ at fixed $N$ to obtain $N^2/(4I_2)$ and eliminate
static $p$ at fixed source. The envelope curvature is
\begin{equation}
 H_{ff}=W_{ff}-\frac{W_{fp}^2}{W_{pp}},\qquad W_{pp}=K+2\beta s p.
\end{equation}
For radial perturbations $\delta f=v/r$, the second variation is the
quadratic form of
\begin{equation}
 \mathcal H_0=-2\partial_r^2+V_{ff}-2\omega^2+H_{ff}
 +\frac{32\pi\omega^2}{I_2}|rf\rangle\langle rf|,
 \label{pi:hessian}
\end{equation}
with Dirichlet $v$ at zero and the box boundary and the radial inner
product $\int dr$. The positive rank-one term is required by fixed charge.
The static $S=0$, $p=0$ curvature is understood by the continuous zero
response limit.

Centered differences at 800 and 1600 interior points and three lowest
symmetric eigenpairs per profile yield positive sampled radial amplitude
curvature for all six refined branches. The smallest eigenvalues at 1600
points span $0.398090$--$0.398097$ for $R_*=20$ and
$0.216946$--$0.216949$ for $R_*=40$. Eigenvector residuals are below
$4\times10^{-10}$ and resolution shifts below $1.1\times10^{-4}$.
This is finite static numerical evidence, not a continuum spectral proof,
angular or nonlinear stability result, or a polarization kinetic/causality
certificate. The phase and charge can still have dynamical obligations
outside this amplitude test.

\section{An exterior-force cost and remaining bridge}
The desired finite exterior is not invisible to gravity. Outside,
\begin{equation}
 \varphi'=\frac DC\left(1+\frac{1}{CK_0}\right),\qquad
 G_{\rm scalar,out}=\frac{q_b^2}{4\pi C}
               \left(1+\frac{1}{CK_0}\right).
\end{equation}
The interior bare inverse-square coefficient is $q_b^2/(4\pi C)$;
the example's exterior coefficient is eleven times larger. Einstein
attraction must be added once in the actual measured dictionary. This is
not itself an observational exclusion, but it forbids importing the old
Newtonian calibration or a lensing verdict unchanged.

There is a conditional thin-wall tradeoff. In the selected units the
Laplace-pressure shift gives $s-1\simeq1/R$, hence $K\sim K_0/R^2$.
Near $r\sim R$ on a putative cubic branch,
$K/(\beta s p)\sim K_0/(\sqrt{\beta S}R)$. Demanding small exterior
scalar enhancement $1/(CK_0)<\eta$ and interior ratio below $\epsilon$
therefore requires parametrically
$R\gtrsim1/(C\eta\epsilon\sqrt{\beta S})$. This estimate assumes a
thin weakly sourced phase and is not a universal bound or a data fit.
It identifies a cost of simply raising $K_0$ to hide the exterior.

The new host needs a charge inventory and formation mechanism. A neutral
population could contain charged droplets of opposite signs, but their
production, separation, collisions and annihilation have not been derived.
Sharing $Q$-dependent scales would not by itself create that inventory.
A finite phase of this type is a candidate particle-level host, not an
identification of the previous MeV excitations or the meV nodal sea.
The sextic cutoff, critical matching, covariant polarization dynamics,
complete stress/lensing, parameter origins, flux/history residual, vacuum
magnitude, abundance and cosmology remain open. No complete theory or
global novelty claim follows from these finite stationary solutions.

\section{Evidence}
\texttt{phase\_interface\_probe\_run\_001}: 87 controls including exact
wall identities, twelve BVP solutions, charge/virial/root checks and
refinement comparisons. \texttt{phase\_radial\_curvature\_run\_001}:
twelve finite eigenpair residual checks, with curvature signs retained as
scientific flags. Scripts, input profiles, bounds and raw outputs are pinned
by their manifests. Both current manifests passed freshness validation.
This report is self-reviewed, without an independent referee.

\begin{thebibliography}{9}
\bibitem{PIqball} J. Heeck, A. Rajaraman, R. Riley and C. B. Verhaaren,
\emph{Understanding Q-Balls Beyond the Thin-Wall Limit},
arXiv:2009.08462v2 (9 February 2021),
\url{https://arxiv.org/html/2009.08462v2}, section II, equations (1)--(8).
The established construction is used with the explicit normalization
translation above, not imported as proof of this coupled candidate.
\end{thebibliography}

\clearpage

\part{An explicit mediator origin and validity of the scalar phase}

\section{What is established}
A canonical real mediator can generate the preceding host's attractive
quartic and stabilizing sextic terms in a renormalizable two-scalar
Lagrangian. An explicit one-parameter family has a unique globally stable
classical vacuum, an exactly characterized homogeneous charged coexistence
phase, and a controlled heavy-mediator limit at tree level. The matching
also changes the host's amplitude kinetic term. These results close the
specified classical mediator-potential question, not the full theory.
The response switching function, quantum thresholds, charge production,
formation, gravity and dark-energy normalization remain unresolved.

Complex charged scalars with real mediators are established soliton models,
including the Friedberg--Lee--Sirlin family \cite{uv:FLS}. Attractive cubic
couplings and quantum corrections have been investigated in related
models \cite{uv:attractive,uv:quantum}. These references establish related
context only; their stability results are not transferred to this potential.
No global novelty claim is made.

\section{Canonical particle model and exact vacuum bound}
Set the host vacuum mass $m_h=1$ as the mass unit and let $M\ge1$ be
the mediator mass in that unit. With $I=|\Psi|^2$ and real $\chi$, take
\begin{align}
 \mathcal L&=-|\partial\Psi|^2-\tfrac12(\partial\chi)^2-V(I,\chi),\\
 V&=I+L I^2+2M\chi I+U(\chi),\qquad L=2-\frac1M,\\
 U(\chi)&=\frac{M^2}{2}\chi^2-\frac M8\chi^3+\frac14\chi^4.
 \label{uv:potential}
\end{align}
All interactions have canonical mass dimension at most four; this is a
renormalizable classical particle model. This does not assert an
all-energy ultraviolet completion or quantum vacuum stability. Dimensionful
coefficients are restored by multiplying mass-dimension-one quantities by
$m_h$ and dimension-two quantities by $m_h^2$; in particular $L$ means
$2-m_h/M_{\rm physical}$.

For $\chi\ge0$ every $I$ term is nonnegative, while
\begin{equation}
 U(\chi)=\chi^2\left[\frac14(\chi-M/4)^2+\frac{31M^2}{64}\right]>0
 \quad(\chi\ne0).
\end{equation}
For $\chi=-x<0$, $U(-x)=M^2x^2/2+Mx^3/8+x^4/4>0$. If
$x\le1/(2M)$, the coefficient $1-2Mx$ of $I$ is nonnegative.
Otherwise minimize exactly over $I\ge0$:
\begin{align}
 I_{\rm min}&=(2Mx-1)/(2L),\\
 4L V_{\rm min}&=Mx\left[\frac L2x^2-2x+4\right]+Lx^4-1.
\end{align}
Since $1\le L<2$,
\begin{equation}
 \frac L2x^2-2x+4=\frac L2(x-2/L)^2+4-2/L\ge2.
\end{equation}
Thus $4LV_{\rm min}\ge2Mx+Lx^4-1>0$ for $x>1/(2M)$.
All field values have been covered: $V\ge0$ with equality only at
$I=\chi=0$, for every $M\ge1$. This is an algebraic global statement,
not an inference from numerical sampling.

Moreover
\begin{equation}
 U''(\chi)=M^2-\tfrac34M\chi+3\chi^2
 =3(\chi-M/8)^2+\frac{61M^2}{64}>0.
\end{equation}
The stationary mediator equation
$U'(\chi)+2MI=0$ therefore has one real solution for every $I\ge0$.
There is no branch ambiguity in its tree-level local elimination.

\section{Tree matching includes gradients}
For general coefficients
$V=m_h^2 I+\lambda_\Psi I^2+a\chi I+M^2\chi^2/2+
\mu\chi^3/3+\lambda_\chi\chi^4/4$, the local stationary branch is
\begin{equation}
 \chi_*(I)=-\frac a{M^2}I-\frac{\mu a^2}{M^6}I^2+O(I^3).
\end{equation}
Substitution yields
\begin{equation}
 V_{\rm eff}=m_h^2 I+
 \left(\lambda_\Psi-\frac{a^2}{2M^2}\right)I^2
 -\frac{\mu a^3}{3M^6}I^3+O(I^4).
\end{equation}
For \eqref{uv:potential}, $a=2M$, $\mu=-3M/8$ and
$\lambda_\chi=1$, hence
\begin{equation}
 V_{\rm eff}=I-\frac{I^2}{M}+\frac{I^3}{M^2}+O(I^4),
 \qquad \chi_*=-2I/M+\frac{3I^2}{2M^3}+O(I^3).
 \label{uv:sextic}
\end{equation}
At $M=1$ the first three terms are the previous illustrative sextic
potential; higher terms need not be negligible there.

Eliminating a propagating mediator also produces derivative operators:
\begin{equation}
 -\tfrac12(\partial\chi_*)^2
 =-\frac{2}{M^2}(\partial I)^2+\cdots.
\end{equation}
For $\Psi=f e^{i\theta}$ the exact adiabatic amplitude kinetic factor is
\begin{equation}
 Z_f=1+2I\left(\frac{2M}{U''(\chi_*)}\right)^2
 =1+\frac{8I}{M^2}+\cdots.
 \label{uv:kinetic}
\end{equation}
The temporal phase term is $-I(\partial\theta)^2$. These expressions
assume sufficiently slow backgrounds; integrating out $\chi$ dynamically
also generates higher derivatives. Therefore accurate potential matching
alone does not justify importing the old wall, radius or curvature results.

\section{Exact charged coexistence and its finite check}
For a homogeneous charged phase, $\Psi=\sqrt I e^{i\omega t}$,
the quantity determining coexistence is
$\omega_*^2=\min_{I>0,\chi}V/I$. At $\chi=-x$, minimizing over $I$
gives $I_*=\sqrt{U(-x)/L}$ and
\begin{equation}
 R(x)=1-2Mx+2\sqrt{L U(-x)}.
\end{equation}
The derivative is negative near $x=0$ and positive at large $x$.
Its unique positive stationary point is the positive root of
\begin{multline}
 Lx^4+\frac{3LM}{4}x^3+
 \left(\frac{137L}{64}-1\right)M^2x^2\\
 +\left(\frac{3L}{4}-\frac12\right)M^3x-M^3=0.
 \label{uv:root}
\end{multline}
All nonconstant coefficients are positive for $L\ge1$; the polynomial
is strictly increasing for $x>0$, from a negative value to infinity.
It follows that this root is the global minimum of $R$ on the negative
mediator branch. Positive $\chi$ cannot improve on $V/I\ge1$.
The vacuum bound and the negative initial slope give
$0<\omega_*^2<1$. The density is $n_*=2\omega_* I_*$ and the pressure
is zero: $V=\omega_*^2 I_*$. These are exact homogeneous classical
results, not a new finite-size two-field droplet solution.

The sextic truncation predicts $I_0=M/2$ and $\omega_0^2=3/4$.
Solving \eqref{uv:root} and the exact mediator branch gives
\begin{center}
\begin{tabular}{rrrrr}
\toprule
$M/m_h$ & $I_*$ & $\omega_*^2$ & potential error & $Z_f$\\
\midrule
1 & 0.530967 & 0.777373 & $-3.398\%$ & $1.574$\\
10 & 4.351594 & 0.774215 & $-2.585\%$ & $1.296$\\
100 & 48.989310 & 0.753423 & $-0.441\%$ & $1.039$\\
1000 & 498.929051 & 0.750358 & $-0.0475\%$ & $1.004$\\
\bottomrule
\end{tabular}
\end{center}
The error is $(V_{\rm sextic}-V_{\rm exact})/V_{\rm exact}$ evaluated
at the exact phase density. The first two fail the chosen one-percent
potential comparison, retained as false flags. The $M=10$ density shift
is approximately $-13\%$ despite its smaller potential error; derivatives
and stationary states require their own check. Exact amplitude/mediator
mass matrices at fixed $\omega_*$ have positive eigenvalues at the four
tested points; their lower eigenvalues are approximately
$0.3392,0.7355,0.9620,0.9960$. This is not a finite-body or full dynamical
stability certificate.

In the large-$M$ limit, \eqref{uv:root} divided by $M^3$ tends $x-1=0$.
Thus $x_*\to1$, $I_*/M\to1/2$, $\chi_*/M\to0$, $\omega_*^2\to3/4$
and $Z_f\to1$. With $\Psi=\sqrt M\,\psi$, the leading effective action
is $M$ times the previous unit-coefficient sextic action. This explains
the classical shape limit; the conserved charge and energy scale with $M$,
so the old fixed inventories must not be silently retained.

\section{What still prevents a full particle bridge}
The attractive quartic is the difference
$\lambda_\Psi-a^2/(2M^2)=-1/M$ of two order-one terms.
Its selected fractional cancellation is $1/(2M)$: $0.5\%$ at $M=100$
and $0.05\%$ at $M=1000$. Increasing the hierarchy improves the tree
expansion but requires increasingly accurate matching of that difference.
No symmetry or renormalization-group argument protecting it has been
provided. Loop thresholds and the independent mass, tadpole and vacuum
counterterms must be calculated before claiming quantum closure. No
scheme-dependent partial loop term is used as a vacuum prediction here.

The mediator potential also does not derive the previous phase-response
function $W(f,p)$ or the nodal medium. A new two-field finite-body solution
with its gradients and full response is needed before transferring any
stationarity or radial-curvature result. Charge creation/separation,
thermal populations, abundance, an energy receiver and source-dependent
formation remain open. The exterior-force calibration cost persists.
This real mediator is not automatically the original source $X$, its
charge-even bridge $\sigma$, or the auxiliary vacuum variable $Q$.
The vacuum magnitude, flux/history residual, acceleration normalization,
relativistic gravity and cosmological/empirical obligations remain open.

\section{Evidence and review}
\texttt{mediator\_phase\_matching\_run\_002} contains 37 exact/finite
controls: vacuum algebra, elimination and root identities, four exact
coexistence comparisons, sixteen stationary-branch samples and four
homogeneous amplitude-matrix tests. Run 001 completed the same core
calculation; run 002 adds the kinetic-factor output. Potential-validity
flags are scientific results, not forced passing assertions. The current
manifest passed input/output freshness validation. This is a self-review
with the explicit proof above, not an independent-referee or quantum
stability certificate.

\begin{thebibliography}{9}
\bibitem{uv:FLS} J. Heeck and M. Sokhashvili,
\emph{Revisiting the Friedberg-Lee-Sirlin soliton model},
arXiv:2303.09566v2, \url{https://arxiv.org/abs/2303.09566v2}.
\bibitem{uv:attractive} Y. Hamada, K. Kawana, T. Kim and P. Lu,
\emph{Q-Balls in the presence of attractive force},
arXiv:2407.11115v2, \url{https://arxiv.org/abs/2407.11115v2}.
\bibitem{uv:quantum} Y.-X. Su, Q.-X. Xie and S.-Y. Zhou,
\emph{Quantum-Corrected Q-balls in the Friedberg-Lee-Sirlin Model},
arXiv:2605.25243, \url{https://arxiv.org/abs/2605.25243}.
Only official primary abstracts/metadata of these three contextual sources
were checked here. Their detailed theorems and numerical claims are not
used as proof of this candidate.
\end{thebibliography}

\clearpage

\part{Quantum control, formation, response cones and vacuum normalization}

\section{Results and their limits}
Four requested gaps are tested on the actual mediator candidate. Its
original order-one couplings do not protect the small attractive quartic
or the light host mass under one-loop running. A separately specified
weak-coupling scaling suppresses those sensitivities while preserving
classical host shapes and instability rates, at the cost of larger charge
and energy inventories. Dilute charged coherent backgrounds have a linear
clumping instability; this does not derive a cosmological charged initial
state or nonlinear droplet production.

The minimal Lorentz-invariant scalar realization of the cubic static law
has a radial characteristic speed approaching $\sqrt2$ times the metric
light speed. Thus the static force calculation alone does not supply a
completion with a common metric cone. Independent medium dynamics may
change that result but have not yet been derived. Finally, the present
particle and gravity equations leave independent response coefficients and
vacuum flux/history data. They do not predict a unique dark-energy density
or acceleration normalization. None of these gaps is relabeled as closed.

\section{Full scalar one-loop divergence and matching sensitivity}
Write the mediator potential in general form
\begin{equation}
 V=V_0+t\chi+m_h^2 I+L I^2+a\chi I+
 \tfrac12 B\chi^2+\tfrac\mu3\chi^3+\tfrac l4\chi^4,
 \qquad I=|\Psi|^2,\quad B=M^2.
\end{equation}
The previous tree family has $V_0=t=0$, $m_h=1$,
$L=2-1/M$, $a=2M$, $\mu=-3M/8$, $l=1$.
The independent vacuum and tadpole counterterms must nevertheless be kept.
For $\Psi=(\phi_1+i\phi_2)/\sqrt2$ and
$\phi_1=\sqrt{2I},\phi_2=0$, the canonical real-scalar Hessian is
\begin{equation}
 H=\begin{pmatrix}
 m_h^2+6LI+a\chi&0&a\sqrt{2I}\\
 0&m_h^2+2LI+a\chi&0\\
 a\sqrt{2I}&0&B+2\mu\chi+3l\chi^2
 \end{pmatrix}.
\end{equation}
The real-scalar MS one-loop convention is
\begin{equation}
 V_1=\frac{1}{64\pi^2}\sum_j h_j^2
       [\ln(h_j/\mu_R^2)-3/2],\qquad h_j=\operatorname{eig}(H).
\end{equation}
Its normalization is checked against Ref.~\cite{qfg:Martin}; the following
polynomial projections are derived here. This expression for a stable
background is not a finite-density fluctuation determinant at nonzero
chemical potential. Negative $h_j$ require the usual instability treatment;
no real global quantum potential is inferred from them.

At one loop the ultraviolet field-strength poles vanish for these scalar
interactions. Scale independence gives
$\sum_i\beta_i\partial_i V=\operatorname{tr}H^2/(32\pi^2)$.
The exact polynomial is
\begin{align}
 \operatorname{tr}H^2={}&40L^2I^2+(16Lm_h^2+4a^2)I+16LaI\chi
 +9l^2\chi^4+12l\mu\chi^3\nonumber\\
 &+(2a^2+6Bl+4\mu^2)\chi^2
 +(4m_h^2a+4B\mu)\chi+2m_h^4+B^2.
 \label{qfg:trace}
\end{align}
Let $\ell_1=16\pi^2$. Its coefficients yield
\begin{align}
 \ell_1\beta_L&=20L^2,&\ell_1\beta_a&=8La,
 &\ell_1\beta_{m_h^2}&=8Lm_h^2+2a^2,\\
 \ell_1\beta_B&=2a^2+6Bl+4\mu^2,
 &\ell_1\beta_\mu&=18l\mu,&\ell_1\beta_l&=18l^2,\\
 \ell_1\beta_t&=2m_h^2a+2B\mu,
 &\beta_{V_0}&=(2m_h^4+B^2)/(32\pi^2).
\end{align}
These are full-theory MS beta functions at this potential's zero quartic
mixed-operator point, not finite threshold matching or running in a
heavy-field-decoupled low-energy EFT.

For the attractive tree coefficient $\lambda_{\rm net}=L-a^2/(2B)$,
\begin{equation}
 \ell_1\beta_{\lambda_{\rm net}}
 =20L^2-\frac{8La^2}{B}
 +\frac{a^2}{2B^2}(2a^2+6Bl+4\mu^2).
\end{equation}
On the selected family this becomes
$[20L^2-32L+233/8]/\ell_1$.
At $M=100$, $\lambda_{\rm net}=-0.01$ but
$\beta_{\lambda_{\rm net}}=0.2827304$;
$\beta_{m_h^2}/m_h^2=506.7067$. The small matching relation is not protected
by the tree vacuum proof. A linearized scale interval
$|\delta\ln\mu_R|\simeq0.03537$ already corresponds to a shift equal to
the chosen quartic magnitude. This diagnostic is not the claim that a
physical observable varies with an arbitrary renormalization scale:
consistent loop matching restores scale invariance, and finite
renormalization can impose measured inputs at a chosen scale.

Tadpoles are not omitted. With $t=0$ initially, following the stationary
vacuum gives $dv/d\ln\mu_R=-\beta_t/B$. At that point the corresponding
curvature-mass and local tree-quartic projections are
\begin{align}
 \beta_{m_h^2}^{\rm follow}&=\beta_{m_h^2}-a\beta_t/B,\\
 \beta_{\lambda_{\rm net}}^{\rm follow}
 &=\beta_{\lambda_{\rm net}}-a^2\mu\beta_t/B^3.
\end{align}
On the family the second correction is
$(6/M^2-9/8)/\ell_1$. It does not remove the order-one-coupling sensitivity.
These local projections are still not pole-mass or scattering threshold
matching conditions.

\section{A controlled classical-limit route}
Introduce a new freely selected parameter $0<\epsilon\ll1$ and replace
\begin{equation}
 L\mapsto\epsilon L,\quad a\mapsto\sqrt\epsilon a,\quad
 \mu\mapsto\sqrt\epsilon\mu,\quad l\mapsto\epsilon l,
 \quad m_h^2,B\ \hbox{fixed}.
\end{equation}
For $\Psi=\psi/\sqrt\epsilon$ and $\chi=X/\sqrt\epsilon$, the exact
scalar action is $S_\epsilon=S_1[\psi,X]/\epsilon$.
Its classical field equations, background Hessian eigenvalues and
formation dispersion relations in the rescaled fields are identical.
The stable-vacuum proof and coexistence phase persist under this invertible
rescaling. Charge, energy and interface tension for a fixed shape grow as
$1/\epsilon$; the radius does not grow unless the specified inventory does.

The matching coefficient scales as $\epsilon\lambda_{\rm net}$,
its one-loop beta function as $\epsilon^2\beta_{\lambda_{\rm net}}$,
and the mass beta as $\epsilon\beta_{m_h^2}$. Consequently a weakly coupled
version can control both diagnostics without changing the classical
host mechanism. At $M=100$:
\begin{center}
\begin{tabular}{rrrr}
\toprule
$\epsilon$ & $|\beta_{\lambda_{\rm net}}/\lambda_{\rm net}|$
 & $|\beta_{m_h^2}/m_h^2|$ & inventory multiplier\\
\midrule
1 & 28.273 & 506.707 & 1\\
$10^{-3}$ & 0.02827 & 0.50671 & $10^3$\\
$10^{-6}$ & $2.8273\times10^{-5}$ & $5.0671\times10^{-4}$ & $10^6$\\
\bottomrule
\end{tabular}
\end{center}
The last row passes a specified one-percent, unit-log sensitivity check;
the first two do not. This is a quantum-control route, not a calculation
of all finite thresholds or a quantum vacuum theorem. Large classical
fields require control of any additional operators in a higher theory.
Portal interactions, response couplings, gravitational backreaction and
cosmological inventories must also be matched anew. The parameter
$\epsilon$ and the increased stock have not been predicted.

\section{Formation: an actual instability, with specified initial state}
At a homogeneous charged background, solve the actual mediator equation
$U'(\chi)+aI=0$ and
$\omega^2=m_h^2+2LI+a\chi>0$. Write
$\Psi=e^{i\omega t}[\sqrt I+(\eta+i\pi)/\sqrt2]$ and perturb $\chi$
by $\xi$. Define
\begin{equation}
 A=4LI,\qquad B_h=U''(\chi),\qquad h=a\sqrt{2I},\qquad
 A_{\rm eff}=A-h^2/B_h.
\end{equation}
For wave number $k$, the exact linear scalar system has potential matrix
\begin{equation}
 K(k)=\begin{pmatrix}k^2+A&0&h\\0&k^2&0\\h&0&k^2+B_h\end{pmatrix}
\end{equation}
and antisymmetric rotation coupling $G_{12}=-G_{21}=-2\omega$ in
$\ddot q+G\dot q+Kq=0$. For $z=\Omega^2$, its dispersion equation is
\begin{equation}
 (k^2+B_h-z)[(k^2+A-z)(k^2-z)-4\omega^2 z]
 -h^2(k^2-z)=0.
 \label{qfg:dispersion}
\end{equation}
The long-wavelength mode has
$c_s^2=A_{\rm eff}/(A_{\rm eff}+4\omega^2)$.
If $A_{\rm eff}<0$ and the denominator is positive, there is a growing
long-wavelength mode. Its unstable band ends at
\begin{equation}
 k_c^2=\frac{2(h^2-AB_h)}
 {\sqrt{(A-B_h)^2+4h^2}+A+B_h}>0.
\end{equation}
A positive $A_{\rm eff}$ instead makes the quadratic energy nonnegative
for all $k$, giving linear stability in this conservative scalar sector.

The real-time six-dimensional generator was diagonalized at 201 wave
numbers $0\le k\le2$, for $M=10,100$ and $I/I_*=0.1,0.5,1,2$.
Its eigenvalues agree with the independently derived stable/unstable-band
classification across the sampled points. At $M=100$:
\begin{center}
\begin{tabular}{rrrrr}
\toprule
$I/I_*$ & $c_s^2$ & $k_c$ & fastest sampled $k$ & growth rate $\gamma$\\
\midrule
0.1 & $-0.04815$ & 0.40795 & 0.29 & 0.04369\\
0.5 & $-0.09924$ & 0.49502 & 0.34 & 0.07436\\
1 & 0.24898 & 0 & --- & no sampled growth\\
2 & 0.50357 & 0 & --- & no sampled growth\\
\bottomrule
\end{tabular}
\end{center}
Units are $m_h$ for wave number/growth and $m_h^{-1}$ for time.
The weak-coupling scaling preserves these rates at the rescaled density.
This is a concrete possible fragmentation starting point, conditional on
preparing a charged coherent state in that range.

The empty vacuum has positive scalar masses and no such instability. A
neutral thermal gas is not this coherent charged state. Its production,
charge separation and occupation cannot be inferred from the table.
Nonlinear evolution, charge/energy conservation and radiation must be
executed before claiming actual droplets or abundance. An isolated system
can radiate into its own scalar waves; an external heat bath is not
logically necessary. The calculation does not establish which available
radiation channels can remove the binding energy at a sufficient rate.
The raw outputs record the initial energy per charge and the bulk binding
budget; finite surfaces and emitted charge must be counted separately.

\section{A covariant coupling test: the minimal scalar cone}
A conformal baryon metric $\widetilde g_{\mu\nu}=e^{2q_b\varphi}g_{\mu\nu}$
reproduces the weak static source energy $q_b\rho_b\varphi$. A Lorentz
auxiliary vector can reproduce the static square through
\begin{equation}
 \mathcal L_{\rm resp}=-\tfrac C2(\partial\varphi)^2
 +P^\mu\partial_\mu\varphi-\frac{P^\mu P_\mu}{2C}-W(P^2).
\end{equation}
On its spacelike static branch choose $W=\beta(P^2)^{3/2}/3$.
This defines a local branch, not a prescribed timelike or vacuum extension.
Eliminating $P$ yields a single-gradient scalar Lagrangian $-F(X)$,
$X=(\partial\varphi)^2/2>0$. If $y=|\nabla\varphi|$, the exact
constitutive equations are
\begin{equation}
 y=p/C+\beta p^2,\qquad F_y=D=C\beta p^2.
\end{equation}
The principal tensor is
$F_Xg^{\mu\nu}+F_{XX}\partial^\mu\varphi\partial^\nu\varphi$.
On the static radial branch it gives
\begin{equation}
 c_r^2=\frac{y\,dD/dy}{D}
 =\frac{2(1+C\beta p)}{1+2C\beta p},\qquad c_\perp^2=1.
 \label{qfg:cone}
\end{equation}
Thus the cubic-dominated limit $C\beta p\ll1$ has $c_r^2\to2$.
The temporal kinetic coefficient is positive for $p>0$ but tends to zero
as the critical branch approaches $p=0$; a uniform perturbative vacuum
limit is not supplied. Conformal coupling leaves the baryonic null cone
unchanged, so this is faster than that metric cone. It fails a requirement
that every mode propagate within the same cone. Superluminal scalar
characteristics alone are not a complete proof of a causal contradiction
or an empirical exclusion of all alternatives.

For the quadratic exterior $W=K_0P^2/2$, the vector can be eliminated
covariantly without a material frame and gives a positive scalar kinetic
coefficient $C_{\rm eff}=C^2K_0/(1+CK_0)$. The exterior therefore need
not rely on the undefined phase current outside a droplet. Its enhanced
Newtonian force remains the earlier $1+1/(CK_0)$ factor.
Independent medium/frame dynamics may alter the interior characteristic
relation; their action and radiative matching still have to be derived.
Nothing in the scalar mediator potential fixes $\beta$, $C$, $q_b$ or the
critical response switching functions. Relativistic lensing and the full
metric field equations remain required as well.

\section{Dark energy: identifiable inputs and a decisive degeneracy}
The local one-loop vacuum divergence includes
$(2m_h^4+M^4)/(32\pi^2)$ in $\beta_{V_0}$, with an independent finite
vacuum counterterm. It is not a predicted cosmological constant. Even the
weak-coupling limit does not remove the free massive vacuum determinant;
its relative smallness compared with a large classical condensate is not
an explanation of the observed vacuum scale.

The previous auxiliary construction instead has
\begin{equation}
 B_{\rm eff}=Q^2/2+\Delta_F+\langle\tau\rangle/4.
\end{equation}
Vacuum-form matter constants can project out of its metric equation under
its stated hypotheses, but their $Q$ derivatives remain in the conjugate
flux equation. That algebra does not select $Q$ or the residual
flux/history term. These are global data, not a result of the host's
local coexistence condition.

For $W=\beta p^3/3$, the annular force coefficient is
\begin{equation}
 a_0=\frac{q_b^3}{4\pi C^3\beta G_{\rm lab}}.
\end{equation}
Holding the particle potential and all $P=0$ phase data fixed while
changing $\beta$ from 1 to 4 changes $a_0$ by a factor of four.
Holding $Q=1$ fixed while changing
$\Delta_F+\langle\tau\rangle/4$ from 0 to 1 changes $B_{\rm eff}$
from $1/2$ to $3/2$. Both variations are allowed by the current parameter
and boundary-data specification. Therefore no unique $\kappa$ or physical
dark-energy normalization follows from it. These counterexamples apply
to this specification, not every possible future completion.

A further mechanism must fix the response coefficients and flux/history
data independently of the measured target. Otherwise the honest result
is a calibrated effective theory. Selecting those data from the observed
$a_0$ and dark-energy density is not a derivation. Flux quantization or a
vacuum-selection proposal would also need its own charge scale and actual
global solution; naming it supplies neither condition.

\section{Evidence and next implication}
\texttt{four\_gap\_audit\_run\_003} has 3219 exact/finite controls,
including 1608 generator spectra and independent band checks. Run 001
completed the core; run 002 added tadpole-following and energy budgets;
run 003 added independent spectral-band controls. No physical thresholds
were loosened. \texttt{weak\_coupling\_control\_run\_001} has five exact
scaling checks and three finite sensitivity comparisons with passing and
failing scientific flags retained. Both current manifests passed freshness
validation. These are self-reviewed results, not independent certificates.

The next useful joined calculation is nonlinear charged-state fragmentation
in the weakly coupled mediator model with explicit outgoing scalar energy
and charge, followed by a physically specified response/metric completion
that meets its declared characteristic-cone and lensing requirements.
Quantum finite matching and independently selected global vacuum data
remain necessary. The full cold-energy/dark-energy theory is still open.

\begin{thebibliography}{9}
\bibitem{qfg:Martin} S. P. Martin,
\emph{Two-loop effective potential for a general renormalizable theory and
softly broken supersymmetry}, arXiv:hep-ph/0111209v2,
\url{https://arxiv.org/pdf/hep-ph/0111209v2}, equations (1.1), (1.3), (3.2).
The previously authenticated real-scalar one-loop convention is reused;
new beta projections, dispersion and response characteristics are derived
above. No external finite-density or full-model stability theorem is imported.
\end{thebibliography}

\clearpage

\end{document}
